% 高一59_华实高中_10月月考.tex
%   —— 高一上数学“2026-2027 年华实高中高一上 10 月月考”（讲评版）
% 试卷版：xelatex 高一59_华实高中_10月月考.tex
% 答案版：xelatex "\def\isanswer{1}% 高一59_华实高中_10月月考.tex
%   —— 高一上数学“2026-2027 年华实高中高一上 10 月月考”（讲评版）
% 试卷版：xelatex 高一59_华实高中_10月月考.tex
% 答案版：xelatex "\def\isanswer{1}% 高一59_华实高中_10月月考.tex
%   —— 高一上数学“2026-2027 年华实高中高一上 10 月月考”（讲评版）
% 试卷版：xelatex 高一59_华实高中_10月月考.tex
% 答案版：xelatex "\def\isanswer{1}% 高一59_华实高中_10月月考.tex
%   —— 高一上数学“2026-2027 年华实高中高一上 10 月月考”（讲评版）
% 试卷版：xelatex 高一59_华实高中_10月月考.tex
% 答案版：xelatex "\def\isanswer{1}\input{高一59_华实高中_10月月考.tex}"
% 紧凑版：xelatex "\def\iscompact{1}\input{高一59_华实高中_10月月考.tex}"
%
% 原始材料：微信公众号“魔都数学加油站”文章，作者破晓，连载“27学年高一59”；
% 原文标题“27学年高一59——2026-2027年上海市华实高中高一上10月月考”；
% 原文链接：https://mp.weixin.qq.com/s/XzyBCmrvpRuMe-NuNk7gkA。
% 页面标注发布时间 2026-10-13 21:45（上海）。
% 正文为 9 张 PNG 页：第 01、02、03、05、06、07、08、09 页 4762×6735，第 04 页 2560×3620
% （**同一篇各页分辨率不同**），均按页序读取。9 页全为试卷正文页（卷首标题在第 1 页页顶，
% 卷末解答题第 21 题解析止于第 9 页），无封面页、目录页、二维码或宣传页。
% 粉色斜排“魔都数学加油站”为卷面水印，与公众号页脚一并未录入。
%
% 卷首标题：照卷面印作“2026--2027 年华实高中高一上 10 月月考”（公众号标题中的“上海市”
% 卷面标题行没有，本稿照卷面），年份区间按全稿惯例排 en dash。原卷只有一行标题、无副标题。
%
% 篇幅与题号：原卷自第 3 题起编号（第 1、2 题不在本卷内）。填空题 3—12（10 题）、
% 选择题 13—16（4 题）、解答题 17—21（5 题），共 19 题。本稿沿用原节与原题号，
% 只改排为全稿统一的大节样式（原卷小节标题为“一、填空题”“二、选择题”“三、解答题”）。
%
% 讲评版：原卷每题下直接印【答案】或【解析】。**只印【答案】的 1 题**（第 3 题），本稿只给 \ans{}；
% 其余 18 题都印【解析】，按全稿口径这类题不另提 \ans{}——第 13、14、15、16 题的结论写在解析末句
% （“故选 D”等），亦照录在解析内。本卷**没有**印在题干括号内的答案。
%
% 副标题：原卷未印副标题。本稿按“知识模块”口径标作“集合与常用逻辑用语、不等式”：
% 19 题中集合与常用逻辑用语 13 题（第 3、9、10、11、12、14、15、16、17、18、19、20、21 题），
% 不等式 10 题（含等式与不等式性质、一元二次不等式、恒成立、基本应用：第 4、5、6、7、8、
% 13、17、18、19、20 题），两类均远超三分之一，故并标。口径同上。
%
% 与原卷的排版差异：
%   ① 原卷每题下直接印【答案】/【解析】，本稿按全稿惯例将答案与解析置于答案版，试卷版保持干净；
%   ② 原卷未印副标题，本稿按“知识模块”口径补出；
%   ③ 原卷填空横线约两个字宽，本稿按全稿惯例统一排 \blankline[4.4em]；
%   ④ 原卷实数集排粗体 $\mathbf{R}$（第 12、13、20 题等），整数集排粗体 $\mathbf{Z}$（第 11、12 题），
%      自然数集排斜体 $N$ 或粗体 $\mathbf{Z}$ 处按本稿惯例改用 $\R$、$\Z$、$\N$；
%      空集记号统一为 \varnothing；
%   ⑤ 原卷在数学式内多处用半角逗号（如 $x_1,x_2$、$a,b\in R$），均照录未强改；
%      句末句点按全稿惯例排西文句点；“漂亮集”“上位点”“下位点”等自定义名词照录；
%   ⑥ 第 17—21 题的子问原卷各自另起一行，本稿按全稿惯例用 enumerate 的 (1)(2)(3)；
%      解析中的子问标号一律用 \stepm{(1)}{…}；第 9、11 题解析中的①、②（③④）分情形或分结论，
%      按全稿惯例作 \stepm{①}{…}、\stepm{②}{…}（第 11 题解析末“故正确的是②④”照录为行内）；
%   ⑦ 第 12 题题干里的“$A_p$”下标原卷作斜体 $p$，照录；
%   ⑧ 每道解答题原卷未留作答空白，本稿按全稿惯例加 \ansspace。
%
% 原卷笔误/存疑（均照录未改，正文对应处另有 % 注）：本卷首轮逐页转录，暂未发现影响答案的笔误。
%
% 插图：无（全卷 9 页无任何图形、无“如图”字样）。
%
\documentclass[a4paper,11pt]{article}
\usepackage{common}

\begin{document}

\examtitlest[3]{2026--2027 年华实高中高一上 10 月月考}{集合与常用逻辑用语、不等式}

\bigsec{一、填空题}

\begin{enumerate}
\setcounter{enumi}{2}

\item “$x\geqslant1$ 且 $y\geqslant1$”的否定形式为\blankline[4.4em].

\ifanswer
\ans{$x<1$ 或 $y<1$}
\fi

\item 已知方程 $2x^2+4x-3=0$ 的两个根为 $x_1,x_2$，则 $\dfrac1{x_1}+\dfrac1{x_2}=$\blankline[4.4em].

\ifanswer
\solbegin
\step{利用韦达定理得 $\begin{cases}x_1+x_2=-2\\ x_1\cdot x_2=-\dfrac32\end{cases}$，则 $\dfrac1{x_1}+\dfrac1{x_2}=\dfrac{x_1+x_2}{x_1\cdot x_2}=\dfrac43$.}
\solend
\fi

\item 设 $a>b>0$，比较大小：$\dfrac1{a-b}$\blankline[2.4em]$\dfrac1a$.

\ifanswer
\solbegin
\step{由 $a>b>0$ 得 $a>0$ 且 $a-b>0$，所以 $\dfrac1{a-b}-\dfrac1a=\dfrac b{a(a-b)}>0$，}
\step{故 $\dfrac1{a-b}>\dfrac1a$.}
\solend
\fi

\item 若不等式 $2ax^2-ax+1>0$ 的解集为 $R$，则实数 $a$ 的取值范围是\blankline[4.4em].

\ifanswer
\solbegin
\step{当 $a\neq0$ 时，由题意得 $\begin{cases}a>0\\ \Delta<0\end{cases}$，即 $\begin{cases}a>0\\ a^2-8a<0\end{cases}$，解得 $0<a<8$.}
\step{当 $a=0$ 时，恒有 $1>0$，不等式也成立.}
\step{综上所述，$a$ 的取值范围是 $[0,8)$.}
\solend
\fi

\item 集合 $A=\{x\mid m-1\leqslant x\leqslant m+7\}$，集合 $B=\{x\mid x^2-3x-10\geqslant0\}$，若 $A\cup B=R$，则实数 $m$ 的取值范围为\blankline[4.4em].

\ifanswer
\solbegin
\step{$A=\{x\mid m-1\leqslant x\leqslant m+7\}$，$B=\{x\mid x^2-3x-10\geqslant0\}=\{x\mid x\leqslant-2\text{ 或 }x\geqslant5\}$，}
\step{若 $A\cup B=R$，则 $\begin{cases}m-1\leqslant-2\\ m+7\geqslant5\end{cases}$，解得 $-2\leqslant m\leqslant-1$，}
\step{所以实数 $m$ 的取值范围为 $[-2,-1]$.}
\solend
\fi

\item 若关于 $x$ 的不等式 $ax^2+bx+c>0$ 的解集为 $(-1,3)$，则关于 $x$ 的不等式 $cx^2+bx+a>0$ 的解集为\blankline[4.4em].

\ifanswer
\solbegin
\step{因为不等式 $ax^2+bx+c>0$ 的解集为 $(-1,3)$，}
\step{所以 $-1$ 和 $3$ 是方程 $ax^2+bx+c=0$ 的实数解，}
\step{所以 $\begin{cases}-1+3=-\dfrac ba\\ -1\times3=\dfrac ca\\ a<0\end{cases}$，解得 $b=-2a$，$c=-3a$，且 $a<0$，}
\step{所以不等式 $cx^2+bx+a>0$ 可化为 $-3ax^2-2ax+a>0$，即 $3x^2+2x-1>0$，}
\step{解得 $x<-1$ 或 $x>\dfrac13$，所以该不等式的解集为 $(-\infty,-1)\cup\left(\dfrac13,+\infty\right)$.}
\solend
\fi

\item 若集合 $A=\{x\mid(a-1)x^2-2x+1=0\}$ 有且仅有 $2$ 个子集，则实数 $a$ 的值为\blankline[4.4em].

\ifanswer
\solbegin
\step{由题意得方程 $(a-1)x^2-2x+1=0$ 只有一个根，}
\stepm{①}{当 $a-1=0$ 时，$a=1$，此时方程化为 $-2x+1=0$，解得 $x=\dfrac12$，符合题意，}
\stepm{②}{当 $a-1\neq0$，即 $a\neq1$ 时，$\Delta=4-4(a-1)=0$，解得 $a=2$，}
\step{综上所述，实数 $a$ 的值为 $1$ 或 $2$.}
\solend
\fi

\item “$x+y\neq5$”是“$x\neq3$ 或 $y\neq2$”的\blankline[4.4em]条件.

\ifanswer
\solbegin
\step{分析充分性：}
\step{若 $x+y\neq5$ 成立，若 $x=3$，则 $y\neq2$，若 $y=2$，则 $x\neq3$，}
\step{所以一定能说明 $x\neq3$ 和 $y\neq2$ 至少有一个成立，即 $x\neq3$ 或 $y\neq2$ 成立，}
\step{充分性成立；}
\step{分析必要性：}
\step{若 $x\neq3$ 或 $y\neq2$，即 $x\neq3$ 和 $y\neq2$ 二者至少一个成立，}
\step{举例：$x=1,y=4$，此时 $x+y=5$，所以必要性不成立.}
\step{所以“$x+y\neq5$”是“$x\neq3$ 或 $y\neq2$”的充分不必要条件.}
\solend
\fi

\item 设 $P$ 为非空实数集满足：对任意给定的 $x,y\in P$（$x,y$ 可以相同），都有 $x+y\in P,x-y\in P,xy\in P$，则称 $P$ 为漂亮集. 关于漂亮集有下列结论：

①集合 $P=\{-2,-1,0,1,2\}$ 为漂亮集；②集合 $P=\{x\mid x=2n,n\in\Z\}$ 为漂亮集；

③若集合 $P_1,P_2$ 为漂亮集，则 $P_1\cup P_2$ 为漂亮集；④若集合 $P$ 为漂亮集，则一定有 $0\in P$.

其中所有正确结论的序号是\blankline[4.4em].

\ifanswer
\solbegin
\step{对于集合 $P=\{-2,-1,0,1,2\}$，取 $x=2,y=-2$，则 $x-y=2-(-2)=4\notin P$，}
\step{不满足漂亮集的定义，所以①错误.}
\step{对于集合 $P=\{x\mid x=2n,n\in Z\}$，设 $x=2n_1,y=2n_2,n_1,n_2\in Z$.}
\step{$x+y=2n_1+2n_2=2(n_1+n_2)$，因为 $n_1+n_2\in Z$，所以 $x+y\in P$.}
\step{$x-y=2n_1-2n_2=2(n_1-n_2)$，因为 $n_1-n_2\in Z$，所以 $x-y\in P$.}
\step{$xy=2n_1\times2n_2=4n_1n_2=2(2n_1n_2)$，因为 $2n_1n_2\in Z$，所以 $xy\in P$.}
\step{满足漂亮集的定义，所以②正确.}
\step{令 $P_1=\{x\mid x=2n,n\in Z\}$,$P_2=\{x\mid x=3n,n\in Z\}$，都是漂亮集.}
\step{$P_1\cup P_2$ 中，取 $x=2(x\in P_1)$,$y=3(y\in P_2)$,$x+y=5\notin P_1\cup P_2$，}
\step{不满足漂亮集的定义，所以③错误.}
\step{因为对任意 $x\in P$，有 $x-x=0\in P$，所以若集合 $P$ 为漂亮集，}
\step{则一定有 $0\in P$，所以④正确.}
\step{故正确的是②④.}
\solend
\fi

\item 已知：集合 $A_p=\{(x,y)\mid y=3x^2-px+5p-65,x\in\R\}$,$M=\{x\mid1\leqslant x\leqslant100,x\in\Z\}$，若 $p\in M$，则所有 $A_p$ 的交集\blankline[4.4em].

\ifanswer
\solbegin
\step{由 $M=\{1,2,\cdots,100\}$，所以所有 $A_p$ 的交集就是这些集合中共同拥有的点.}
\step{若 $(x,y)$ 是所有集合的一个公共点，}
\step{那么它特别要同时属于 $p=1$ 和 $p=2$ 时对应的两个集合.}
\step{当 $p=1$ 时，$y=3x^2-x-60$；}
\step{当 $p=2$ 时，$y=3x^2-2x-55$；联立得 $x=5,y=10$.}
\step{因此，如果公共点存在，只可能是 $(5,10)$.}
\step{接下来验证这个点是否属于所有集合，}
\step{对任意 $p\in M$，把 $x=5$ 代入集合的关系式，}
\step{得 $3x^2-px+5p-65=75-5p+5p-65=10$，与 $y=10$ 相符，}
\step{所以 $(5,10)\in A_p$ 对每一个 $p$ 都成立，因此所有 $A_p$ 的交集为 $\{(5,10)\}$.}
\solend
\fi

\end{enumerate}

% 守卫须放在 \bigsec 之前，否则标题仍会落页末（见 高一40 的教训）
\needvspace{5cm}

\bigsec{二、选择题}

\begin{enumerate}
\setcounter{enumi}{12}

\item 已知 $a,b,c\in\R$ 且 $a>b$，则下列不等式一定成立的是\mbox{（\quad）}

\keepwithprev
A. $\dfrac1a<\dfrac1b$\qquad B. $a^2>b^2$

C. $a|c|>b|c|$\qquad D. $\dfrac a{c^2+1}>\dfrac b{c^2+1}$

\ifanswer
\solbegin
\step{$A$：当 $a>0>b$ 时，$\dfrac1a>\dfrac1b$，故 $A$ 错误；}
\step{$B$：当 $a=-1,b=-2$ 时，满足 $a>b$，但 $a^2>b^2$ 不成立，故 $B$ 错误；}
\step{$C$：当 $c=0$ 时，$a|c|=b|c|$，故 $C$ 错误；}
\step{$D$：由 $a>b,c^2+1>0$，得 $\dfrac a{c^2+1}>\dfrac b{c^2+1}$，故 $D$ 正确；}
\step{故选 $D$.}
\solend
\fi

\item 已知集合 $A=\{a^2,0,-1\}$，$B=\{a,b,0\}$，若 $A=B$，则 $(ab)^{2026}$ 的值为\mbox{（\quad）}

\keepwithprev
A. $0$\qquad B. $-1$

C. $1$\qquad D. $\pm1$

\ifanswer
\solbegin
\step{因为集合 $A=\{a^2,0,-1\}$，$B=\{a,b,0\}$，$A=B$，所以 $\begin{cases}a=-1\\ a^2=b\end{cases}$，}
\step{所以 $a=-1$，$b=1$，所以 $(ab)^{2026}=1$，故选 $C$.}
\solend
\fi

\item 设全集 $U=\R$，$A=\{x\mid x<-1\text{ 或 }x>2\}$,$B=\{y\mid y=x^2+a\}$，若 $\overline A\subseteq B$，则实数 $a$ 的取值范围是\mbox{（\quad）}

\keepwithprev
A. $a\leqslant-1$\qquad B. $a>2$

C. $a<-1$\qquad D. $a\geqslant-2$

\ifanswer
\solbegin
\step{因为 $A=\{x\mid x<-1\text{ 或 }x>2\}$，则 $\overline A=\{x\mid-1\leqslant x\leqslant2\}$，}
\step{又 $B=\{y\mid y=x^2+a\}=\{y\mid y\geqslant a\}$，且 $\overline A\subseteq B$，所以 $a\leqslant-1$，故选 $A$.}
\solend
\fi

\item 设 $P$ 和 $Q$ 是两个集合，定义集合 $P-Q=\{x\mid x\in P$，且 $x\notin Q\}$，如果 $Q=\{x\mid(1-x)(x-3)>0\}$，$P=\left\{x\left|\dfrac x{2-x}>0\right.\right\}$，那么 $P-Q$ 等于\mbox{（\quad）}

\keepwithprev
A. $\{x\mid0<x<1\}$\qquad B. $\{x\mid1\leqslant x<2\}$

C. $\{x\mid0<x\leqslant1\}$\qquad D. $\{x\mid2\leqslant x<3\}$

\ifanswer
\solbegin
\step{因为 $P=\left\{x\left|\dfrac x{2-x}>0\right.\right\}=\{x\mid x(x-2)<0\}=\{x\mid0<x<2\}$，}
\step{$Q=\{x\mid(1-x)(x-3)>0\}=\{x\mid(x-3)(x-1)<0\}=\{x\mid1<x<3\}$，}
\step{所以 $P-Q=\{x\mid0<x\leqslant1\}$，故选 $C$.}
\solend
\fi

\end{enumerate}

\bigsec{三、解答题}

\begin{enumerate}
\setcounter{enumi}{16}

\item 已知集合 $A=\{x\mid x^2-9\geqslant0\}$，$B=\{x\mid|x-4|<2\}$，$C=\left\{x\left|\dfrac{x-8}{x+2}<0\right.\right\}$.

（1）求 $A\cup C$；

（2）若全集 $U=R$，求 $\overline A\cap B$.

\ifanswer
\solbegin
\stepm{(1)}{集合 $A=\{x\mid x^2-9\geqslant0\}=\{x\mid x\geqslant3\text{ 或 }x\leqslant-3\}$，}
\step{$B=\{x\mid|x-4|<2\}=\{x\mid2<x<6\}$，$C=\left\{x\left|\dfrac{x-8}{x+2}<0\right.\right\}=\{x\mid-2<x<8\}$，}
\step{所以 $A\cup C=\{x\mid x\leqslant-3\text{ 或 }x>-2\}$；}
\stepm{(2)}{全集 $U=R$，集合 $A=\{x\mid x^2-9\geqslant0\}=\{x\mid x\geqslant3\text{ 或 }x\leqslant-3\}$，}
\step{$B=\{x\mid|x-4|<2\}=\{x\mid2<x<6\}$，}
\step{所以 $\overline A=\{x\mid-3<x<3\}$，所以 $\overline A\cap B=\{x\mid2<x<3\}$.}
\solend
\fi

\ansspace[6]

\item 已知集合 $A=\left\{x\left|\dfrac{2x-1}{x+1}<1,x\in\R\right.\right\}$，集合 $B=\{x\mid(x-a-1)(x-a+1)\leqslant0\}$.

（1）求集合 $A$；

（2）若 $B\subseteq A$，求实数 $a$ 的取值范围.

\ifanswer
\solbegin
\stepm{(1)}{由 $\dfrac{2x-1}{x+1}<1$，得 $\dfrac{x-2}{x+1}<0$ 所以 $(x-2)(x+1)<0$，解得 $-1<x<2$，}
\step{即 $A=(-1,2)$；}
\stepm{(2)}{由 $(x-a-1)(x-a+1)\leqslant0$，得 $a-1\leqslant x\leqslant a+1$，}
\step{所以 $B=\{x\mid a-1\leqslant x\leqslant a+1\}$，}
\step{又 $B\subseteq A$，且 $a+1>a-1$ 即 $B\neq\varnothing$，所以 $\begin{cases}a-1>-1\\ a+1<2\end{cases}$，所以 $0<a<1$，}
\step{所以实数 $a$ 的取值范围是 $\{a\mid0<a<1\}$.}
\solend
\fi

\ansspace[6]

\item 已知集合 $A=\{x\mid x^2-4x+3=0\}$，$B=\{x\mid x^2-ax+3=0\}$.

（1）若 $A\cup B=B$，求实数 $a$ 的值；

（2）若 $A\cap B=B$，求实数 $a$ 的取值范围.

\ifanswer
\solbegin
\stepm{(1)}{集合 $A=\{x\mid x^2-4x+3=0\}=\{1,3\}$，$B=\{x\mid x^2-ax+3=0\}$.}
\step{$A\cup B=B$，所以 $1\in B$，$3\in B$，所以 $\begin{cases}1-a+3=0\\ 9-3a+3=0\end{cases}$，解得实数 $a=4$；}
\stepm{(2)}{因为 $A\cap B=B$，所以 $B\subseteq A$，}
\step{当 $B=\varnothing$ 时，$\Delta=(-a)^2-12<0$，解得 $-2\sqrt3<a<2\sqrt3$，}
\step{当 $B$ 是单元素集时，$\Delta=a^2-12=0$，解得 $a=2\sqrt3$，此时 $B=\{\sqrt3\}$，}
\step{不合题意；}
\step{当 $B=\{1,3\}$ 时，由 (1) 得 $a=4$，}
\step{综上，实数 $a$ 的取值范围是 $(-2\sqrt3,2\sqrt3)\cup\{4\}$.}
\solend
\fi

\ansspace[6]

\item 已知全集为 $R$，实数 $a$、$b$ 满足 $a\neq b$，集合 $M=\{x\mid x^2-3x-4\leqslant0\}$，集合 $A=\{x\mid(x+a)(x+b)>0\}$，集合 $B=\{x\mid x^2+(a-2)x-2a>0\}$.

（1）若 $\overline A=M$，求 $a$、$b$ 的值；

（2）若 $a>b>-1$，求 $A\cap B$；

（3）若 $a^2+a\in\overline B$，求 $a$ 的取值范围.

\ifanswer
\solbegin
\stepm{(1)}{$M=\{x\mid x^2-3x-4\leqslant0\}=\{x\mid-1\leqslant x\leqslant4\}$，$A=\{x\mid(x+a)(x+b)\leqslant0\}$，}
\step{因为 $\overline A=M$，所以 $a=1$，$b=-4$ 或 $a=-4$，$b=1$.}
\stepm{(2)}{因为 $a>b>-1$，所以 $-a<-b<1$，所以 $A=\{x\mid x<-a\text{ 或 }x>-b\}$，}
\step{因为 $B=\{x\mid x^2+(a-2)x-2a>0\}=\{x\mid x<-a\text{ 或 }x>2\}$，}
\step{所以 $A\cap B=\{x\mid x<-a\text{ 或 }x>2\}$，}
\stepm{(3)}{$\overline B=\{x\mid(x-2)(x+a)\leqslant0\}$，}
\step{由 $a^2+a\in\overline B$，得 $(a^2+a-2)(a^2+2a)\leqslant0\Leftrightarrow a(a-1)(a-2)^2\leqslant0$，}
\step{解得 $a=2$ 或 $0\leqslant a\leqslant1$，故 $a$ 的取值范围是 $[0,1]\cup\{2\}$.}
\solend
\fi

\ansspace[8]

\item 对平面直角坐标系的第一象限内的任意两点作如下定义：若 $\dfrac ab>\dfrac cd$，那么称点 $(a,b)$ 是点 $(c,d)$ 的“上位点”，同时点 $(c,d)$ 是点 $(a,b)$ 的“下位点”.

（1）试写出点 $(2,4)$ 的一个“上位点”坐标；

（2）已知点 $(a,b)$ 是点 $(c,d)$ 的“上位点”，判断点 $P(a+c,b+d)$ 是点 $(c,d)$ 的“上位点”还是“下位点”，并证明你的结论；

（3）设正整数 $n$ 满足以下条件：对集合 $\{t\mid0<t<2020,t\in\Z\}$ 内的任意元素 $m$，总存在正整数 $k=2m+1$，使得点 $(n,k)$ 既是点 $(2025,m)$ 的“下位点”，又是点 $(2026,m+1)$ 的“上位点”，问这样的正整数 $n$ 是否存在？若存在，求正整数 $n$ 的值；若不存在，请说明理由.

\ifanswer
\solbegin
\stepm{(1)}{由定义，点 $(a,b)$ 是点 $(2,4)$ 的“上位点”，}
\step{当且仅当 $\dfrac ab>\dfrac24=\dfrac12$ 且 $a>0,b>0$.}
\step{取 $(a,b)=(1,1)$，因为 $\dfrac11>\dfrac12$，}
\step{所以点 $(1,1)$ 是点 $(2,4)$ 的一个“上位点”.（答案不唯一）}
\stepm{(2)}{点 $P(a+c,b+d)$ 是点 $(c,d)$ 的“上位点”.}
\step{$(a,b)$ 是点 $(c,d)$ 的“上位点”，所以 $\dfrac ab>\dfrac cd$.}
\step{因为 $a,b,c,d>0$，两边同乘 $bd$ 得 $ad>bc$，即 $ad-bc>0$.}
\step{作差比较 $\dfrac{a+c}{b+d}$ 与 $\dfrac cd$：$\dfrac{a+c}{b+d}-\dfrac cd=\dfrac{d(a+c)-c(b+d)}{d(b+d)}=\dfrac{ad-bc}{d(b+d)}>0$，}
\step{所以点 $P(a+c,b+d)$ 是点 $(c,d)$ 的“上位点”.}
\stepm{(3)}{由定义，点 $(n,k)$ 是点 $(2025,m)$ 的下位点即 $\dfrac{2025}m>\dfrac nk$；}
\step{点 $(n,k)$ 是点 $(2026,m+1)$ 的上位点即 $\dfrac nk>\dfrac{2026}{m+1}$.}
\step{代入 $k=2m+1$，得对任意 $m\in\{1,2,\cdots,2019\}$，}
\step{都有 $4052-\dfrac{2026}{m+1}<n<4050+\dfrac{2025}m$.}
\step{右边 $4050+\dfrac{2025}m$ 随 $m$ 增大而减小，故最近处在 $m=2019$，}
\step{得 $n<4050+\dfrac{2025}{2019}<4051.01$，所以 $n\leqslant4051$.}
\step{左边 $4052-\dfrac{2026}{m+1}$ 随 $m$ 增大而增大，故最近处也在 $m=2019$，}
\step{得 $n>4052-\dfrac{2026}{2020}>4050.99$，所以 $n\geqslant4051$.}
\step{综上，$n=4051$.}
\step{检验：当 $n=4051$ 时，对任意 $1\leqslant m\leqslant2019,k=2m+1$ 是正整数，}
\step{且 $\dfrac{4051}{2m+1}<\dfrac{2025}m$,$4051m<4050m+2025,m<2025$，成立；}
\step{$\dfrac{4051}{2m+1}>\dfrac{2026}{m+1}$,$4051(m+1)>4052m+2026,m<2025$，成立.}
\step{又各点横纵坐标均为正，都在第一象限内，符合定义，}
\step{所以满足条件的正整数 $n$ 存在，其值为 $n=4051$.}
\solend
\fi

\ansspace*[10]

\end{enumerate}

\end{document}
"
% 紧凑版：xelatex "\def\iscompact{1}% 高一59_华实高中_10月月考.tex
%   —— 高一上数学“2026-2027 年华实高中高一上 10 月月考”（讲评版）
% 试卷版：xelatex 高一59_华实高中_10月月考.tex
% 答案版：xelatex "\def\isanswer{1}\input{高一59_华实高中_10月月考.tex}"
% 紧凑版：xelatex "\def\iscompact{1}\input{高一59_华实高中_10月月考.tex}"
%
% 原始材料：微信公众号“魔都数学加油站”文章，作者破晓，连载“27学年高一59”；
% 原文标题“27学年高一59——2026-2027年上海市华实高中高一上10月月考”；
% 原文链接：https://mp.weixin.qq.com/s/XzyBCmrvpRuMe-NuNk7gkA。
% 页面标注发布时间 2026-10-13 21:45（上海）。
% 正文为 9 张 PNG 页：第 01、02、03、05、06、07、08、09 页 4762×6735，第 04 页 2560×3620
% （**同一篇各页分辨率不同**），均按页序读取。9 页全为试卷正文页（卷首标题在第 1 页页顶，
% 卷末解答题第 21 题解析止于第 9 页），无封面页、目录页、二维码或宣传页。
% 粉色斜排“魔都数学加油站”为卷面水印，与公众号页脚一并未录入。
%
% 卷首标题：照卷面印作“2026--2027 年华实高中高一上 10 月月考”（公众号标题中的“上海市”
% 卷面标题行没有，本稿照卷面），年份区间按全稿惯例排 en dash。原卷只有一行标题、无副标题。
%
% 篇幅与题号：原卷自第 3 题起编号（第 1、2 题不在本卷内）。填空题 3—12（10 题）、
% 选择题 13—16（4 题）、解答题 17—21（5 题），共 19 题。本稿沿用原节与原题号，
% 只改排为全稿统一的大节样式（原卷小节标题为“一、填空题”“二、选择题”“三、解答题”）。
%
% 讲评版：原卷每题下直接印【答案】或【解析】。**只印【答案】的 1 题**（第 3 题），本稿只给 \ans{}；
% 其余 18 题都印【解析】，按全稿口径这类题不另提 \ans{}——第 13、14、15、16 题的结论写在解析末句
% （“故选 D”等），亦照录在解析内。本卷**没有**印在题干括号内的答案。
%
% 副标题：原卷未印副标题。本稿按“知识模块”口径标作“集合与常用逻辑用语、不等式”：
% 19 题中集合与常用逻辑用语 13 题（第 3、9、10、11、12、14、15、16、17、18、19、20、21 题），
% 不等式 10 题（含等式与不等式性质、一元二次不等式、恒成立、基本应用：第 4、5、6、7、8、
% 13、17、18、19、20 题），两类均远超三分之一，故并标。口径同上。
%
% 与原卷的排版差异：
%   ① 原卷每题下直接印【答案】/【解析】，本稿按全稿惯例将答案与解析置于答案版，试卷版保持干净；
%   ② 原卷未印副标题，本稿按“知识模块”口径补出；
%   ③ 原卷填空横线约两个字宽，本稿按全稿惯例统一排 \blankline[4.4em]；
%   ④ 原卷实数集排粗体 $\mathbf{R}$（第 12、13、20 题等），整数集排粗体 $\mathbf{Z}$（第 11、12 题），
%      自然数集排斜体 $N$ 或粗体 $\mathbf{Z}$ 处按本稿惯例改用 $\R$、$\Z$、$\N$；
%      空集记号统一为 \varnothing；
%   ⑤ 原卷在数学式内多处用半角逗号（如 $x_1,x_2$、$a,b\in R$），均照录未强改；
%      句末句点按全稿惯例排西文句点；“漂亮集”“上位点”“下位点”等自定义名词照录；
%   ⑥ 第 17—21 题的子问原卷各自另起一行，本稿按全稿惯例用 enumerate 的 (1)(2)(3)；
%      解析中的子问标号一律用 \stepm{(1)}{…}；第 9、11 题解析中的①、②（③④）分情形或分结论，
%      按全稿惯例作 \stepm{①}{…}、\stepm{②}{…}（第 11 题解析末“故正确的是②④”照录为行内）；
%   ⑦ 第 12 题题干里的“$A_p$”下标原卷作斜体 $p$，照录；
%   ⑧ 每道解答题原卷未留作答空白，本稿按全稿惯例加 \ansspace。
%
% 原卷笔误/存疑（均照录未改，正文对应处另有 % 注）：本卷首轮逐页转录，暂未发现影响答案的笔误。
%
% 插图：无（全卷 9 页无任何图形、无“如图”字样）。
%
\documentclass[a4paper,11pt]{article}
\usepackage{common}

\begin{document}

\examtitlest[3]{2026--2027 年华实高中高一上 10 月月考}{集合与常用逻辑用语、不等式}

\bigsec{一、填空题}

\begin{enumerate}
\setcounter{enumi}{2}

\item “$x\geqslant1$ 且 $y\geqslant1$”的否定形式为\blankline[4.4em].

\ifanswer
\ans{$x<1$ 或 $y<1$}
\fi

\item 已知方程 $2x^2+4x-3=0$ 的两个根为 $x_1,x_2$，则 $\dfrac1{x_1}+\dfrac1{x_2}=$\blankline[4.4em].

\ifanswer
\solbegin
\step{利用韦达定理得 $\begin{cases}x_1+x_2=-2\\ x_1\cdot x_2=-\dfrac32\end{cases}$，则 $\dfrac1{x_1}+\dfrac1{x_2}=\dfrac{x_1+x_2}{x_1\cdot x_2}=\dfrac43$.}
\solend
\fi

\item 设 $a>b>0$，比较大小：$\dfrac1{a-b}$\blankline[2.4em]$\dfrac1a$.

\ifanswer
\solbegin
\step{由 $a>b>0$ 得 $a>0$ 且 $a-b>0$，所以 $\dfrac1{a-b}-\dfrac1a=\dfrac b{a(a-b)}>0$，}
\step{故 $\dfrac1{a-b}>\dfrac1a$.}
\solend
\fi

\item 若不等式 $2ax^2-ax+1>0$ 的解集为 $R$，则实数 $a$ 的取值范围是\blankline[4.4em].

\ifanswer
\solbegin
\step{当 $a\neq0$ 时，由题意得 $\begin{cases}a>0\\ \Delta<0\end{cases}$，即 $\begin{cases}a>0\\ a^2-8a<0\end{cases}$，解得 $0<a<8$.}
\step{当 $a=0$ 时，恒有 $1>0$，不等式也成立.}
\step{综上所述，$a$ 的取值范围是 $[0,8)$.}
\solend
\fi

\item 集合 $A=\{x\mid m-1\leqslant x\leqslant m+7\}$，集合 $B=\{x\mid x^2-3x-10\geqslant0\}$，若 $A\cup B=R$，则实数 $m$ 的取值范围为\blankline[4.4em].

\ifanswer
\solbegin
\step{$A=\{x\mid m-1\leqslant x\leqslant m+7\}$，$B=\{x\mid x^2-3x-10\geqslant0\}=\{x\mid x\leqslant-2\text{ 或 }x\geqslant5\}$，}
\step{若 $A\cup B=R$，则 $\begin{cases}m-1\leqslant-2\\ m+7\geqslant5\end{cases}$，解得 $-2\leqslant m\leqslant-1$，}
\step{所以实数 $m$ 的取值范围为 $[-2,-1]$.}
\solend
\fi

\item 若关于 $x$ 的不等式 $ax^2+bx+c>0$ 的解集为 $(-1,3)$，则关于 $x$ 的不等式 $cx^2+bx+a>0$ 的解集为\blankline[4.4em].

\ifanswer
\solbegin
\step{因为不等式 $ax^2+bx+c>0$ 的解集为 $(-1,3)$，}
\step{所以 $-1$ 和 $3$ 是方程 $ax^2+bx+c=0$ 的实数解，}
\step{所以 $\begin{cases}-1+3=-\dfrac ba\\ -1\times3=\dfrac ca\\ a<0\end{cases}$，解得 $b=-2a$，$c=-3a$，且 $a<0$，}
\step{所以不等式 $cx^2+bx+a>0$ 可化为 $-3ax^2-2ax+a>0$，即 $3x^2+2x-1>0$，}
\step{解得 $x<-1$ 或 $x>\dfrac13$，所以该不等式的解集为 $(-\infty,-1)\cup\left(\dfrac13,+\infty\right)$.}
\solend
\fi

\item 若集合 $A=\{x\mid(a-1)x^2-2x+1=0\}$ 有且仅有 $2$ 个子集，则实数 $a$ 的值为\blankline[4.4em].

\ifanswer
\solbegin
\step{由题意得方程 $(a-1)x^2-2x+1=0$ 只有一个根，}
\stepm{①}{当 $a-1=0$ 时，$a=1$，此时方程化为 $-2x+1=0$，解得 $x=\dfrac12$，符合题意，}
\stepm{②}{当 $a-1\neq0$，即 $a\neq1$ 时，$\Delta=4-4(a-1)=0$，解得 $a=2$，}
\step{综上所述，实数 $a$ 的值为 $1$ 或 $2$.}
\solend
\fi

\item “$x+y\neq5$”是“$x\neq3$ 或 $y\neq2$”的\blankline[4.4em]条件.

\ifanswer
\solbegin
\step{分析充分性：}
\step{若 $x+y\neq5$ 成立，若 $x=3$，则 $y\neq2$，若 $y=2$，则 $x\neq3$，}
\step{所以一定能说明 $x\neq3$ 和 $y\neq2$ 至少有一个成立，即 $x\neq3$ 或 $y\neq2$ 成立，}
\step{充分性成立；}
\step{分析必要性：}
\step{若 $x\neq3$ 或 $y\neq2$，即 $x\neq3$ 和 $y\neq2$ 二者至少一个成立，}
\step{举例：$x=1,y=4$，此时 $x+y=5$，所以必要性不成立.}
\step{所以“$x+y\neq5$”是“$x\neq3$ 或 $y\neq2$”的充分不必要条件.}
\solend
\fi

\item 设 $P$ 为非空实数集满足：对任意给定的 $x,y\in P$（$x,y$ 可以相同），都有 $x+y\in P,x-y\in P,xy\in P$，则称 $P$ 为漂亮集. 关于漂亮集有下列结论：

①集合 $P=\{-2,-1,0,1,2\}$ 为漂亮集；②集合 $P=\{x\mid x=2n,n\in\Z\}$ 为漂亮集；

③若集合 $P_1,P_2$ 为漂亮集，则 $P_1\cup P_2$ 为漂亮集；④若集合 $P$ 为漂亮集，则一定有 $0\in P$.

其中所有正确结论的序号是\blankline[4.4em].

\ifanswer
\solbegin
\step{对于集合 $P=\{-2,-1,0,1,2\}$，取 $x=2,y=-2$，则 $x-y=2-(-2)=4\notin P$，}
\step{不满足漂亮集的定义，所以①错误.}
\step{对于集合 $P=\{x\mid x=2n,n\in Z\}$，设 $x=2n_1,y=2n_2,n_1,n_2\in Z$.}
\step{$x+y=2n_1+2n_2=2(n_1+n_2)$，因为 $n_1+n_2\in Z$，所以 $x+y\in P$.}
\step{$x-y=2n_1-2n_2=2(n_1-n_2)$，因为 $n_1-n_2\in Z$，所以 $x-y\in P$.}
\step{$xy=2n_1\times2n_2=4n_1n_2=2(2n_1n_2)$，因为 $2n_1n_2\in Z$，所以 $xy\in P$.}
\step{满足漂亮集的定义，所以②正确.}
\step{令 $P_1=\{x\mid x=2n,n\in Z\}$,$P_2=\{x\mid x=3n,n\in Z\}$，都是漂亮集.}
\step{$P_1\cup P_2$ 中，取 $x=2(x\in P_1)$,$y=3(y\in P_2)$,$x+y=5\notin P_1\cup P_2$，}
\step{不满足漂亮集的定义，所以③错误.}
\step{因为对任意 $x\in P$，有 $x-x=0\in P$，所以若集合 $P$ 为漂亮集，}
\step{则一定有 $0\in P$，所以④正确.}
\step{故正确的是②④.}
\solend
\fi

\item 已知：集合 $A_p=\{(x,y)\mid y=3x^2-px+5p-65,x\in\R\}$,$M=\{x\mid1\leqslant x\leqslant100,x\in\Z\}$，若 $p\in M$，则所有 $A_p$ 的交集\blankline[4.4em].

\ifanswer
\solbegin
\step{由 $M=\{1,2,\cdots,100\}$，所以所有 $A_p$ 的交集就是这些集合中共同拥有的点.}
\step{若 $(x,y)$ 是所有集合的一个公共点，}
\step{那么它特别要同时属于 $p=1$ 和 $p=2$ 时对应的两个集合.}
\step{当 $p=1$ 时，$y=3x^2-x-60$；}
\step{当 $p=2$ 时，$y=3x^2-2x-55$；联立得 $x=5,y=10$.}
\step{因此，如果公共点存在，只可能是 $(5,10)$.}
\step{接下来验证这个点是否属于所有集合，}
\step{对任意 $p\in M$，把 $x=5$ 代入集合的关系式，}
\step{得 $3x^2-px+5p-65=75-5p+5p-65=10$，与 $y=10$ 相符，}
\step{所以 $(5,10)\in A_p$ 对每一个 $p$ 都成立，因此所有 $A_p$ 的交集为 $\{(5,10)\}$.}
\solend
\fi

\end{enumerate}

% 守卫须放在 \bigsec 之前，否则标题仍会落页末（见 高一40 的教训）
\needvspace{5cm}

\bigsec{二、选择题}

\begin{enumerate}
\setcounter{enumi}{12}

\item 已知 $a,b,c\in\R$ 且 $a>b$，则下列不等式一定成立的是\mbox{（\quad）}

\keepwithprev
A. $\dfrac1a<\dfrac1b$\qquad B. $a^2>b^2$

C. $a|c|>b|c|$\qquad D. $\dfrac a{c^2+1}>\dfrac b{c^2+1}$

\ifanswer
\solbegin
\step{$A$：当 $a>0>b$ 时，$\dfrac1a>\dfrac1b$，故 $A$ 错误；}
\step{$B$：当 $a=-1,b=-2$ 时，满足 $a>b$，但 $a^2>b^2$ 不成立，故 $B$ 错误；}
\step{$C$：当 $c=0$ 时，$a|c|=b|c|$，故 $C$ 错误；}
\step{$D$：由 $a>b,c^2+1>0$，得 $\dfrac a{c^2+1}>\dfrac b{c^2+1}$，故 $D$ 正确；}
\step{故选 $D$.}
\solend
\fi

\item 已知集合 $A=\{a^2,0,-1\}$，$B=\{a,b,0\}$，若 $A=B$，则 $(ab)^{2026}$ 的值为\mbox{（\quad）}

\keepwithprev
A. $0$\qquad B. $-1$

C. $1$\qquad D. $\pm1$

\ifanswer
\solbegin
\step{因为集合 $A=\{a^2,0,-1\}$，$B=\{a,b,0\}$，$A=B$，所以 $\begin{cases}a=-1\\ a^2=b\end{cases}$，}
\step{所以 $a=-1$，$b=1$，所以 $(ab)^{2026}=1$，故选 $C$.}
\solend
\fi

\item 设全集 $U=\R$，$A=\{x\mid x<-1\text{ 或 }x>2\}$,$B=\{y\mid y=x^2+a\}$，若 $\overline A\subseteq B$，则实数 $a$ 的取值范围是\mbox{（\quad）}

\keepwithprev
A. $a\leqslant-1$\qquad B. $a>2$

C. $a<-1$\qquad D. $a\geqslant-2$

\ifanswer
\solbegin
\step{因为 $A=\{x\mid x<-1\text{ 或 }x>2\}$，则 $\overline A=\{x\mid-1\leqslant x\leqslant2\}$，}
\step{又 $B=\{y\mid y=x^2+a\}=\{y\mid y\geqslant a\}$，且 $\overline A\subseteq B$，所以 $a\leqslant-1$，故选 $A$.}
\solend
\fi

\item 设 $P$ 和 $Q$ 是两个集合，定义集合 $P-Q=\{x\mid x\in P$，且 $x\notin Q\}$，如果 $Q=\{x\mid(1-x)(x-3)>0\}$，$P=\left\{x\left|\dfrac x{2-x}>0\right.\right\}$，那么 $P-Q$ 等于\mbox{（\quad）}

\keepwithprev
A. $\{x\mid0<x<1\}$\qquad B. $\{x\mid1\leqslant x<2\}$

C. $\{x\mid0<x\leqslant1\}$\qquad D. $\{x\mid2\leqslant x<3\}$

\ifanswer
\solbegin
\step{因为 $P=\left\{x\left|\dfrac x{2-x}>0\right.\right\}=\{x\mid x(x-2)<0\}=\{x\mid0<x<2\}$，}
\step{$Q=\{x\mid(1-x)(x-3)>0\}=\{x\mid(x-3)(x-1)<0\}=\{x\mid1<x<3\}$，}
\step{所以 $P-Q=\{x\mid0<x\leqslant1\}$，故选 $C$.}
\solend
\fi

\end{enumerate}

\bigsec{三、解答题}

\begin{enumerate}
\setcounter{enumi}{16}

\item 已知集合 $A=\{x\mid x^2-9\geqslant0\}$，$B=\{x\mid|x-4|<2\}$，$C=\left\{x\left|\dfrac{x-8}{x+2}<0\right.\right\}$.

（1）求 $A\cup C$；

（2）若全集 $U=R$，求 $\overline A\cap B$.

\ifanswer
\solbegin
\stepm{(1)}{集合 $A=\{x\mid x^2-9\geqslant0\}=\{x\mid x\geqslant3\text{ 或 }x\leqslant-3\}$，}
\step{$B=\{x\mid|x-4|<2\}=\{x\mid2<x<6\}$，$C=\left\{x\left|\dfrac{x-8}{x+2}<0\right.\right\}=\{x\mid-2<x<8\}$，}
\step{所以 $A\cup C=\{x\mid x\leqslant-3\text{ 或 }x>-2\}$；}
\stepm{(2)}{全集 $U=R$，集合 $A=\{x\mid x^2-9\geqslant0\}=\{x\mid x\geqslant3\text{ 或 }x\leqslant-3\}$，}
\step{$B=\{x\mid|x-4|<2\}=\{x\mid2<x<6\}$，}
\step{所以 $\overline A=\{x\mid-3<x<3\}$，所以 $\overline A\cap B=\{x\mid2<x<3\}$.}
\solend
\fi

\ansspace[6]

\item 已知集合 $A=\left\{x\left|\dfrac{2x-1}{x+1}<1,x\in\R\right.\right\}$，集合 $B=\{x\mid(x-a-1)(x-a+1)\leqslant0\}$.

（1）求集合 $A$；

（2）若 $B\subseteq A$，求实数 $a$ 的取值范围.

\ifanswer
\solbegin
\stepm{(1)}{由 $\dfrac{2x-1}{x+1}<1$，得 $\dfrac{x-2}{x+1}<0$ 所以 $(x-2)(x+1)<0$，解得 $-1<x<2$，}
\step{即 $A=(-1,2)$；}
\stepm{(2)}{由 $(x-a-1)(x-a+1)\leqslant0$，得 $a-1\leqslant x\leqslant a+1$，}
\step{所以 $B=\{x\mid a-1\leqslant x\leqslant a+1\}$，}
\step{又 $B\subseteq A$，且 $a+1>a-1$ 即 $B\neq\varnothing$，所以 $\begin{cases}a-1>-1\\ a+1<2\end{cases}$，所以 $0<a<1$，}
\step{所以实数 $a$ 的取值范围是 $\{a\mid0<a<1\}$.}
\solend
\fi

\ansspace[6]

\item 已知集合 $A=\{x\mid x^2-4x+3=0\}$，$B=\{x\mid x^2-ax+3=0\}$.

（1）若 $A\cup B=B$，求实数 $a$ 的值；

（2）若 $A\cap B=B$，求实数 $a$ 的取值范围.

\ifanswer
\solbegin
\stepm{(1)}{集合 $A=\{x\mid x^2-4x+3=0\}=\{1,3\}$，$B=\{x\mid x^2-ax+3=0\}$.}
\step{$A\cup B=B$，所以 $1\in B$，$3\in B$，所以 $\begin{cases}1-a+3=0\\ 9-3a+3=0\end{cases}$，解得实数 $a=4$；}
\stepm{(2)}{因为 $A\cap B=B$，所以 $B\subseteq A$，}
\step{当 $B=\varnothing$ 时，$\Delta=(-a)^2-12<0$，解得 $-2\sqrt3<a<2\sqrt3$，}
\step{当 $B$ 是单元素集时，$\Delta=a^2-12=0$，解得 $a=2\sqrt3$，此时 $B=\{\sqrt3\}$，}
\step{不合题意；}
\step{当 $B=\{1,3\}$ 时，由 (1) 得 $a=4$，}
\step{综上，实数 $a$ 的取值范围是 $(-2\sqrt3,2\sqrt3)\cup\{4\}$.}
\solend
\fi

\ansspace[6]

\item 已知全集为 $R$，实数 $a$、$b$ 满足 $a\neq b$，集合 $M=\{x\mid x^2-3x-4\leqslant0\}$，集合 $A=\{x\mid(x+a)(x+b)>0\}$，集合 $B=\{x\mid x^2+(a-2)x-2a>0\}$.

（1）若 $\overline A=M$，求 $a$、$b$ 的值；

（2）若 $a>b>-1$，求 $A\cap B$；

（3）若 $a^2+a\in\overline B$，求 $a$ 的取值范围.

\ifanswer
\solbegin
\stepm{(1)}{$M=\{x\mid x^2-3x-4\leqslant0\}=\{x\mid-1\leqslant x\leqslant4\}$，$A=\{x\mid(x+a)(x+b)\leqslant0\}$，}
\step{因为 $\overline A=M$，所以 $a=1$，$b=-4$ 或 $a=-4$，$b=1$.}
\stepm{(2)}{因为 $a>b>-1$，所以 $-a<-b<1$，所以 $A=\{x\mid x<-a\text{ 或 }x>-b\}$，}
\step{因为 $B=\{x\mid x^2+(a-2)x-2a>0\}=\{x\mid x<-a\text{ 或 }x>2\}$，}
\step{所以 $A\cap B=\{x\mid x<-a\text{ 或 }x>2\}$，}
\stepm{(3)}{$\overline B=\{x\mid(x-2)(x+a)\leqslant0\}$，}
\step{由 $a^2+a\in\overline B$，得 $(a^2+a-2)(a^2+2a)\leqslant0\Leftrightarrow a(a-1)(a-2)^2\leqslant0$，}
\step{解得 $a=2$ 或 $0\leqslant a\leqslant1$，故 $a$ 的取值范围是 $[0,1]\cup\{2\}$.}
\solend
\fi

\ansspace[8]

\item 对平面直角坐标系的第一象限内的任意两点作如下定义：若 $\dfrac ab>\dfrac cd$，那么称点 $(a,b)$ 是点 $(c,d)$ 的“上位点”，同时点 $(c,d)$ 是点 $(a,b)$ 的“下位点”.

（1）试写出点 $(2,4)$ 的一个“上位点”坐标；

（2）已知点 $(a,b)$ 是点 $(c,d)$ 的“上位点”，判断点 $P(a+c,b+d)$ 是点 $(c,d)$ 的“上位点”还是“下位点”，并证明你的结论；

（3）设正整数 $n$ 满足以下条件：对集合 $\{t\mid0<t<2020,t\in\Z\}$ 内的任意元素 $m$，总存在正整数 $k=2m+1$，使得点 $(n,k)$ 既是点 $(2025,m)$ 的“下位点”，又是点 $(2026,m+1)$ 的“上位点”，问这样的正整数 $n$ 是否存在？若存在，求正整数 $n$ 的值；若不存在，请说明理由.

\ifanswer
\solbegin
\stepm{(1)}{由定义，点 $(a,b)$ 是点 $(2,4)$ 的“上位点”，}
\step{当且仅当 $\dfrac ab>\dfrac24=\dfrac12$ 且 $a>0,b>0$.}
\step{取 $(a,b)=(1,1)$，因为 $\dfrac11>\dfrac12$，}
\step{所以点 $(1,1)$ 是点 $(2,4)$ 的一个“上位点”.（答案不唯一）}
\stepm{(2)}{点 $P(a+c,b+d)$ 是点 $(c,d)$ 的“上位点”.}
\step{$(a,b)$ 是点 $(c,d)$ 的“上位点”，所以 $\dfrac ab>\dfrac cd$.}
\step{因为 $a,b,c,d>0$，两边同乘 $bd$ 得 $ad>bc$，即 $ad-bc>0$.}
\step{作差比较 $\dfrac{a+c}{b+d}$ 与 $\dfrac cd$：$\dfrac{a+c}{b+d}-\dfrac cd=\dfrac{d(a+c)-c(b+d)}{d(b+d)}=\dfrac{ad-bc}{d(b+d)}>0$，}
\step{所以点 $P(a+c,b+d)$ 是点 $(c,d)$ 的“上位点”.}
\stepm{(3)}{由定义，点 $(n,k)$ 是点 $(2025,m)$ 的下位点即 $\dfrac{2025}m>\dfrac nk$；}
\step{点 $(n,k)$ 是点 $(2026,m+1)$ 的上位点即 $\dfrac nk>\dfrac{2026}{m+1}$.}
\step{代入 $k=2m+1$，得对任意 $m\in\{1,2,\cdots,2019\}$，}
\step{都有 $4052-\dfrac{2026}{m+1}<n<4050+\dfrac{2025}m$.}
\step{右边 $4050+\dfrac{2025}m$ 随 $m$ 增大而减小，故最近处在 $m=2019$，}
\step{得 $n<4050+\dfrac{2025}{2019}<4051.01$，所以 $n\leqslant4051$.}
\step{左边 $4052-\dfrac{2026}{m+1}$ 随 $m$ 增大而增大，故最近处也在 $m=2019$，}
\step{得 $n>4052-\dfrac{2026}{2020}>4050.99$，所以 $n\geqslant4051$.}
\step{综上，$n=4051$.}
\step{检验：当 $n=4051$ 时，对任意 $1\leqslant m\leqslant2019,k=2m+1$ 是正整数，}
\step{且 $\dfrac{4051}{2m+1}<\dfrac{2025}m$,$4051m<4050m+2025,m<2025$，成立；}
\step{$\dfrac{4051}{2m+1}>\dfrac{2026}{m+1}$,$4051(m+1)>4052m+2026,m<2025$，成立.}
\step{又各点横纵坐标均为正，都在第一象限内，符合定义，}
\step{所以满足条件的正整数 $n$ 存在，其值为 $n=4051$.}
\solend
\fi

\ansspace*[10]

\end{enumerate}

\end{document}
"
%
% 原始材料：微信公众号“魔都数学加油站”文章，作者破晓，连载“27学年高一59”；
% 原文标题“27学年高一59——2026-2027年上海市华实高中高一上10月月考”；
% 原文链接：https://mp.weixin.qq.com/s/XzyBCmrvpRuMe-NuNk7gkA。
% 页面标注发布时间 2026-10-13 21:45（上海）。
% 正文为 9 张 PNG 页：第 01、02、03、05、06、07、08、09 页 4762×6735，第 04 页 2560×3620
% （**同一篇各页分辨率不同**），均按页序读取。9 页全为试卷正文页（卷首标题在第 1 页页顶，
% 卷末解答题第 21 题解析止于第 9 页），无封面页、目录页、二维码或宣传页。
% 粉色斜排“魔都数学加油站”为卷面水印，与公众号页脚一并未录入。
%
% 卷首标题：照卷面印作“2026--2027 年华实高中高一上 10 月月考”（公众号标题中的“上海市”
% 卷面标题行没有，本稿照卷面），年份区间按全稿惯例排 en dash。原卷只有一行标题、无副标题。
%
% 篇幅与题号：原卷自第 3 题起编号（第 1、2 题不在本卷内）。填空题 3—12（10 题）、
% 选择题 13—16（4 题）、解答题 17—21（5 题），共 19 题。本稿沿用原节与原题号，
% 只改排为全稿统一的大节样式（原卷小节标题为“一、填空题”“二、选择题”“三、解答题”）。
%
% 讲评版：原卷每题下直接印【答案】或【解析】。**只印【答案】的 1 题**（第 3 题），本稿只给 \ans{}；
% 其余 18 题都印【解析】，按全稿口径这类题不另提 \ans{}——第 13、14、15、16 题的结论写在解析末句
% （“故选 D”等），亦照录在解析内。本卷**没有**印在题干括号内的答案。
%
% 副标题：原卷未印副标题。本稿按“知识模块”口径标作“集合与常用逻辑用语、不等式”：
% 19 题中集合与常用逻辑用语 13 题（第 3、9、10、11、12、14、15、16、17、18、19、20、21 题），
% 不等式 10 题（含等式与不等式性质、一元二次不等式、恒成立、基本应用：第 4、5、6、7、8、
% 13、17、18、19、20 题），两类均远超三分之一，故并标。口径同上。
%
% 与原卷的排版差异：
%   ① 原卷每题下直接印【答案】/【解析】，本稿按全稿惯例将答案与解析置于答案版，试卷版保持干净；
%   ② 原卷未印副标题，本稿按“知识模块”口径补出；
%   ③ 原卷填空横线约两个字宽，本稿按全稿惯例统一排 \blankline[4.4em]；
%   ④ 原卷实数集排粗体 $\mathbf{R}$（第 12、13、20 题等），整数集排粗体 $\mathbf{Z}$（第 11、12 题），
%      自然数集排斜体 $N$ 或粗体 $\mathbf{Z}$ 处按本稿惯例改用 $\R$、$\Z$、$\N$；
%      空集记号统一为 \varnothing；
%   ⑤ 原卷在数学式内多处用半角逗号（如 $x_1,x_2$、$a,b\in R$），均照录未强改；
%      句末句点按全稿惯例排西文句点；“漂亮集”“上位点”“下位点”等自定义名词照录；
%   ⑥ 第 17—21 题的子问原卷各自另起一行，本稿按全稿惯例用 enumerate 的 (1)(2)(3)；
%      解析中的子问标号一律用 \stepm{(1)}{…}；第 9、11 题解析中的①、②（③④）分情形或分结论，
%      按全稿惯例作 \stepm{①}{…}、\stepm{②}{…}（第 11 题解析末“故正确的是②④”照录为行内）；
%   ⑦ 第 12 题题干里的“$A_p$”下标原卷作斜体 $p$，照录；
%   ⑧ 每道解答题原卷未留作答空白，本稿按全稿惯例加 \ansspace。
%
% 原卷笔误/存疑（均照录未改，正文对应处另有 % 注）：本卷首轮逐页转录，暂未发现影响答案的笔误。
%
% 插图：无（全卷 9 页无任何图形、无“如图”字样）。
%
\documentclass[a4paper,11pt]{article}
\usepackage{common}

\begin{document}

\examtitlest[3]{2026--2027 年华实高中高一上 10 月月考}{集合与常用逻辑用语、不等式}

\bigsec{一、填空题}

\begin{enumerate}
\setcounter{enumi}{2}

\item “$x\geqslant1$ 且 $y\geqslant1$”的否定形式为\blankline[4.4em].

\ifanswer
\ans{$x<1$ 或 $y<1$}
\fi

\item 已知方程 $2x^2+4x-3=0$ 的两个根为 $x_1,x_2$，则 $\dfrac1{x_1}+\dfrac1{x_2}=$\blankline[4.4em].

\ifanswer
\solbegin
\step{利用韦达定理得 $\begin{cases}x_1+x_2=-2\\ x_1\cdot x_2=-\dfrac32\end{cases}$，则 $\dfrac1{x_1}+\dfrac1{x_2}=\dfrac{x_1+x_2}{x_1\cdot x_2}=\dfrac43$.}
\solend
\fi

\item 设 $a>b>0$，比较大小：$\dfrac1{a-b}$\blankline[2.4em]$\dfrac1a$.

\ifanswer
\solbegin
\step{由 $a>b>0$ 得 $a>0$ 且 $a-b>0$，所以 $\dfrac1{a-b}-\dfrac1a=\dfrac b{a(a-b)}>0$，}
\step{故 $\dfrac1{a-b}>\dfrac1a$.}
\solend
\fi

\item 若不等式 $2ax^2-ax+1>0$ 的解集为 $R$，则实数 $a$ 的取值范围是\blankline[4.4em].

\ifanswer
\solbegin
\step{当 $a\neq0$ 时，由题意得 $\begin{cases}a>0\\ \Delta<0\end{cases}$，即 $\begin{cases}a>0\\ a^2-8a<0\end{cases}$，解得 $0<a<8$.}
\step{当 $a=0$ 时，恒有 $1>0$，不等式也成立.}
\step{综上所述，$a$ 的取值范围是 $[0,8)$.}
\solend
\fi

\item 集合 $A=\{x\mid m-1\leqslant x\leqslant m+7\}$，集合 $B=\{x\mid x^2-3x-10\geqslant0\}$，若 $A\cup B=R$，则实数 $m$ 的取值范围为\blankline[4.4em].

\ifanswer
\solbegin
\step{$A=\{x\mid m-1\leqslant x\leqslant m+7\}$，$B=\{x\mid x^2-3x-10\geqslant0\}=\{x\mid x\leqslant-2\text{ 或 }x\geqslant5\}$，}
\step{若 $A\cup B=R$，则 $\begin{cases}m-1\leqslant-2\\ m+7\geqslant5\end{cases}$，解得 $-2\leqslant m\leqslant-1$，}
\step{所以实数 $m$ 的取值范围为 $[-2,-1]$.}
\solend
\fi

\item 若关于 $x$ 的不等式 $ax^2+bx+c>0$ 的解集为 $(-1,3)$，则关于 $x$ 的不等式 $cx^2+bx+a>0$ 的解集为\blankline[4.4em].

\ifanswer
\solbegin
\step{因为不等式 $ax^2+bx+c>0$ 的解集为 $(-1,3)$，}
\step{所以 $-1$ 和 $3$ 是方程 $ax^2+bx+c=0$ 的实数解，}
\step{所以 $\begin{cases}-1+3=-\dfrac ba\\ -1\times3=\dfrac ca\\ a<0\end{cases}$，解得 $b=-2a$，$c=-3a$，且 $a<0$，}
\step{所以不等式 $cx^2+bx+a>0$ 可化为 $-3ax^2-2ax+a>0$，即 $3x^2+2x-1>0$，}
\step{解得 $x<-1$ 或 $x>\dfrac13$，所以该不等式的解集为 $(-\infty,-1)\cup\left(\dfrac13,+\infty\right)$.}
\solend
\fi

\item 若集合 $A=\{x\mid(a-1)x^2-2x+1=0\}$ 有且仅有 $2$ 个子集，则实数 $a$ 的值为\blankline[4.4em].

\ifanswer
\solbegin
\step{由题意得方程 $(a-1)x^2-2x+1=0$ 只有一个根，}
\stepm{①}{当 $a-1=0$ 时，$a=1$，此时方程化为 $-2x+1=0$，解得 $x=\dfrac12$，符合题意，}
\stepm{②}{当 $a-1\neq0$，即 $a\neq1$ 时，$\Delta=4-4(a-1)=0$，解得 $a=2$，}
\step{综上所述，实数 $a$ 的值为 $1$ 或 $2$.}
\solend
\fi

\item “$x+y\neq5$”是“$x\neq3$ 或 $y\neq2$”的\blankline[4.4em]条件.

\ifanswer
\solbegin
\step{分析充分性：}
\step{若 $x+y\neq5$ 成立，若 $x=3$，则 $y\neq2$，若 $y=2$，则 $x\neq3$，}
\step{所以一定能说明 $x\neq3$ 和 $y\neq2$ 至少有一个成立，即 $x\neq3$ 或 $y\neq2$ 成立，}
\step{充分性成立；}
\step{分析必要性：}
\step{若 $x\neq3$ 或 $y\neq2$，即 $x\neq3$ 和 $y\neq2$ 二者至少一个成立，}
\step{举例：$x=1,y=4$，此时 $x+y=5$，所以必要性不成立.}
\step{所以“$x+y\neq5$”是“$x\neq3$ 或 $y\neq2$”的充分不必要条件.}
\solend
\fi

\item 设 $P$ 为非空实数集满足：对任意给定的 $x,y\in P$（$x,y$ 可以相同），都有 $x+y\in P,x-y\in P,xy\in P$，则称 $P$ 为漂亮集. 关于漂亮集有下列结论：

①集合 $P=\{-2,-1,0,1,2\}$ 为漂亮集；②集合 $P=\{x\mid x=2n,n\in\Z\}$ 为漂亮集；

③若集合 $P_1,P_2$ 为漂亮集，则 $P_1\cup P_2$ 为漂亮集；④若集合 $P$ 为漂亮集，则一定有 $0\in P$.

其中所有正确结论的序号是\blankline[4.4em].

\ifanswer
\solbegin
\step{对于集合 $P=\{-2,-1,0,1,2\}$，取 $x=2,y=-2$，则 $x-y=2-(-2)=4\notin P$，}
\step{不满足漂亮集的定义，所以①错误.}
\step{对于集合 $P=\{x\mid x=2n,n\in Z\}$，设 $x=2n_1,y=2n_2,n_1,n_2\in Z$.}
\step{$x+y=2n_1+2n_2=2(n_1+n_2)$，因为 $n_1+n_2\in Z$，所以 $x+y\in P$.}
\step{$x-y=2n_1-2n_2=2(n_1-n_2)$，因为 $n_1-n_2\in Z$，所以 $x-y\in P$.}
\step{$xy=2n_1\times2n_2=4n_1n_2=2(2n_1n_2)$，因为 $2n_1n_2\in Z$，所以 $xy\in P$.}
\step{满足漂亮集的定义，所以②正确.}
\step{令 $P_1=\{x\mid x=2n,n\in Z\}$,$P_2=\{x\mid x=3n,n\in Z\}$，都是漂亮集.}
\step{$P_1\cup P_2$ 中，取 $x=2(x\in P_1)$,$y=3(y\in P_2)$,$x+y=5\notin P_1\cup P_2$，}
\step{不满足漂亮集的定义，所以③错误.}
\step{因为对任意 $x\in P$，有 $x-x=0\in P$，所以若集合 $P$ 为漂亮集，}
\step{则一定有 $0\in P$，所以④正确.}
\step{故正确的是②④.}
\solend
\fi

\item 已知：集合 $A_p=\{(x,y)\mid y=3x^2-px+5p-65,x\in\R\}$,$M=\{x\mid1\leqslant x\leqslant100,x\in\Z\}$，若 $p\in M$，则所有 $A_p$ 的交集\blankline[4.4em].

\ifanswer
\solbegin
\step{由 $M=\{1,2,\cdots,100\}$，所以所有 $A_p$ 的交集就是这些集合中共同拥有的点.}
\step{若 $(x,y)$ 是所有集合的一个公共点，}
\step{那么它特别要同时属于 $p=1$ 和 $p=2$ 时对应的两个集合.}
\step{当 $p=1$ 时，$y=3x^2-x-60$；}
\step{当 $p=2$ 时，$y=3x^2-2x-55$；联立得 $x=5,y=10$.}
\step{因此，如果公共点存在，只可能是 $(5,10)$.}
\step{接下来验证这个点是否属于所有集合，}
\step{对任意 $p\in M$，把 $x=5$ 代入集合的关系式，}
\step{得 $3x^2-px+5p-65=75-5p+5p-65=10$，与 $y=10$ 相符，}
\step{所以 $(5,10)\in A_p$ 对每一个 $p$ 都成立，因此所有 $A_p$ 的交集为 $\{(5,10)\}$.}
\solend
\fi

\end{enumerate}

% 守卫须放在 \bigsec 之前，否则标题仍会落页末（见 高一40 的教训）
\needvspace{5cm}

\bigsec{二、选择题}

\begin{enumerate}
\setcounter{enumi}{12}

\item 已知 $a,b,c\in\R$ 且 $a>b$，则下列不等式一定成立的是\mbox{（\quad）}

\keepwithprev
A. $\dfrac1a<\dfrac1b$\qquad B. $a^2>b^2$

C. $a|c|>b|c|$\qquad D. $\dfrac a{c^2+1}>\dfrac b{c^2+1}$

\ifanswer
\solbegin
\step{$A$：当 $a>0>b$ 时，$\dfrac1a>\dfrac1b$，故 $A$ 错误；}
\step{$B$：当 $a=-1,b=-2$ 时，满足 $a>b$，但 $a^2>b^2$ 不成立，故 $B$ 错误；}
\step{$C$：当 $c=0$ 时，$a|c|=b|c|$，故 $C$ 错误；}
\step{$D$：由 $a>b,c^2+1>0$，得 $\dfrac a{c^2+1}>\dfrac b{c^2+1}$，故 $D$ 正确；}
\step{故选 $D$.}
\solend
\fi

\item 已知集合 $A=\{a^2,0,-1\}$，$B=\{a,b,0\}$，若 $A=B$，则 $(ab)^{2026}$ 的值为\mbox{（\quad）}

\keepwithprev
A. $0$\qquad B. $-1$

C. $1$\qquad D. $\pm1$

\ifanswer
\solbegin
\step{因为集合 $A=\{a^2,0,-1\}$，$B=\{a,b,0\}$，$A=B$，所以 $\begin{cases}a=-1\\ a^2=b\end{cases}$，}
\step{所以 $a=-1$，$b=1$，所以 $(ab)^{2026}=1$，故选 $C$.}
\solend
\fi

\item 设全集 $U=\R$，$A=\{x\mid x<-1\text{ 或 }x>2\}$,$B=\{y\mid y=x^2+a\}$，若 $\overline A\subseteq B$，则实数 $a$ 的取值范围是\mbox{（\quad）}

\keepwithprev
A. $a\leqslant-1$\qquad B. $a>2$

C. $a<-1$\qquad D. $a\geqslant-2$

\ifanswer
\solbegin
\step{因为 $A=\{x\mid x<-1\text{ 或 }x>2\}$，则 $\overline A=\{x\mid-1\leqslant x\leqslant2\}$，}
\step{又 $B=\{y\mid y=x^2+a\}=\{y\mid y\geqslant a\}$，且 $\overline A\subseteq B$，所以 $a\leqslant-1$，故选 $A$.}
\solend
\fi

\item 设 $P$ 和 $Q$ 是两个集合，定义集合 $P-Q=\{x\mid x\in P$，且 $x\notin Q\}$，如果 $Q=\{x\mid(1-x)(x-3)>0\}$，$P=\left\{x\left|\dfrac x{2-x}>0\right.\right\}$，那么 $P-Q$ 等于\mbox{（\quad）}

\keepwithprev
A. $\{x\mid0<x<1\}$\qquad B. $\{x\mid1\leqslant x<2\}$

C. $\{x\mid0<x\leqslant1\}$\qquad D. $\{x\mid2\leqslant x<3\}$

\ifanswer
\solbegin
\step{因为 $P=\left\{x\left|\dfrac x{2-x}>0\right.\right\}=\{x\mid x(x-2)<0\}=\{x\mid0<x<2\}$，}
\step{$Q=\{x\mid(1-x)(x-3)>0\}=\{x\mid(x-3)(x-1)<0\}=\{x\mid1<x<3\}$，}
\step{所以 $P-Q=\{x\mid0<x\leqslant1\}$，故选 $C$.}
\solend
\fi

\end{enumerate}

\bigsec{三、解答题}

\begin{enumerate}
\setcounter{enumi}{16}

\item 已知集合 $A=\{x\mid x^2-9\geqslant0\}$，$B=\{x\mid|x-4|<2\}$，$C=\left\{x\left|\dfrac{x-8}{x+2}<0\right.\right\}$.

（1）求 $A\cup C$；

（2）若全集 $U=R$，求 $\overline A\cap B$.

\ifanswer
\solbegin
\stepm{(1)}{集合 $A=\{x\mid x^2-9\geqslant0\}=\{x\mid x\geqslant3\text{ 或 }x\leqslant-3\}$，}
\step{$B=\{x\mid|x-4|<2\}=\{x\mid2<x<6\}$，$C=\left\{x\left|\dfrac{x-8}{x+2}<0\right.\right\}=\{x\mid-2<x<8\}$，}
\step{所以 $A\cup C=\{x\mid x\leqslant-3\text{ 或 }x>-2\}$；}
\stepm{(2)}{全集 $U=R$，集合 $A=\{x\mid x^2-9\geqslant0\}=\{x\mid x\geqslant3\text{ 或 }x\leqslant-3\}$，}
\step{$B=\{x\mid|x-4|<2\}=\{x\mid2<x<6\}$，}
\step{所以 $\overline A=\{x\mid-3<x<3\}$，所以 $\overline A\cap B=\{x\mid2<x<3\}$.}
\solend
\fi

\ansspace[6]

\item 已知集合 $A=\left\{x\left|\dfrac{2x-1}{x+1}<1,x\in\R\right.\right\}$，集合 $B=\{x\mid(x-a-1)(x-a+1)\leqslant0\}$.

（1）求集合 $A$；

（2）若 $B\subseteq A$，求实数 $a$ 的取值范围.

\ifanswer
\solbegin
\stepm{(1)}{由 $\dfrac{2x-1}{x+1}<1$，得 $\dfrac{x-2}{x+1}<0$ 所以 $(x-2)(x+1)<0$，解得 $-1<x<2$，}
\step{即 $A=(-1,2)$；}
\stepm{(2)}{由 $(x-a-1)(x-a+1)\leqslant0$，得 $a-1\leqslant x\leqslant a+1$，}
\step{所以 $B=\{x\mid a-1\leqslant x\leqslant a+1\}$，}
\step{又 $B\subseteq A$，且 $a+1>a-1$ 即 $B\neq\varnothing$，所以 $\begin{cases}a-1>-1\\ a+1<2\end{cases}$，所以 $0<a<1$，}
\step{所以实数 $a$ 的取值范围是 $\{a\mid0<a<1\}$.}
\solend
\fi

\ansspace[6]

\item 已知集合 $A=\{x\mid x^2-4x+3=0\}$，$B=\{x\mid x^2-ax+3=0\}$.

（1）若 $A\cup B=B$，求实数 $a$ 的值；

（2）若 $A\cap B=B$，求实数 $a$ 的取值范围.

\ifanswer
\solbegin
\stepm{(1)}{集合 $A=\{x\mid x^2-4x+3=0\}=\{1,3\}$，$B=\{x\mid x^2-ax+3=0\}$.}
\step{$A\cup B=B$，所以 $1\in B$，$3\in B$，所以 $\begin{cases}1-a+3=0\\ 9-3a+3=0\end{cases}$，解得实数 $a=4$；}
\stepm{(2)}{因为 $A\cap B=B$，所以 $B\subseteq A$，}
\step{当 $B=\varnothing$ 时，$\Delta=(-a)^2-12<0$，解得 $-2\sqrt3<a<2\sqrt3$，}
\step{当 $B$ 是单元素集时，$\Delta=a^2-12=0$，解得 $a=2\sqrt3$，此时 $B=\{\sqrt3\}$，}
\step{不合题意；}
\step{当 $B=\{1,3\}$ 时，由 (1) 得 $a=4$，}
\step{综上，实数 $a$ 的取值范围是 $(-2\sqrt3,2\sqrt3)\cup\{4\}$.}
\solend
\fi

\ansspace[6]

\item 已知全集为 $R$，实数 $a$、$b$ 满足 $a\neq b$，集合 $M=\{x\mid x^2-3x-4\leqslant0\}$，集合 $A=\{x\mid(x+a)(x+b)>0\}$，集合 $B=\{x\mid x^2+(a-2)x-2a>0\}$.

（1）若 $\overline A=M$，求 $a$、$b$ 的值；

（2）若 $a>b>-1$，求 $A\cap B$；

（3）若 $a^2+a\in\overline B$，求 $a$ 的取值范围.

\ifanswer
\solbegin
\stepm{(1)}{$M=\{x\mid x^2-3x-4\leqslant0\}=\{x\mid-1\leqslant x\leqslant4\}$，$A=\{x\mid(x+a)(x+b)\leqslant0\}$，}
\step{因为 $\overline A=M$，所以 $a=1$，$b=-4$ 或 $a=-4$，$b=1$.}
\stepm{(2)}{因为 $a>b>-1$，所以 $-a<-b<1$，所以 $A=\{x\mid x<-a\text{ 或 }x>-b\}$，}
\step{因为 $B=\{x\mid x^2+(a-2)x-2a>0\}=\{x\mid x<-a\text{ 或 }x>2\}$，}
\step{所以 $A\cap B=\{x\mid x<-a\text{ 或 }x>2\}$，}
\stepm{(3)}{$\overline B=\{x\mid(x-2)(x+a)\leqslant0\}$，}
\step{由 $a^2+a\in\overline B$，得 $(a^2+a-2)(a^2+2a)\leqslant0\Leftrightarrow a(a-1)(a-2)^2\leqslant0$，}
\step{解得 $a=2$ 或 $0\leqslant a\leqslant1$，故 $a$ 的取值范围是 $[0,1]\cup\{2\}$.}
\solend
\fi

\ansspace[8]

\item 对平面直角坐标系的第一象限内的任意两点作如下定义：若 $\dfrac ab>\dfrac cd$，那么称点 $(a,b)$ 是点 $(c,d)$ 的“上位点”，同时点 $(c,d)$ 是点 $(a,b)$ 的“下位点”.

（1）试写出点 $(2,4)$ 的一个“上位点”坐标；

（2）已知点 $(a,b)$ 是点 $(c,d)$ 的“上位点”，判断点 $P(a+c,b+d)$ 是点 $(c,d)$ 的“上位点”还是“下位点”，并证明你的结论；

（3）设正整数 $n$ 满足以下条件：对集合 $\{t\mid0<t<2020,t\in\Z\}$ 内的任意元素 $m$，总存在正整数 $k=2m+1$，使得点 $(n,k)$ 既是点 $(2025,m)$ 的“下位点”，又是点 $(2026,m+1)$ 的“上位点”，问这样的正整数 $n$ 是否存在？若存在，求正整数 $n$ 的值；若不存在，请说明理由.

\ifanswer
\solbegin
\stepm{(1)}{由定义，点 $(a,b)$ 是点 $(2,4)$ 的“上位点”，}
\step{当且仅当 $\dfrac ab>\dfrac24=\dfrac12$ 且 $a>0,b>0$.}
\step{取 $(a,b)=(1,1)$，因为 $\dfrac11>\dfrac12$，}
\step{所以点 $(1,1)$ 是点 $(2,4)$ 的一个“上位点”.（答案不唯一）}
\stepm{(2)}{点 $P(a+c,b+d)$ 是点 $(c,d)$ 的“上位点”.}
\step{$(a,b)$ 是点 $(c,d)$ 的“上位点”，所以 $\dfrac ab>\dfrac cd$.}
\step{因为 $a,b,c,d>0$，两边同乘 $bd$ 得 $ad>bc$，即 $ad-bc>0$.}
\step{作差比较 $\dfrac{a+c}{b+d}$ 与 $\dfrac cd$：$\dfrac{a+c}{b+d}-\dfrac cd=\dfrac{d(a+c)-c(b+d)}{d(b+d)}=\dfrac{ad-bc}{d(b+d)}>0$，}
\step{所以点 $P(a+c,b+d)$ 是点 $(c,d)$ 的“上位点”.}
\stepm{(3)}{由定义，点 $(n,k)$ 是点 $(2025,m)$ 的下位点即 $\dfrac{2025}m>\dfrac nk$；}
\step{点 $(n,k)$ 是点 $(2026,m+1)$ 的上位点即 $\dfrac nk>\dfrac{2026}{m+1}$.}
\step{代入 $k=2m+1$，得对任意 $m\in\{1,2,\cdots,2019\}$，}
\step{都有 $4052-\dfrac{2026}{m+1}<n<4050+\dfrac{2025}m$.}
\step{右边 $4050+\dfrac{2025}m$ 随 $m$ 增大而减小，故最近处在 $m=2019$，}
\step{得 $n<4050+\dfrac{2025}{2019}<4051.01$，所以 $n\leqslant4051$.}
\step{左边 $4052-\dfrac{2026}{m+1}$ 随 $m$ 增大而增大，故最近处也在 $m=2019$，}
\step{得 $n>4052-\dfrac{2026}{2020}>4050.99$，所以 $n\geqslant4051$.}
\step{综上，$n=4051$.}
\step{检验：当 $n=4051$ 时，对任意 $1\leqslant m\leqslant2019,k=2m+1$ 是正整数，}
\step{且 $\dfrac{4051}{2m+1}<\dfrac{2025}m$,$4051m<4050m+2025,m<2025$，成立；}
\step{$\dfrac{4051}{2m+1}>\dfrac{2026}{m+1}$,$4051(m+1)>4052m+2026,m<2025$，成立.}
\step{又各点横纵坐标均为正，都在第一象限内，符合定义，}
\step{所以满足条件的正整数 $n$ 存在，其值为 $n=4051$.}
\solend
\fi

\ansspace*[10]

\end{enumerate}

\end{document}
"
% 紧凑版：xelatex "\def\iscompact{1}% 高一59_华实高中_10月月考.tex
%   —— 高一上数学“2026-2027 年华实高中高一上 10 月月考”（讲评版）
% 试卷版：xelatex 高一59_华实高中_10月月考.tex
% 答案版：xelatex "\def\isanswer{1}% 高一59_华实高中_10月月考.tex
%   —— 高一上数学“2026-2027 年华实高中高一上 10 月月考”（讲评版）
% 试卷版：xelatex 高一59_华实高中_10月月考.tex
% 答案版：xelatex "\def\isanswer{1}\input{高一59_华实高中_10月月考.tex}"
% 紧凑版：xelatex "\def\iscompact{1}\input{高一59_华实高中_10月月考.tex}"
%
% 原始材料：微信公众号“魔都数学加油站”文章，作者破晓，连载“27学年高一59”；
% 原文标题“27学年高一59——2026-2027年上海市华实高中高一上10月月考”；
% 原文链接：https://mp.weixin.qq.com/s/XzyBCmrvpRuMe-NuNk7gkA。
% 页面标注发布时间 2026-10-13 21:45（上海）。
% 正文为 9 张 PNG 页：第 01、02、03、05、06、07、08、09 页 4762×6735，第 04 页 2560×3620
% （**同一篇各页分辨率不同**），均按页序读取。9 页全为试卷正文页（卷首标题在第 1 页页顶，
% 卷末解答题第 21 题解析止于第 9 页），无封面页、目录页、二维码或宣传页。
% 粉色斜排“魔都数学加油站”为卷面水印，与公众号页脚一并未录入。
%
% 卷首标题：照卷面印作“2026--2027 年华实高中高一上 10 月月考”（公众号标题中的“上海市”
% 卷面标题行没有，本稿照卷面），年份区间按全稿惯例排 en dash。原卷只有一行标题、无副标题。
%
% 篇幅与题号：原卷自第 3 题起编号（第 1、2 题不在本卷内）。填空题 3—12（10 题）、
% 选择题 13—16（4 题）、解答题 17—21（5 题），共 19 题。本稿沿用原节与原题号，
% 只改排为全稿统一的大节样式（原卷小节标题为“一、填空题”“二、选择题”“三、解答题”）。
%
% 讲评版：原卷每题下直接印【答案】或【解析】。**只印【答案】的 1 题**（第 3 题），本稿只给 \ans{}；
% 其余 18 题都印【解析】，按全稿口径这类题不另提 \ans{}——第 13、14、15、16 题的结论写在解析末句
% （“故选 D”等），亦照录在解析内。本卷**没有**印在题干括号内的答案。
%
% 副标题：原卷未印副标题。本稿按“知识模块”口径标作“集合与常用逻辑用语、不等式”：
% 19 题中集合与常用逻辑用语 13 题（第 3、9、10、11、12、14、15、16、17、18、19、20、21 题），
% 不等式 10 题（含等式与不等式性质、一元二次不等式、恒成立、基本应用：第 4、5、6、7、8、
% 13、17、18、19、20 题），两类均远超三分之一，故并标。口径同上。
%
% 与原卷的排版差异：
%   ① 原卷每题下直接印【答案】/【解析】，本稿按全稿惯例将答案与解析置于答案版，试卷版保持干净；
%   ② 原卷未印副标题，本稿按“知识模块”口径补出；
%   ③ 原卷填空横线约两个字宽，本稿按全稿惯例统一排 \blankline[4.4em]；
%   ④ 原卷实数集排粗体 $\mathbf{R}$（第 12、13、20 题等），整数集排粗体 $\mathbf{Z}$（第 11、12 题），
%      自然数集排斜体 $N$ 或粗体 $\mathbf{Z}$ 处按本稿惯例改用 $\R$、$\Z$、$\N$；
%      空集记号统一为 \varnothing；
%   ⑤ 原卷在数学式内多处用半角逗号（如 $x_1,x_2$、$a,b\in R$），均照录未强改；
%      句末句点按全稿惯例排西文句点；“漂亮集”“上位点”“下位点”等自定义名词照录；
%   ⑥ 第 17—21 题的子问原卷各自另起一行，本稿按全稿惯例用 enumerate 的 (1)(2)(3)；
%      解析中的子问标号一律用 \stepm{(1)}{…}；第 9、11 题解析中的①、②（③④）分情形或分结论，
%      按全稿惯例作 \stepm{①}{…}、\stepm{②}{…}（第 11 题解析末“故正确的是②④”照录为行内）；
%   ⑦ 第 12 题题干里的“$A_p$”下标原卷作斜体 $p$，照录；
%   ⑧ 每道解答题原卷未留作答空白，本稿按全稿惯例加 \ansspace。
%
% 原卷笔误/存疑（均照录未改，正文对应处另有 % 注）：本卷首轮逐页转录，暂未发现影响答案的笔误。
%
% 插图：无（全卷 9 页无任何图形、无“如图”字样）。
%
\documentclass[a4paper,11pt]{article}
\usepackage{common}

\begin{document}

\examtitlest[3]{2026--2027 年华实高中高一上 10 月月考}{集合与常用逻辑用语、不等式}

\bigsec{一、填空题}

\begin{enumerate}
\setcounter{enumi}{2}

\item “$x\geqslant1$ 且 $y\geqslant1$”的否定形式为\blankline[4.4em].

\ifanswer
\ans{$x<1$ 或 $y<1$}
\fi

\item 已知方程 $2x^2+4x-3=0$ 的两个根为 $x_1,x_2$，则 $\dfrac1{x_1}+\dfrac1{x_2}=$\blankline[4.4em].

\ifanswer
\solbegin
\step{利用韦达定理得 $\begin{cases}x_1+x_2=-2\\ x_1\cdot x_2=-\dfrac32\end{cases}$，则 $\dfrac1{x_1}+\dfrac1{x_2}=\dfrac{x_1+x_2}{x_1\cdot x_2}=\dfrac43$.}
\solend
\fi

\item 设 $a>b>0$，比较大小：$\dfrac1{a-b}$\blankline[2.4em]$\dfrac1a$.

\ifanswer
\solbegin
\step{由 $a>b>0$ 得 $a>0$ 且 $a-b>0$，所以 $\dfrac1{a-b}-\dfrac1a=\dfrac b{a(a-b)}>0$，}
\step{故 $\dfrac1{a-b}>\dfrac1a$.}
\solend
\fi

\item 若不等式 $2ax^2-ax+1>0$ 的解集为 $R$，则实数 $a$ 的取值范围是\blankline[4.4em].

\ifanswer
\solbegin
\step{当 $a\neq0$ 时，由题意得 $\begin{cases}a>0\\ \Delta<0\end{cases}$，即 $\begin{cases}a>0\\ a^2-8a<0\end{cases}$，解得 $0<a<8$.}
\step{当 $a=0$ 时，恒有 $1>0$，不等式也成立.}
\step{综上所述，$a$ 的取值范围是 $[0,8)$.}
\solend
\fi

\item 集合 $A=\{x\mid m-1\leqslant x\leqslant m+7\}$，集合 $B=\{x\mid x^2-3x-10\geqslant0\}$，若 $A\cup B=R$，则实数 $m$ 的取值范围为\blankline[4.4em].

\ifanswer
\solbegin
\step{$A=\{x\mid m-1\leqslant x\leqslant m+7\}$，$B=\{x\mid x^2-3x-10\geqslant0\}=\{x\mid x\leqslant-2\text{ 或 }x\geqslant5\}$，}
\step{若 $A\cup B=R$，则 $\begin{cases}m-1\leqslant-2\\ m+7\geqslant5\end{cases}$，解得 $-2\leqslant m\leqslant-1$，}
\step{所以实数 $m$ 的取值范围为 $[-2,-1]$.}
\solend
\fi

\item 若关于 $x$ 的不等式 $ax^2+bx+c>0$ 的解集为 $(-1,3)$，则关于 $x$ 的不等式 $cx^2+bx+a>0$ 的解集为\blankline[4.4em].

\ifanswer
\solbegin
\step{因为不等式 $ax^2+bx+c>0$ 的解集为 $(-1,3)$，}
\step{所以 $-1$ 和 $3$ 是方程 $ax^2+bx+c=0$ 的实数解，}
\step{所以 $\begin{cases}-1+3=-\dfrac ba\\ -1\times3=\dfrac ca\\ a<0\end{cases}$，解得 $b=-2a$，$c=-3a$，且 $a<0$，}
\step{所以不等式 $cx^2+bx+a>0$ 可化为 $-3ax^2-2ax+a>0$，即 $3x^2+2x-1>0$，}
\step{解得 $x<-1$ 或 $x>\dfrac13$，所以该不等式的解集为 $(-\infty,-1)\cup\left(\dfrac13,+\infty\right)$.}
\solend
\fi

\item 若集合 $A=\{x\mid(a-1)x^2-2x+1=0\}$ 有且仅有 $2$ 个子集，则实数 $a$ 的值为\blankline[4.4em].

\ifanswer
\solbegin
\step{由题意得方程 $(a-1)x^2-2x+1=0$ 只有一个根，}
\stepm{①}{当 $a-1=0$ 时，$a=1$，此时方程化为 $-2x+1=0$，解得 $x=\dfrac12$，符合题意，}
\stepm{②}{当 $a-1\neq0$，即 $a\neq1$ 时，$\Delta=4-4(a-1)=0$，解得 $a=2$，}
\step{综上所述，实数 $a$ 的值为 $1$ 或 $2$.}
\solend
\fi

\item “$x+y\neq5$”是“$x\neq3$ 或 $y\neq2$”的\blankline[4.4em]条件.

\ifanswer
\solbegin
\step{分析充分性：}
\step{若 $x+y\neq5$ 成立，若 $x=3$，则 $y\neq2$，若 $y=2$，则 $x\neq3$，}
\step{所以一定能说明 $x\neq3$ 和 $y\neq2$ 至少有一个成立，即 $x\neq3$ 或 $y\neq2$ 成立，}
\step{充分性成立；}
\step{分析必要性：}
\step{若 $x\neq3$ 或 $y\neq2$，即 $x\neq3$ 和 $y\neq2$ 二者至少一个成立，}
\step{举例：$x=1,y=4$，此时 $x+y=5$，所以必要性不成立.}
\step{所以“$x+y\neq5$”是“$x\neq3$ 或 $y\neq2$”的充分不必要条件.}
\solend
\fi

\item 设 $P$ 为非空实数集满足：对任意给定的 $x,y\in P$（$x,y$ 可以相同），都有 $x+y\in P,x-y\in P,xy\in P$，则称 $P$ 为漂亮集. 关于漂亮集有下列结论：

①集合 $P=\{-2,-1,0,1,2\}$ 为漂亮集；②集合 $P=\{x\mid x=2n,n\in\Z\}$ 为漂亮集；

③若集合 $P_1,P_2$ 为漂亮集，则 $P_1\cup P_2$ 为漂亮集；④若集合 $P$ 为漂亮集，则一定有 $0\in P$.

其中所有正确结论的序号是\blankline[4.4em].

\ifanswer
\solbegin
\step{对于集合 $P=\{-2,-1,0,1,2\}$，取 $x=2,y=-2$，则 $x-y=2-(-2)=4\notin P$，}
\step{不满足漂亮集的定义，所以①错误.}
\step{对于集合 $P=\{x\mid x=2n,n\in Z\}$，设 $x=2n_1,y=2n_2,n_1,n_2\in Z$.}
\step{$x+y=2n_1+2n_2=2(n_1+n_2)$，因为 $n_1+n_2\in Z$，所以 $x+y\in P$.}
\step{$x-y=2n_1-2n_2=2(n_1-n_2)$，因为 $n_1-n_2\in Z$，所以 $x-y\in P$.}
\step{$xy=2n_1\times2n_2=4n_1n_2=2(2n_1n_2)$，因为 $2n_1n_2\in Z$，所以 $xy\in P$.}
\step{满足漂亮集的定义，所以②正确.}
\step{令 $P_1=\{x\mid x=2n,n\in Z\}$,$P_2=\{x\mid x=3n,n\in Z\}$，都是漂亮集.}
\step{$P_1\cup P_2$ 中，取 $x=2(x\in P_1)$,$y=3(y\in P_2)$,$x+y=5\notin P_1\cup P_2$，}
\step{不满足漂亮集的定义，所以③错误.}
\step{因为对任意 $x\in P$，有 $x-x=0\in P$，所以若集合 $P$ 为漂亮集，}
\step{则一定有 $0\in P$，所以④正确.}
\step{故正确的是②④.}
\solend
\fi

\item 已知：集合 $A_p=\{(x,y)\mid y=3x^2-px+5p-65,x\in\R\}$,$M=\{x\mid1\leqslant x\leqslant100,x\in\Z\}$，若 $p\in M$，则所有 $A_p$ 的交集\blankline[4.4em].

\ifanswer
\solbegin
\step{由 $M=\{1,2,\cdots,100\}$，所以所有 $A_p$ 的交集就是这些集合中共同拥有的点.}
\step{若 $(x,y)$ 是所有集合的一个公共点，}
\step{那么它特别要同时属于 $p=1$ 和 $p=2$ 时对应的两个集合.}
\step{当 $p=1$ 时，$y=3x^2-x-60$；}
\step{当 $p=2$ 时，$y=3x^2-2x-55$；联立得 $x=5,y=10$.}
\step{因此，如果公共点存在，只可能是 $(5,10)$.}
\step{接下来验证这个点是否属于所有集合，}
\step{对任意 $p\in M$，把 $x=5$ 代入集合的关系式，}
\step{得 $3x^2-px+5p-65=75-5p+5p-65=10$，与 $y=10$ 相符，}
\step{所以 $(5,10)\in A_p$ 对每一个 $p$ 都成立，因此所有 $A_p$ 的交集为 $\{(5,10)\}$.}
\solend
\fi

\end{enumerate}

% 守卫须放在 \bigsec 之前，否则标题仍会落页末（见 高一40 的教训）
\needvspace{5cm}

\bigsec{二、选择题}

\begin{enumerate}
\setcounter{enumi}{12}

\item 已知 $a,b,c\in\R$ 且 $a>b$，则下列不等式一定成立的是\mbox{（\quad）}

\keepwithprev
A. $\dfrac1a<\dfrac1b$\qquad B. $a^2>b^2$

C. $a|c|>b|c|$\qquad D. $\dfrac a{c^2+1}>\dfrac b{c^2+1}$

\ifanswer
\solbegin
\step{$A$：当 $a>0>b$ 时，$\dfrac1a>\dfrac1b$，故 $A$ 错误；}
\step{$B$：当 $a=-1,b=-2$ 时，满足 $a>b$，但 $a^2>b^2$ 不成立，故 $B$ 错误；}
\step{$C$：当 $c=0$ 时，$a|c|=b|c|$，故 $C$ 错误；}
\step{$D$：由 $a>b,c^2+1>0$，得 $\dfrac a{c^2+1}>\dfrac b{c^2+1}$，故 $D$ 正确；}
\step{故选 $D$.}
\solend
\fi

\item 已知集合 $A=\{a^2,0,-1\}$，$B=\{a,b,0\}$，若 $A=B$，则 $(ab)^{2026}$ 的值为\mbox{（\quad）}

\keepwithprev
A. $0$\qquad B. $-1$

C. $1$\qquad D. $\pm1$

\ifanswer
\solbegin
\step{因为集合 $A=\{a^2,0,-1\}$，$B=\{a,b,0\}$，$A=B$，所以 $\begin{cases}a=-1\\ a^2=b\end{cases}$，}
\step{所以 $a=-1$，$b=1$，所以 $(ab)^{2026}=1$，故选 $C$.}
\solend
\fi

\item 设全集 $U=\R$，$A=\{x\mid x<-1\text{ 或 }x>2\}$,$B=\{y\mid y=x^2+a\}$，若 $\overline A\subseteq B$，则实数 $a$ 的取值范围是\mbox{（\quad）}

\keepwithprev
A. $a\leqslant-1$\qquad B. $a>2$

C. $a<-1$\qquad D. $a\geqslant-2$

\ifanswer
\solbegin
\step{因为 $A=\{x\mid x<-1\text{ 或 }x>2\}$，则 $\overline A=\{x\mid-1\leqslant x\leqslant2\}$，}
\step{又 $B=\{y\mid y=x^2+a\}=\{y\mid y\geqslant a\}$，且 $\overline A\subseteq B$，所以 $a\leqslant-1$，故选 $A$.}
\solend
\fi

\item 设 $P$ 和 $Q$ 是两个集合，定义集合 $P-Q=\{x\mid x\in P$，且 $x\notin Q\}$，如果 $Q=\{x\mid(1-x)(x-3)>0\}$，$P=\left\{x\left|\dfrac x{2-x}>0\right.\right\}$，那么 $P-Q$ 等于\mbox{（\quad）}

\keepwithprev
A. $\{x\mid0<x<1\}$\qquad B. $\{x\mid1\leqslant x<2\}$

C. $\{x\mid0<x\leqslant1\}$\qquad D. $\{x\mid2\leqslant x<3\}$

\ifanswer
\solbegin
\step{因为 $P=\left\{x\left|\dfrac x{2-x}>0\right.\right\}=\{x\mid x(x-2)<0\}=\{x\mid0<x<2\}$，}
\step{$Q=\{x\mid(1-x)(x-3)>0\}=\{x\mid(x-3)(x-1)<0\}=\{x\mid1<x<3\}$，}
\step{所以 $P-Q=\{x\mid0<x\leqslant1\}$，故选 $C$.}
\solend
\fi

\end{enumerate}

\bigsec{三、解答题}

\begin{enumerate}
\setcounter{enumi}{16}

\item 已知集合 $A=\{x\mid x^2-9\geqslant0\}$，$B=\{x\mid|x-4|<2\}$，$C=\left\{x\left|\dfrac{x-8}{x+2}<0\right.\right\}$.

（1）求 $A\cup C$；

（2）若全集 $U=R$，求 $\overline A\cap B$.

\ifanswer
\solbegin
\stepm{(1)}{集合 $A=\{x\mid x^2-9\geqslant0\}=\{x\mid x\geqslant3\text{ 或 }x\leqslant-3\}$，}
\step{$B=\{x\mid|x-4|<2\}=\{x\mid2<x<6\}$，$C=\left\{x\left|\dfrac{x-8}{x+2}<0\right.\right\}=\{x\mid-2<x<8\}$，}
\step{所以 $A\cup C=\{x\mid x\leqslant-3\text{ 或 }x>-2\}$；}
\stepm{(2)}{全集 $U=R$，集合 $A=\{x\mid x^2-9\geqslant0\}=\{x\mid x\geqslant3\text{ 或 }x\leqslant-3\}$，}
\step{$B=\{x\mid|x-4|<2\}=\{x\mid2<x<6\}$，}
\step{所以 $\overline A=\{x\mid-3<x<3\}$，所以 $\overline A\cap B=\{x\mid2<x<3\}$.}
\solend
\fi

\ansspace[6]

\item 已知集合 $A=\left\{x\left|\dfrac{2x-1}{x+1}<1,x\in\R\right.\right\}$，集合 $B=\{x\mid(x-a-1)(x-a+1)\leqslant0\}$.

（1）求集合 $A$；

（2）若 $B\subseteq A$，求实数 $a$ 的取值范围.

\ifanswer
\solbegin
\stepm{(1)}{由 $\dfrac{2x-1}{x+1}<1$，得 $\dfrac{x-2}{x+1}<0$ 所以 $(x-2)(x+1)<0$，解得 $-1<x<2$，}
\step{即 $A=(-1,2)$；}
\stepm{(2)}{由 $(x-a-1)(x-a+1)\leqslant0$，得 $a-1\leqslant x\leqslant a+1$，}
\step{所以 $B=\{x\mid a-1\leqslant x\leqslant a+1\}$，}
\step{又 $B\subseteq A$，且 $a+1>a-1$ 即 $B\neq\varnothing$，所以 $\begin{cases}a-1>-1\\ a+1<2\end{cases}$，所以 $0<a<1$，}
\step{所以实数 $a$ 的取值范围是 $\{a\mid0<a<1\}$.}
\solend
\fi

\ansspace[6]

\item 已知集合 $A=\{x\mid x^2-4x+3=0\}$，$B=\{x\mid x^2-ax+3=0\}$.

（1）若 $A\cup B=B$，求实数 $a$ 的值；

（2）若 $A\cap B=B$，求实数 $a$ 的取值范围.

\ifanswer
\solbegin
\stepm{(1)}{集合 $A=\{x\mid x^2-4x+3=0\}=\{1,3\}$，$B=\{x\mid x^2-ax+3=0\}$.}
\step{$A\cup B=B$，所以 $1\in B$，$3\in B$，所以 $\begin{cases}1-a+3=0\\ 9-3a+3=0\end{cases}$，解得实数 $a=4$；}
\stepm{(2)}{因为 $A\cap B=B$，所以 $B\subseteq A$，}
\step{当 $B=\varnothing$ 时，$\Delta=(-a)^2-12<0$，解得 $-2\sqrt3<a<2\sqrt3$，}
\step{当 $B$ 是单元素集时，$\Delta=a^2-12=0$，解得 $a=2\sqrt3$，此时 $B=\{\sqrt3\}$，}
\step{不合题意；}
\step{当 $B=\{1,3\}$ 时，由 (1) 得 $a=4$，}
\step{综上，实数 $a$ 的取值范围是 $(-2\sqrt3,2\sqrt3)\cup\{4\}$.}
\solend
\fi

\ansspace[6]

\item 已知全集为 $R$，实数 $a$、$b$ 满足 $a\neq b$，集合 $M=\{x\mid x^2-3x-4\leqslant0\}$，集合 $A=\{x\mid(x+a)(x+b)>0\}$，集合 $B=\{x\mid x^2+(a-2)x-2a>0\}$.

（1）若 $\overline A=M$，求 $a$、$b$ 的值；

（2）若 $a>b>-1$，求 $A\cap B$；

（3）若 $a^2+a\in\overline B$，求 $a$ 的取值范围.

\ifanswer
\solbegin
\stepm{(1)}{$M=\{x\mid x^2-3x-4\leqslant0\}=\{x\mid-1\leqslant x\leqslant4\}$，$A=\{x\mid(x+a)(x+b)\leqslant0\}$，}
\step{因为 $\overline A=M$，所以 $a=1$，$b=-4$ 或 $a=-4$，$b=1$.}
\stepm{(2)}{因为 $a>b>-1$，所以 $-a<-b<1$，所以 $A=\{x\mid x<-a\text{ 或 }x>-b\}$，}
\step{因为 $B=\{x\mid x^2+(a-2)x-2a>0\}=\{x\mid x<-a\text{ 或 }x>2\}$，}
\step{所以 $A\cap B=\{x\mid x<-a\text{ 或 }x>2\}$，}
\stepm{(3)}{$\overline B=\{x\mid(x-2)(x+a)\leqslant0\}$，}
\step{由 $a^2+a\in\overline B$，得 $(a^2+a-2)(a^2+2a)\leqslant0\Leftrightarrow a(a-1)(a-2)^2\leqslant0$，}
\step{解得 $a=2$ 或 $0\leqslant a\leqslant1$，故 $a$ 的取值范围是 $[0,1]\cup\{2\}$.}
\solend
\fi

\ansspace[8]

\item 对平面直角坐标系的第一象限内的任意两点作如下定义：若 $\dfrac ab>\dfrac cd$，那么称点 $(a,b)$ 是点 $(c,d)$ 的“上位点”，同时点 $(c,d)$ 是点 $(a,b)$ 的“下位点”.

（1）试写出点 $(2,4)$ 的一个“上位点”坐标；

（2）已知点 $(a,b)$ 是点 $(c,d)$ 的“上位点”，判断点 $P(a+c,b+d)$ 是点 $(c,d)$ 的“上位点”还是“下位点”，并证明你的结论；

（3）设正整数 $n$ 满足以下条件：对集合 $\{t\mid0<t<2020,t\in\Z\}$ 内的任意元素 $m$，总存在正整数 $k=2m+1$，使得点 $(n,k)$ 既是点 $(2025,m)$ 的“下位点”，又是点 $(2026,m+1)$ 的“上位点”，问这样的正整数 $n$ 是否存在？若存在，求正整数 $n$ 的值；若不存在，请说明理由.

\ifanswer
\solbegin
\stepm{(1)}{由定义，点 $(a,b)$ 是点 $(2,4)$ 的“上位点”，}
\step{当且仅当 $\dfrac ab>\dfrac24=\dfrac12$ 且 $a>0,b>0$.}
\step{取 $(a,b)=(1,1)$，因为 $\dfrac11>\dfrac12$，}
\step{所以点 $(1,1)$ 是点 $(2,4)$ 的一个“上位点”.（答案不唯一）}
\stepm{(2)}{点 $P(a+c,b+d)$ 是点 $(c,d)$ 的“上位点”.}
\step{$(a,b)$ 是点 $(c,d)$ 的“上位点”，所以 $\dfrac ab>\dfrac cd$.}
\step{因为 $a,b,c,d>0$，两边同乘 $bd$ 得 $ad>bc$，即 $ad-bc>0$.}
\step{作差比较 $\dfrac{a+c}{b+d}$ 与 $\dfrac cd$：$\dfrac{a+c}{b+d}-\dfrac cd=\dfrac{d(a+c)-c(b+d)}{d(b+d)}=\dfrac{ad-bc}{d(b+d)}>0$，}
\step{所以点 $P(a+c,b+d)$ 是点 $(c,d)$ 的“上位点”.}
\stepm{(3)}{由定义，点 $(n,k)$ 是点 $(2025,m)$ 的下位点即 $\dfrac{2025}m>\dfrac nk$；}
\step{点 $(n,k)$ 是点 $(2026,m+1)$ 的上位点即 $\dfrac nk>\dfrac{2026}{m+1}$.}
\step{代入 $k=2m+1$，得对任意 $m\in\{1,2,\cdots,2019\}$，}
\step{都有 $4052-\dfrac{2026}{m+1}<n<4050+\dfrac{2025}m$.}
\step{右边 $4050+\dfrac{2025}m$ 随 $m$ 增大而减小，故最近处在 $m=2019$，}
\step{得 $n<4050+\dfrac{2025}{2019}<4051.01$，所以 $n\leqslant4051$.}
\step{左边 $4052-\dfrac{2026}{m+1}$ 随 $m$ 增大而增大，故最近处也在 $m=2019$，}
\step{得 $n>4052-\dfrac{2026}{2020}>4050.99$，所以 $n\geqslant4051$.}
\step{综上，$n=4051$.}
\step{检验：当 $n=4051$ 时，对任意 $1\leqslant m\leqslant2019,k=2m+1$ 是正整数，}
\step{且 $\dfrac{4051}{2m+1}<\dfrac{2025}m$,$4051m<4050m+2025,m<2025$，成立；}
\step{$\dfrac{4051}{2m+1}>\dfrac{2026}{m+1}$,$4051(m+1)>4052m+2026,m<2025$，成立.}
\step{又各点横纵坐标均为正，都在第一象限内，符合定义，}
\step{所以满足条件的正整数 $n$ 存在，其值为 $n=4051$.}
\solend
\fi

\ansspace*[10]

\end{enumerate}

\end{document}
"
% 紧凑版：xelatex "\def\iscompact{1}% 高一59_华实高中_10月月考.tex
%   —— 高一上数学“2026-2027 年华实高中高一上 10 月月考”（讲评版）
% 试卷版：xelatex 高一59_华实高中_10月月考.tex
% 答案版：xelatex "\def\isanswer{1}\input{高一59_华实高中_10月月考.tex}"
% 紧凑版：xelatex "\def\iscompact{1}\input{高一59_华实高中_10月月考.tex}"
%
% 原始材料：微信公众号“魔都数学加油站”文章，作者破晓，连载“27学年高一59”；
% 原文标题“27学年高一59——2026-2027年上海市华实高中高一上10月月考”；
% 原文链接：https://mp.weixin.qq.com/s/XzyBCmrvpRuMe-NuNk7gkA。
% 页面标注发布时间 2026-10-13 21:45（上海）。
% 正文为 9 张 PNG 页：第 01、02、03、05、06、07、08、09 页 4762×6735，第 04 页 2560×3620
% （**同一篇各页分辨率不同**），均按页序读取。9 页全为试卷正文页（卷首标题在第 1 页页顶，
% 卷末解答题第 21 题解析止于第 9 页），无封面页、目录页、二维码或宣传页。
% 粉色斜排“魔都数学加油站”为卷面水印，与公众号页脚一并未录入。
%
% 卷首标题：照卷面印作“2026--2027 年华实高中高一上 10 月月考”（公众号标题中的“上海市”
% 卷面标题行没有，本稿照卷面），年份区间按全稿惯例排 en dash。原卷只有一行标题、无副标题。
%
% 篇幅与题号：原卷自第 3 题起编号（第 1、2 题不在本卷内）。填空题 3—12（10 题）、
% 选择题 13—16（4 题）、解答题 17—21（5 题），共 19 题。本稿沿用原节与原题号，
% 只改排为全稿统一的大节样式（原卷小节标题为“一、填空题”“二、选择题”“三、解答题”）。
%
% 讲评版：原卷每题下直接印【答案】或【解析】。**只印【答案】的 1 题**（第 3 题），本稿只给 \ans{}；
% 其余 18 题都印【解析】，按全稿口径这类题不另提 \ans{}——第 13、14、15、16 题的结论写在解析末句
% （“故选 D”等），亦照录在解析内。本卷**没有**印在题干括号内的答案。
%
% 副标题：原卷未印副标题。本稿按“知识模块”口径标作“集合与常用逻辑用语、不等式”：
% 19 题中集合与常用逻辑用语 13 题（第 3、9、10、11、12、14、15、16、17、18、19、20、21 题），
% 不等式 10 题（含等式与不等式性质、一元二次不等式、恒成立、基本应用：第 4、5、6、7、8、
% 13、17、18、19、20 题），两类均远超三分之一，故并标。口径同上。
%
% 与原卷的排版差异：
%   ① 原卷每题下直接印【答案】/【解析】，本稿按全稿惯例将答案与解析置于答案版，试卷版保持干净；
%   ② 原卷未印副标题，本稿按“知识模块”口径补出；
%   ③ 原卷填空横线约两个字宽，本稿按全稿惯例统一排 \blankline[4.4em]；
%   ④ 原卷实数集排粗体 $\mathbf{R}$（第 12、13、20 题等），整数集排粗体 $\mathbf{Z}$（第 11、12 题），
%      自然数集排斜体 $N$ 或粗体 $\mathbf{Z}$ 处按本稿惯例改用 $\R$、$\Z$、$\N$；
%      空集记号统一为 \varnothing；
%   ⑤ 原卷在数学式内多处用半角逗号（如 $x_1,x_2$、$a,b\in R$），均照录未强改；
%      句末句点按全稿惯例排西文句点；“漂亮集”“上位点”“下位点”等自定义名词照录；
%   ⑥ 第 17—21 题的子问原卷各自另起一行，本稿按全稿惯例用 enumerate 的 (1)(2)(3)；
%      解析中的子问标号一律用 \stepm{(1)}{…}；第 9、11 题解析中的①、②（③④）分情形或分结论，
%      按全稿惯例作 \stepm{①}{…}、\stepm{②}{…}（第 11 题解析末“故正确的是②④”照录为行内）；
%   ⑦ 第 12 题题干里的“$A_p$”下标原卷作斜体 $p$，照录；
%   ⑧ 每道解答题原卷未留作答空白，本稿按全稿惯例加 \ansspace。
%
% 原卷笔误/存疑（均照录未改，正文对应处另有 % 注）：本卷首轮逐页转录，暂未发现影响答案的笔误。
%
% 插图：无（全卷 9 页无任何图形、无“如图”字样）。
%
\documentclass[a4paper,11pt]{article}
\usepackage{common}

\begin{document}

\examtitlest[3]{2026--2027 年华实高中高一上 10 月月考}{集合与常用逻辑用语、不等式}

\bigsec{一、填空题}

\begin{enumerate}
\setcounter{enumi}{2}

\item “$x\geqslant1$ 且 $y\geqslant1$”的否定形式为\blankline[4.4em].

\ifanswer
\ans{$x<1$ 或 $y<1$}
\fi

\item 已知方程 $2x^2+4x-3=0$ 的两个根为 $x_1,x_2$，则 $\dfrac1{x_1}+\dfrac1{x_2}=$\blankline[4.4em].

\ifanswer
\solbegin
\step{利用韦达定理得 $\begin{cases}x_1+x_2=-2\\ x_1\cdot x_2=-\dfrac32\end{cases}$，则 $\dfrac1{x_1}+\dfrac1{x_2}=\dfrac{x_1+x_2}{x_1\cdot x_2}=\dfrac43$.}
\solend
\fi

\item 设 $a>b>0$，比较大小：$\dfrac1{a-b}$\blankline[2.4em]$\dfrac1a$.

\ifanswer
\solbegin
\step{由 $a>b>0$ 得 $a>0$ 且 $a-b>0$，所以 $\dfrac1{a-b}-\dfrac1a=\dfrac b{a(a-b)}>0$，}
\step{故 $\dfrac1{a-b}>\dfrac1a$.}
\solend
\fi

\item 若不等式 $2ax^2-ax+1>0$ 的解集为 $R$，则实数 $a$ 的取值范围是\blankline[4.4em].

\ifanswer
\solbegin
\step{当 $a\neq0$ 时，由题意得 $\begin{cases}a>0\\ \Delta<0\end{cases}$，即 $\begin{cases}a>0\\ a^2-8a<0\end{cases}$，解得 $0<a<8$.}
\step{当 $a=0$ 时，恒有 $1>0$，不等式也成立.}
\step{综上所述，$a$ 的取值范围是 $[0,8)$.}
\solend
\fi

\item 集合 $A=\{x\mid m-1\leqslant x\leqslant m+7\}$，集合 $B=\{x\mid x^2-3x-10\geqslant0\}$，若 $A\cup B=R$，则实数 $m$ 的取值范围为\blankline[4.4em].

\ifanswer
\solbegin
\step{$A=\{x\mid m-1\leqslant x\leqslant m+7\}$，$B=\{x\mid x^2-3x-10\geqslant0\}=\{x\mid x\leqslant-2\text{ 或 }x\geqslant5\}$，}
\step{若 $A\cup B=R$，则 $\begin{cases}m-1\leqslant-2\\ m+7\geqslant5\end{cases}$，解得 $-2\leqslant m\leqslant-1$，}
\step{所以实数 $m$ 的取值范围为 $[-2,-1]$.}
\solend
\fi

\item 若关于 $x$ 的不等式 $ax^2+bx+c>0$ 的解集为 $(-1,3)$，则关于 $x$ 的不等式 $cx^2+bx+a>0$ 的解集为\blankline[4.4em].

\ifanswer
\solbegin
\step{因为不等式 $ax^2+bx+c>0$ 的解集为 $(-1,3)$，}
\step{所以 $-1$ 和 $3$ 是方程 $ax^2+bx+c=0$ 的实数解，}
\step{所以 $\begin{cases}-1+3=-\dfrac ba\\ -1\times3=\dfrac ca\\ a<0\end{cases}$，解得 $b=-2a$，$c=-3a$，且 $a<0$，}
\step{所以不等式 $cx^2+bx+a>0$ 可化为 $-3ax^2-2ax+a>0$，即 $3x^2+2x-1>0$，}
\step{解得 $x<-1$ 或 $x>\dfrac13$，所以该不等式的解集为 $(-\infty,-1)\cup\left(\dfrac13,+\infty\right)$.}
\solend
\fi

\item 若集合 $A=\{x\mid(a-1)x^2-2x+1=0\}$ 有且仅有 $2$ 个子集，则实数 $a$ 的值为\blankline[4.4em].

\ifanswer
\solbegin
\step{由题意得方程 $(a-1)x^2-2x+1=0$ 只有一个根，}
\stepm{①}{当 $a-1=0$ 时，$a=1$，此时方程化为 $-2x+1=0$，解得 $x=\dfrac12$，符合题意，}
\stepm{②}{当 $a-1\neq0$，即 $a\neq1$ 时，$\Delta=4-4(a-1)=0$，解得 $a=2$，}
\step{综上所述，实数 $a$ 的值为 $1$ 或 $2$.}
\solend
\fi

\item “$x+y\neq5$”是“$x\neq3$ 或 $y\neq2$”的\blankline[4.4em]条件.

\ifanswer
\solbegin
\step{分析充分性：}
\step{若 $x+y\neq5$ 成立，若 $x=3$，则 $y\neq2$，若 $y=2$，则 $x\neq3$，}
\step{所以一定能说明 $x\neq3$ 和 $y\neq2$ 至少有一个成立，即 $x\neq3$ 或 $y\neq2$ 成立，}
\step{充分性成立；}
\step{分析必要性：}
\step{若 $x\neq3$ 或 $y\neq2$，即 $x\neq3$ 和 $y\neq2$ 二者至少一个成立，}
\step{举例：$x=1,y=4$，此时 $x+y=5$，所以必要性不成立.}
\step{所以“$x+y\neq5$”是“$x\neq3$ 或 $y\neq2$”的充分不必要条件.}
\solend
\fi

\item 设 $P$ 为非空实数集满足：对任意给定的 $x,y\in P$（$x,y$ 可以相同），都有 $x+y\in P,x-y\in P,xy\in P$，则称 $P$ 为漂亮集. 关于漂亮集有下列结论：

①集合 $P=\{-2,-1,0,1,2\}$ 为漂亮集；②集合 $P=\{x\mid x=2n,n\in\Z\}$ 为漂亮集；

③若集合 $P_1,P_2$ 为漂亮集，则 $P_1\cup P_2$ 为漂亮集；④若集合 $P$ 为漂亮集，则一定有 $0\in P$.

其中所有正确结论的序号是\blankline[4.4em].

\ifanswer
\solbegin
\step{对于集合 $P=\{-2,-1,0,1,2\}$，取 $x=2,y=-2$，则 $x-y=2-(-2)=4\notin P$，}
\step{不满足漂亮集的定义，所以①错误.}
\step{对于集合 $P=\{x\mid x=2n,n\in Z\}$，设 $x=2n_1,y=2n_2,n_1,n_2\in Z$.}
\step{$x+y=2n_1+2n_2=2(n_1+n_2)$，因为 $n_1+n_2\in Z$，所以 $x+y\in P$.}
\step{$x-y=2n_1-2n_2=2(n_1-n_2)$，因为 $n_1-n_2\in Z$，所以 $x-y\in P$.}
\step{$xy=2n_1\times2n_2=4n_1n_2=2(2n_1n_2)$，因为 $2n_1n_2\in Z$，所以 $xy\in P$.}
\step{满足漂亮集的定义，所以②正确.}
\step{令 $P_1=\{x\mid x=2n,n\in Z\}$,$P_2=\{x\mid x=3n,n\in Z\}$，都是漂亮集.}
\step{$P_1\cup P_2$ 中，取 $x=2(x\in P_1)$,$y=3(y\in P_2)$,$x+y=5\notin P_1\cup P_2$，}
\step{不满足漂亮集的定义，所以③错误.}
\step{因为对任意 $x\in P$，有 $x-x=0\in P$，所以若集合 $P$ 为漂亮集，}
\step{则一定有 $0\in P$，所以④正确.}
\step{故正确的是②④.}
\solend
\fi

\item 已知：集合 $A_p=\{(x,y)\mid y=3x^2-px+5p-65,x\in\R\}$,$M=\{x\mid1\leqslant x\leqslant100,x\in\Z\}$，若 $p\in M$，则所有 $A_p$ 的交集\blankline[4.4em].

\ifanswer
\solbegin
\step{由 $M=\{1,2,\cdots,100\}$，所以所有 $A_p$ 的交集就是这些集合中共同拥有的点.}
\step{若 $(x,y)$ 是所有集合的一个公共点，}
\step{那么它特别要同时属于 $p=1$ 和 $p=2$ 时对应的两个集合.}
\step{当 $p=1$ 时，$y=3x^2-x-60$；}
\step{当 $p=2$ 时，$y=3x^2-2x-55$；联立得 $x=5,y=10$.}
\step{因此，如果公共点存在，只可能是 $(5,10)$.}
\step{接下来验证这个点是否属于所有集合，}
\step{对任意 $p\in M$，把 $x=5$ 代入集合的关系式，}
\step{得 $3x^2-px+5p-65=75-5p+5p-65=10$，与 $y=10$ 相符，}
\step{所以 $(5,10)\in A_p$ 对每一个 $p$ 都成立，因此所有 $A_p$ 的交集为 $\{(5,10)\}$.}
\solend
\fi

\end{enumerate}

% 守卫须放在 \bigsec 之前，否则标题仍会落页末（见 高一40 的教训）
\needvspace{5cm}

\bigsec{二、选择题}

\begin{enumerate}
\setcounter{enumi}{12}

\item 已知 $a,b,c\in\R$ 且 $a>b$，则下列不等式一定成立的是\mbox{（\quad）}

\keepwithprev
A. $\dfrac1a<\dfrac1b$\qquad B. $a^2>b^2$

C. $a|c|>b|c|$\qquad D. $\dfrac a{c^2+1}>\dfrac b{c^2+1}$

\ifanswer
\solbegin
\step{$A$：当 $a>0>b$ 时，$\dfrac1a>\dfrac1b$，故 $A$ 错误；}
\step{$B$：当 $a=-1,b=-2$ 时，满足 $a>b$，但 $a^2>b^2$ 不成立，故 $B$ 错误；}
\step{$C$：当 $c=0$ 时，$a|c|=b|c|$，故 $C$ 错误；}
\step{$D$：由 $a>b,c^2+1>0$，得 $\dfrac a{c^2+1}>\dfrac b{c^2+1}$，故 $D$ 正确；}
\step{故选 $D$.}
\solend
\fi

\item 已知集合 $A=\{a^2,0,-1\}$，$B=\{a,b,0\}$，若 $A=B$，则 $(ab)^{2026}$ 的值为\mbox{（\quad）}

\keepwithprev
A. $0$\qquad B. $-1$

C. $1$\qquad D. $\pm1$

\ifanswer
\solbegin
\step{因为集合 $A=\{a^2,0,-1\}$，$B=\{a,b,0\}$，$A=B$，所以 $\begin{cases}a=-1\\ a^2=b\end{cases}$，}
\step{所以 $a=-1$，$b=1$，所以 $(ab)^{2026}=1$，故选 $C$.}
\solend
\fi

\item 设全集 $U=\R$，$A=\{x\mid x<-1\text{ 或 }x>2\}$,$B=\{y\mid y=x^2+a\}$，若 $\overline A\subseteq B$，则实数 $a$ 的取值范围是\mbox{（\quad）}

\keepwithprev
A. $a\leqslant-1$\qquad B. $a>2$

C. $a<-1$\qquad D. $a\geqslant-2$

\ifanswer
\solbegin
\step{因为 $A=\{x\mid x<-1\text{ 或 }x>2\}$，则 $\overline A=\{x\mid-1\leqslant x\leqslant2\}$，}
\step{又 $B=\{y\mid y=x^2+a\}=\{y\mid y\geqslant a\}$，且 $\overline A\subseteq B$，所以 $a\leqslant-1$，故选 $A$.}
\solend
\fi

\item 设 $P$ 和 $Q$ 是两个集合，定义集合 $P-Q=\{x\mid x\in P$，且 $x\notin Q\}$，如果 $Q=\{x\mid(1-x)(x-3)>0\}$，$P=\left\{x\left|\dfrac x{2-x}>0\right.\right\}$，那么 $P-Q$ 等于\mbox{（\quad）}

\keepwithprev
A. $\{x\mid0<x<1\}$\qquad B. $\{x\mid1\leqslant x<2\}$

C. $\{x\mid0<x\leqslant1\}$\qquad D. $\{x\mid2\leqslant x<3\}$

\ifanswer
\solbegin
\step{因为 $P=\left\{x\left|\dfrac x{2-x}>0\right.\right\}=\{x\mid x(x-2)<0\}=\{x\mid0<x<2\}$，}
\step{$Q=\{x\mid(1-x)(x-3)>0\}=\{x\mid(x-3)(x-1)<0\}=\{x\mid1<x<3\}$，}
\step{所以 $P-Q=\{x\mid0<x\leqslant1\}$，故选 $C$.}
\solend
\fi

\end{enumerate}

\bigsec{三、解答题}

\begin{enumerate}
\setcounter{enumi}{16}

\item 已知集合 $A=\{x\mid x^2-9\geqslant0\}$，$B=\{x\mid|x-4|<2\}$，$C=\left\{x\left|\dfrac{x-8}{x+2}<0\right.\right\}$.

（1）求 $A\cup C$；

（2）若全集 $U=R$，求 $\overline A\cap B$.

\ifanswer
\solbegin
\stepm{(1)}{集合 $A=\{x\mid x^2-9\geqslant0\}=\{x\mid x\geqslant3\text{ 或 }x\leqslant-3\}$，}
\step{$B=\{x\mid|x-4|<2\}=\{x\mid2<x<6\}$，$C=\left\{x\left|\dfrac{x-8}{x+2}<0\right.\right\}=\{x\mid-2<x<8\}$，}
\step{所以 $A\cup C=\{x\mid x\leqslant-3\text{ 或 }x>-2\}$；}
\stepm{(2)}{全集 $U=R$，集合 $A=\{x\mid x^2-9\geqslant0\}=\{x\mid x\geqslant3\text{ 或 }x\leqslant-3\}$，}
\step{$B=\{x\mid|x-4|<2\}=\{x\mid2<x<6\}$，}
\step{所以 $\overline A=\{x\mid-3<x<3\}$，所以 $\overline A\cap B=\{x\mid2<x<3\}$.}
\solend
\fi

\ansspace[6]

\item 已知集合 $A=\left\{x\left|\dfrac{2x-1}{x+1}<1,x\in\R\right.\right\}$，集合 $B=\{x\mid(x-a-1)(x-a+1)\leqslant0\}$.

（1）求集合 $A$；

（2）若 $B\subseteq A$，求实数 $a$ 的取值范围.

\ifanswer
\solbegin
\stepm{(1)}{由 $\dfrac{2x-1}{x+1}<1$，得 $\dfrac{x-2}{x+1}<0$ 所以 $(x-2)(x+1)<0$，解得 $-1<x<2$，}
\step{即 $A=(-1,2)$；}
\stepm{(2)}{由 $(x-a-1)(x-a+1)\leqslant0$，得 $a-1\leqslant x\leqslant a+1$，}
\step{所以 $B=\{x\mid a-1\leqslant x\leqslant a+1\}$，}
\step{又 $B\subseteq A$，且 $a+1>a-1$ 即 $B\neq\varnothing$，所以 $\begin{cases}a-1>-1\\ a+1<2\end{cases}$，所以 $0<a<1$，}
\step{所以实数 $a$ 的取值范围是 $\{a\mid0<a<1\}$.}
\solend
\fi

\ansspace[6]

\item 已知集合 $A=\{x\mid x^2-4x+3=0\}$，$B=\{x\mid x^2-ax+3=0\}$.

（1）若 $A\cup B=B$，求实数 $a$ 的值；

（2）若 $A\cap B=B$，求实数 $a$ 的取值范围.

\ifanswer
\solbegin
\stepm{(1)}{集合 $A=\{x\mid x^2-4x+3=0\}=\{1,3\}$，$B=\{x\mid x^2-ax+3=0\}$.}
\step{$A\cup B=B$，所以 $1\in B$，$3\in B$，所以 $\begin{cases}1-a+3=0\\ 9-3a+3=0\end{cases}$，解得实数 $a=4$；}
\stepm{(2)}{因为 $A\cap B=B$，所以 $B\subseteq A$，}
\step{当 $B=\varnothing$ 时，$\Delta=(-a)^2-12<0$，解得 $-2\sqrt3<a<2\sqrt3$，}
\step{当 $B$ 是单元素集时，$\Delta=a^2-12=0$，解得 $a=2\sqrt3$，此时 $B=\{\sqrt3\}$，}
\step{不合题意；}
\step{当 $B=\{1,3\}$ 时，由 (1) 得 $a=4$，}
\step{综上，实数 $a$ 的取值范围是 $(-2\sqrt3,2\sqrt3)\cup\{4\}$.}
\solend
\fi

\ansspace[6]

\item 已知全集为 $R$，实数 $a$、$b$ 满足 $a\neq b$，集合 $M=\{x\mid x^2-3x-4\leqslant0\}$，集合 $A=\{x\mid(x+a)(x+b)>0\}$，集合 $B=\{x\mid x^2+(a-2)x-2a>0\}$.

（1）若 $\overline A=M$，求 $a$、$b$ 的值；

（2）若 $a>b>-1$，求 $A\cap B$；

（3）若 $a^2+a\in\overline B$，求 $a$ 的取值范围.

\ifanswer
\solbegin
\stepm{(1)}{$M=\{x\mid x^2-3x-4\leqslant0\}=\{x\mid-1\leqslant x\leqslant4\}$，$A=\{x\mid(x+a)(x+b)\leqslant0\}$，}
\step{因为 $\overline A=M$，所以 $a=1$，$b=-4$ 或 $a=-4$，$b=1$.}
\stepm{(2)}{因为 $a>b>-1$，所以 $-a<-b<1$，所以 $A=\{x\mid x<-a\text{ 或 }x>-b\}$，}
\step{因为 $B=\{x\mid x^2+(a-2)x-2a>0\}=\{x\mid x<-a\text{ 或 }x>2\}$，}
\step{所以 $A\cap B=\{x\mid x<-a\text{ 或 }x>2\}$，}
\stepm{(3)}{$\overline B=\{x\mid(x-2)(x+a)\leqslant0\}$，}
\step{由 $a^2+a\in\overline B$，得 $(a^2+a-2)(a^2+2a)\leqslant0\Leftrightarrow a(a-1)(a-2)^2\leqslant0$，}
\step{解得 $a=2$ 或 $0\leqslant a\leqslant1$，故 $a$ 的取值范围是 $[0,1]\cup\{2\}$.}
\solend
\fi

\ansspace[8]

\item 对平面直角坐标系的第一象限内的任意两点作如下定义：若 $\dfrac ab>\dfrac cd$，那么称点 $(a,b)$ 是点 $(c,d)$ 的“上位点”，同时点 $(c,d)$ 是点 $(a,b)$ 的“下位点”.

（1）试写出点 $(2,4)$ 的一个“上位点”坐标；

（2）已知点 $(a,b)$ 是点 $(c,d)$ 的“上位点”，判断点 $P(a+c,b+d)$ 是点 $(c,d)$ 的“上位点”还是“下位点”，并证明你的结论；

（3）设正整数 $n$ 满足以下条件：对集合 $\{t\mid0<t<2020,t\in\Z\}$ 内的任意元素 $m$，总存在正整数 $k=2m+1$，使得点 $(n,k)$ 既是点 $(2025,m)$ 的“下位点”，又是点 $(2026,m+1)$ 的“上位点”，问这样的正整数 $n$ 是否存在？若存在，求正整数 $n$ 的值；若不存在，请说明理由.

\ifanswer
\solbegin
\stepm{(1)}{由定义，点 $(a,b)$ 是点 $(2,4)$ 的“上位点”，}
\step{当且仅当 $\dfrac ab>\dfrac24=\dfrac12$ 且 $a>0,b>0$.}
\step{取 $(a,b)=(1,1)$，因为 $\dfrac11>\dfrac12$，}
\step{所以点 $(1,1)$ 是点 $(2,4)$ 的一个“上位点”.（答案不唯一）}
\stepm{(2)}{点 $P(a+c,b+d)$ 是点 $(c,d)$ 的“上位点”.}
\step{$(a,b)$ 是点 $(c,d)$ 的“上位点”，所以 $\dfrac ab>\dfrac cd$.}
\step{因为 $a,b,c,d>0$，两边同乘 $bd$ 得 $ad>bc$，即 $ad-bc>0$.}
\step{作差比较 $\dfrac{a+c}{b+d}$ 与 $\dfrac cd$：$\dfrac{a+c}{b+d}-\dfrac cd=\dfrac{d(a+c)-c(b+d)}{d(b+d)}=\dfrac{ad-bc}{d(b+d)}>0$，}
\step{所以点 $P(a+c,b+d)$ 是点 $(c,d)$ 的“上位点”.}
\stepm{(3)}{由定义，点 $(n,k)$ 是点 $(2025,m)$ 的下位点即 $\dfrac{2025}m>\dfrac nk$；}
\step{点 $(n,k)$ 是点 $(2026,m+1)$ 的上位点即 $\dfrac nk>\dfrac{2026}{m+1}$.}
\step{代入 $k=2m+1$，得对任意 $m\in\{1,2,\cdots,2019\}$，}
\step{都有 $4052-\dfrac{2026}{m+1}<n<4050+\dfrac{2025}m$.}
\step{右边 $4050+\dfrac{2025}m$ 随 $m$ 增大而减小，故最近处在 $m=2019$，}
\step{得 $n<4050+\dfrac{2025}{2019}<4051.01$，所以 $n\leqslant4051$.}
\step{左边 $4052-\dfrac{2026}{m+1}$ 随 $m$ 增大而增大，故最近处也在 $m=2019$，}
\step{得 $n>4052-\dfrac{2026}{2020}>4050.99$，所以 $n\geqslant4051$.}
\step{综上，$n=4051$.}
\step{检验：当 $n=4051$ 时，对任意 $1\leqslant m\leqslant2019,k=2m+1$ 是正整数，}
\step{且 $\dfrac{4051}{2m+1}<\dfrac{2025}m$,$4051m<4050m+2025,m<2025$，成立；}
\step{$\dfrac{4051}{2m+1}>\dfrac{2026}{m+1}$,$4051(m+1)>4052m+2026,m<2025$，成立.}
\step{又各点横纵坐标均为正，都在第一象限内，符合定义，}
\step{所以满足条件的正整数 $n$ 存在，其值为 $n=4051$.}
\solend
\fi

\ansspace*[10]

\end{enumerate}

\end{document}
"
%
% 原始材料：微信公众号“魔都数学加油站”文章，作者破晓，连载“27学年高一59”；
% 原文标题“27学年高一59——2026-2027年上海市华实高中高一上10月月考”；
% 原文链接：https://mp.weixin.qq.com/s/XzyBCmrvpRuMe-NuNk7gkA。
% 页面标注发布时间 2026-10-13 21:45（上海）。
% 正文为 9 张 PNG 页：第 01、02、03、05、06、07、08、09 页 4762×6735，第 04 页 2560×3620
% （**同一篇各页分辨率不同**），均按页序读取。9 页全为试卷正文页（卷首标题在第 1 页页顶，
% 卷末解答题第 21 题解析止于第 9 页），无封面页、目录页、二维码或宣传页。
% 粉色斜排“魔都数学加油站”为卷面水印，与公众号页脚一并未录入。
%
% 卷首标题：照卷面印作“2026--2027 年华实高中高一上 10 月月考”（公众号标题中的“上海市”
% 卷面标题行没有，本稿照卷面），年份区间按全稿惯例排 en dash。原卷只有一行标题、无副标题。
%
% 篇幅与题号：原卷自第 3 题起编号（第 1、2 题不在本卷内）。填空题 3—12（10 题）、
% 选择题 13—16（4 题）、解答题 17—21（5 题），共 19 题。本稿沿用原节与原题号，
% 只改排为全稿统一的大节样式（原卷小节标题为“一、填空题”“二、选择题”“三、解答题”）。
%
% 讲评版：原卷每题下直接印【答案】或【解析】。**只印【答案】的 1 题**（第 3 题），本稿只给 \ans{}；
% 其余 18 题都印【解析】，按全稿口径这类题不另提 \ans{}——第 13、14、15、16 题的结论写在解析末句
% （“故选 D”等），亦照录在解析内。本卷**没有**印在题干括号内的答案。
%
% 副标题：原卷未印副标题。本稿按“知识模块”口径标作“集合与常用逻辑用语、不等式”：
% 19 题中集合与常用逻辑用语 13 题（第 3、9、10、11、12、14、15、16、17、18、19、20、21 题），
% 不等式 10 题（含等式与不等式性质、一元二次不等式、恒成立、基本应用：第 4、5、6、7、8、
% 13、17、18、19、20 题），两类均远超三分之一，故并标。口径同上。
%
% 与原卷的排版差异：
%   ① 原卷每题下直接印【答案】/【解析】，本稿按全稿惯例将答案与解析置于答案版，试卷版保持干净；
%   ② 原卷未印副标题，本稿按“知识模块”口径补出；
%   ③ 原卷填空横线约两个字宽，本稿按全稿惯例统一排 \blankline[4.4em]；
%   ④ 原卷实数集排粗体 $\mathbf{R}$（第 12、13、20 题等），整数集排粗体 $\mathbf{Z}$（第 11、12 题），
%      自然数集排斜体 $N$ 或粗体 $\mathbf{Z}$ 处按本稿惯例改用 $\R$、$\Z$、$\N$；
%      空集记号统一为 \varnothing；
%   ⑤ 原卷在数学式内多处用半角逗号（如 $x_1,x_2$、$a,b\in R$），均照录未强改；
%      句末句点按全稿惯例排西文句点；“漂亮集”“上位点”“下位点”等自定义名词照录；
%   ⑥ 第 17—21 题的子问原卷各自另起一行，本稿按全稿惯例用 enumerate 的 (1)(2)(3)；
%      解析中的子问标号一律用 \stepm{(1)}{…}；第 9、11 题解析中的①、②（③④）分情形或分结论，
%      按全稿惯例作 \stepm{①}{…}、\stepm{②}{…}（第 11 题解析末“故正确的是②④”照录为行内）；
%   ⑦ 第 12 题题干里的“$A_p$”下标原卷作斜体 $p$，照录；
%   ⑧ 每道解答题原卷未留作答空白，本稿按全稿惯例加 \ansspace。
%
% 原卷笔误/存疑（均照录未改，正文对应处另有 % 注）：本卷首轮逐页转录，暂未发现影响答案的笔误。
%
% 插图：无（全卷 9 页无任何图形、无“如图”字样）。
%
\documentclass[a4paper,11pt]{article}
\usepackage{common}

\begin{document}

\examtitlest[3]{2026--2027 年华实高中高一上 10 月月考}{集合与常用逻辑用语、不等式}

\bigsec{一、填空题}

\begin{enumerate}
\setcounter{enumi}{2}

\item “$x\geqslant1$ 且 $y\geqslant1$”的否定形式为\blankline[4.4em].

\ifanswer
\ans{$x<1$ 或 $y<1$}
\fi

\item 已知方程 $2x^2+4x-3=0$ 的两个根为 $x_1,x_2$，则 $\dfrac1{x_1}+\dfrac1{x_2}=$\blankline[4.4em].

\ifanswer
\solbegin
\step{利用韦达定理得 $\begin{cases}x_1+x_2=-2\\ x_1\cdot x_2=-\dfrac32\end{cases}$，则 $\dfrac1{x_1}+\dfrac1{x_2}=\dfrac{x_1+x_2}{x_1\cdot x_2}=\dfrac43$.}
\solend
\fi

\item 设 $a>b>0$，比较大小：$\dfrac1{a-b}$\blankline[2.4em]$\dfrac1a$.

\ifanswer
\solbegin
\step{由 $a>b>0$ 得 $a>0$ 且 $a-b>0$，所以 $\dfrac1{a-b}-\dfrac1a=\dfrac b{a(a-b)}>0$，}
\step{故 $\dfrac1{a-b}>\dfrac1a$.}
\solend
\fi

\item 若不等式 $2ax^2-ax+1>0$ 的解集为 $R$，则实数 $a$ 的取值范围是\blankline[4.4em].

\ifanswer
\solbegin
\step{当 $a\neq0$ 时，由题意得 $\begin{cases}a>0\\ \Delta<0\end{cases}$，即 $\begin{cases}a>0\\ a^2-8a<0\end{cases}$，解得 $0<a<8$.}
\step{当 $a=0$ 时，恒有 $1>0$，不等式也成立.}
\step{综上所述，$a$ 的取值范围是 $[0,8)$.}
\solend
\fi

\item 集合 $A=\{x\mid m-1\leqslant x\leqslant m+7\}$，集合 $B=\{x\mid x^2-3x-10\geqslant0\}$，若 $A\cup B=R$，则实数 $m$ 的取值范围为\blankline[4.4em].

\ifanswer
\solbegin
\step{$A=\{x\mid m-1\leqslant x\leqslant m+7\}$，$B=\{x\mid x^2-3x-10\geqslant0\}=\{x\mid x\leqslant-2\text{ 或 }x\geqslant5\}$，}
\step{若 $A\cup B=R$，则 $\begin{cases}m-1\leqslant-2\\ m+7\geqslant5\end{cases}$，解得 $-2\leqslant m\leqslant-1$，}
\step{所以实数 $m$ 的取值范围为 $[-2,-1]$.}
\solend
\fi

\item 若关于 $x$ 的不等式 $ax^2+bx+c>0$ 的解集为 $(-1,3)$，则关于 $x$ 的不等式 $cx^2+bx+a>0$ 的解集为\blankline[4.4em].

\ifanswer
\solbegin
\step{因为不等式 $ax^2+bx+c>0$ 的解集为 $(-1,3)$，}
\step{所以 $-1$ 和 $3$ 是方程 $ax^2+bx+c=0$ 的实数解，}
\step{所以 $\begin{cases}-1+3=-\dfrac ba\\ -1\times3=\dfrac ca\\ a<0\end{cases}$，解得 $b=-2a$，$c=-3a$，且 $a<0$，}
\step{所以不等式 $cx^2+bx+a>0$ 可化为 $-3ax^2-2ax+a>0$，即 $3x^2+2x-1>0$，}
\step{解得 $x<-1$ 或 $x>\dfrac13$，所以该不等式的解集为 $(-\infty,-1)\cup\left(\dfrac13,+\infty\right)$.}
\solend
\fi

\item 若集合 $A=\{x\mid(a-1)x^2-2x+1=0\}$ 有且仅有 $2$ 个子集，则实数 $a$ 的值为\blankline[4.4em].

\ifanswer
\solbegin
\step{由题意得方程 $(a-1)x^2-2x+1=0$ 只有一个根，}
\stepm{①}{当 $a-1=0$ 时，$a=1$，此时方程化为 $-2x+1=0$，解得 $x=\dfrac12$，符合题意，}
\stepm{②}{当 $a-1\neq0$，即 $a\neq1$ 时，$\Delta=4-4(a-1)=0$，解得 $a=2$，}
\step{综上所述，实数 $a$ 的值为 $1$ 或 $2$.}
\solend
\fi

\item “$x+y\neq5$”是“$x\neq3$ 或 $y\neq2$”的\blankline[4.4em]条件.

\ifanswer
\solbegin
\step{分析充分性：}
\step{若 $x+y\neq5$ 成立，若 $x=3$，则 $y\neq2$，若 $y=2$，则 $x\neq3$，}
\step{所以一定能说明 $x\neq3$ 和 $y\neq2$ 至少有一个成立，即 $x\neq3$ 或 $y\neq2$ 成立，}
\step{充分性成立；}
\step{分析必要性：}
\step{若 $x\neq3$ 或 $y\neq2$，即 $x\neq3$ 和 $y\neq2$ 二者至少一个成立，}
\step{举例：$x=1,y=4$，此时 $x+y=5$，所以必要性不成立.}
\step{所以“$x+y\neq5$”是“$x\neq3$ 或 $y\neq2$”的充分不必要条件.}
\solend
\fi

\item 设 $P$ 为非空实数集满足：对任意给定的 $x,y\in P$（$x,y$ 可以相同），都有 $x+y\in P,x-y\in P,xy\in P$，则称 $P$ 为漂亮集. 关于漂亮集有下列结论：

①集合 $P=\{-2,-1,0,1,2\}$ 为漂亮集；②集合 $P=\{x\mid x=2n,n\in\Z\}$ 为漂亮集；

③若集合 $P_1,P_2$ 为漂亮集，则 $P_1\cup P_2$ 为漂亮集；④若集合 $P$ 为漂亮集，则一定有 $0\in P$.

其中所有正确结论的序号是\blankline[4.4em].

\ifanswer
\solbegin
\step{对于集合 $P=\{-2,-1,0,1,2\}$，取 $x=2,y=-2$，则 $x-y=2-(-2)=4\notin P$，}
\step{不满足漂亮集的定义，所以①错误.}
\step{对于集合 $P=\{x\mid x=2n,n\in Z\}$，设 $x=2n_1,y=2n_2,n_1,n_2\in Z$.}
\step{$x+y=2n_1+2n_2=2(n_1+n_2)$，因为 $n_1+n_2\in Z$，所以 $x+y\in P$.}
\step{$x-y=2n_1-2n_2=2(n_1-n_2)$，因为 $n_1-n_2\in Z$，所以 $x-y\in P$.}
\step{$xy=2n_1\times2n_2=4n_1n_2=2(2n_1n_2)$，因为 $2n_1n_2\in Z$，所以 $xy\in P$.}
\step{满足漂亮集的定义，所以②正确.}
\step{令 $P_1=\{x\mid x=2n,n\in Z\}$,$P_2=\{x\mid x=3n,n\in Z\}$，都是漂亮集.}
\step{$P_1\cup P_2$ 中，取 $x=2(x\in P_1)$,$y=3(y\in P_2)$,$x+y=5\notin P_1\cup P_2$，}
\step{不满足漂亮集的定义，所以③错误.}
\step{因为对任意 $x\in P$，有 $x-x=0\in P$，所以若集合 $P$ 为漂亮集，}
\step{则一定有 $0\in P$，所以④正确.}
\step{故正确的是②④.}
\solend
\fi

\item 已知：集合 $A_p=\{(x,y)\mid y=3x^2-px+5p-65,x\in\R\}$,$M=\{x\mid1\leqslant x\leqslant100,x\in\Z\}$，若 $p\in M$，则所有 $A_p$ 的交集\blankline[4.4em].

\ifanswer
\solbegin
\step{由 $M=\{1,2,\cdots,100\}$，所以所有 $A_p$ 的交集就是这些集合中共同拥有的点.}
\step{若 $(x,y)$ 是所有集合的一个公共点，}
\step{那么它特别要同时属于 $p=1$ 和 $p=2$ 时对应的两个集合.}
\step{当 $p=1$ 时，$y=3x^2-x-60$；}
\step{当 $p=2$ 时，$y=3x^2-2x-55$；联立得 $x=5,y=10$.}
\step{因此，如果公共点存在，只可能是 $(5,10)$.}
\step{接下来验证这个点是否属于所有集合，}
\step{对任意 $p\in M$，把 $x=5$ 代入集合的关系式，}
\step{得 $3x^2-px+5p-65=75-5p+5p-65=10$，与 $y=10$ 相符，}
\step{所以 $(5,10)\in A_p$ 对每一个 $p$ 都成立，因此所有 $A_p$ 的交集为 $\{(5,10)\}$.}
\solend
\fi

\end{enumerate}

% 守卫须放在 \bigsec 之前，否则标题仍会落页末（见 高一40 的教训）
\needvspace{5cm}

\bigsec{二、选择题}

\begin{enumerate}
\setcounter{enumi}{12}

\item 已知 $a,b,c\in\R$ 且 $a>b$，则下列不等式一定成立的是\mbox{（\quad）}

\keepwithprev
A. $\dfrac1a<\dfrac1b$\qquad B. $a^2>b^2$

C. $a|c|>b|c|$\qquad D. $\dfrac a{c^2+1}>\dfrac b{c^2+1}$

\ifanswer
\solbegin
\step{$A$：当 $a>0>b$ 时，$\dfrac1a>\dfrac1b$，故 $A$ 错误；}
\step{$B$：当 $a=-1,b=-2$ 时，满足 $a>b$，但 $a^2>b^2$ 不成立，故 $B$ 错误；}
\step{$C$：当 $c=0$ 时，$a|c|=b|c|$，故 $C$ 错误；}
\step{$D$：由 $a>b,c^2+1>0$，得 $\dfrac a{c^2+1}>\dfrac b{c^2+1}$，故 $D$ 正确；}
\step{故选 $D$.}
\solend
\fi

\item 已知集合 $A=\{a^2,0,-1\}$，$B=\{a,b,0\}$，若 $A=B$，则 $(ab)^{2026}$ 的值为\mbox{（\quad）}

\keepwithprev
A. $0$\qquad B. $-1$

C. $1$\qquad D. $\pm1$

\ifanswer
\solbegin
\step{因为集合 $A=\{a^2,0,-1\}$，$B=\{a,b,0\}$，$A=B$，所以 $\begin{cases}a=-1\\ a^2=b\end{cases}$，}
\step{所以 $a=-1$，$b=1$，所以 $(ab)^{2026}=1$，故选 $C$.}
\solend
\fi

\item 设全集 $U=\R$，$A=\{x\mid x<-1\text{ 或 }x>2\}$,$B=\{y\mid y=x^2+a\}$，若 $\overline A\subseteq B$，则实数 $a$ 的取值范围是\mbox{（\quad）}

\keepwithprev
A. $a\leqslant-1$\qquad B. $a>2$

C. $a<-1$\qquad D. $a\geqslant-2$

\ifanswer
\solbegin
\step{因为 $A=\{x\mid x<-1\text{ 或 }x>2\}$，则 $\overline A=\{x\mid-1\leqslant x\leqslant2\}$，}
\step{又 $B=\{y\mid y=x^2+a\}=\{y\mid y\geqslant a\}$，且 $\overline A\subseteq B$，所以 $a\leqslant-1$，故选 $A$.}
\solend
\fi

\item 设 $P$ 和 $Q$ 是两个集合，定义集合 $P-Q=\{x\mid x\in P$，且 $x\notin Q\}$，如果 $Q=\{x\mid(1-x)(x-3)>0\}$，$P=\left\{x\left|\dfrac x{2-x}>0\right.\right\}$，那么 $P-Q$ 等于\mbox{（\quad）}

\keepwithprev
A. $\{x\mid0<x<1\}$\qquad B. $\{x\mid1\leqslant x<2\}$

C. $\{x\mid0<x\leqslant1\}$\qquad D. $\{x\mid2\leqslant x<3\}$

\ifanswer
\solbegin
\step{因为 $P=\left\{x\left|\dfrac x{2-x}>0\right.\right\}=\{x\mid x(x-2)<0\}=\{x\mid0<x<2\}$，}
\step{$Q=\{x\mid(1-x)(x-3)>0\}=\{x\mid(x-3)(x-1)<0\}=\{x\mid1<x<3\}$，}
\step{所以 $P-Q=\{x\mid0<x\leqslant1\}$，故选 $C$.}
\solend
\fi

\end{enumerate}

\bigsec{三、解答题}

\begin{enumerate}
\setcounter{enumi}{16}

\item 已知集合 $A=\{x\mid x^2-9\geqslant0\}$，$B=\{x\mid|x-4|<2\}$，$C=\left\{x\left|\dfrac{x-8}{x+2}<0\right.\right\}$.

（1）求 $A\cup C$；

（2）若全集 $U=R$，求 $\overline A\cap B$.

\ifanswer
\solbegin
\stepm{(1)}{集合 $A=\{x\mid x^2-9\geqslant0\}=\{x\mid x\geqslant3\text{ 或 }x\leqslant-3\}$，}
\step{$B=\{x\mid|x-4|<2\}=\{x\mid2<x<6\}$，$C=\left\{x\left|\dfrac{x-8}{x+2}<0\right.\right\}=\{x\mid-2<x<8\}$，}
\step{所以 $A\cup C=\{x\mid x\leqslant-3\text{ 或 }x>-2\}$；}
\stepm{(2)}{全集 $U=R$，集合 $A=\{x\mid x^2-9\geqslant0\}=\{x\mid x\geqslant3\text{ 或 }x\leqslant-3\}$，}
\step{$B=\{x\mid|x-4|<2\}=\{x\mid2<x<6\}$，}
\step{所以 $\overline A=\{x\mid-3<x<3\}$，所以 $\overline A\cap B=\{x\mid2<x<3\}$.}
\solend
\fi

\ansspace[6]

\item 已知集合 $A=\left\{x\left|\dfrac{2x-1}{x+1}<1,x\in\R\right.\right\}$，集合 $B=\{x\mid(x-a-1)(x-a+1)\leqslant0\}$.

（1）求集合 $A$；

（2）若 $B\subseteq A$，求实数 $a$ 的取值范围.

\ifanswer
\solbegin
\stepm{(1)}{由 $\dfrac{2x-1}{x+1}<1$，得 $\dfrac{x-2}{x+1}<0$ 所以 $(x-2)(x+1)<0$，解得 $-1<x<2$，}
\step{即 $A=(-1,2)$；}
\stepm{(2)}{由 $(x-a-1)(x-a+1)\leqslant0$，得 $a-1\leqslant x\leqslant a+1$，}
\step{所以 $B=\{x\mid a-1\leqslant x\leqslant a+1\}$，}
\step{又 $B\subseteq A$，且 $a+1>a-1$ 即 $B\neq\varnothing$，所以 $\begin{cases}a-1>-1\\ a+1<2\end{cases}$，所以 $0<a<1$，}
\step{所以实数 $a$ 的取值范围是 $\{a\mid0<a<1\}$.}
\solend
\fi

\ansspace[6]

\item 已知集合 $A=\{x\mid x^2-4x+3=0\}$，$B=\{x\mid x^2-ax+3=0\}$.

（1）若 $A\cup B=B$，求实数 $a$ 的值；

（2）若 $A\cap B=B$，求实数 $a$ 的取值范围.

\ifanswer
\solbegin
\stepm{(1)}{集合 $A=\{x\mid x^2-4x+3=0\}=\{1,3\}$，$B=\{x\mid x^2-ax+3=0\}$.}
\step{$A\cup B=B$，所以 $1\in B$，$3\in B$，所以 $\begin{cases}1-a+3=0\\ 9-3a+3=0\end{cases}$，解得实数 $a=4$；}
\stepm{(2)}{因为 $A\cap B=B$，所以 $B\subseteq A$，}
\step{当 $B=\varnothing$ 时，$\Delta=(-a)^2-12<0$，解得 $-2\sqrt3<a<2\sqrt3$，}
\step{当 $B$ 是单元素集时，$\Delta=a^2-12=0$，解得 $a=2\sqrt3$，此时 $B=\{\sqrt3\}$，}
\step{不合题意；}
\step{当 $B=\{1,3\}$ 时，由 (1) 得 $a=4$，}
\step{综上，实数 $a$ 的取值范围是 $(-2\sqrt3,2\sqrt3)\cup\{4\}$.}
\solend
\fi

\ansspace[6]

\item 已知全集为 $R$，实数 $a$、$b$ 满足 $a\neq b$，集合 $M=\{x\mid x^2-3x-4\leqslant0\}$，集合 $A=\{x\mid(x+a)(x+b)>0\}$，集合 $B=\{x\mid x^2+(a-2)x-2a>0\}$.

（1）若 $\overline A=M$，求 $a$、$b$ 的值；

（2）若 $a>b>-1$，求 $A\cap B$；

（3）若 $a^2+a\in\overline B$，求 $a$ 的取值范围.

\ifanswer
\solbegin
\stepm{(1)}{$M=\{x\mid x^2-3x-4\leqslant0\}=\{x\mid-1\leqslant x\leqslant4\}$，$A=\{x\mid(x+a)(x+b)\leqslant0\}$，}
\step{因为 $\overline A=M$，所以 $a=1$，$b=-4$ 或 $a=-4$，$b=1$.}
\stepm{(2)}{因为 $a>b>-1$，所以 $-a<-b<1$，所以 $A=\{x\mid x<-a\text{ 或 }x>-b\}$，}
\step{因为 $B=\{x\mid x^2+(a-2)x-2a>0\}=\{x\mid x<-a\text{ 或 }x>2\}$，}
\step{所以 $A\cap B=\{x\mid x<-a\text{ 或 }x>2\}$，}
\stepm{(3)}{$\overline B=\{x\mid(x-2)(x+a)\leqslant0\}$，}
\step{由 $a^2+a\in\overline B$，得 $(a^2+a-2)(a^2+2a)\leqslant0\Leftrightarrow a(a-1)(a-2)^2\leqslant0$，}
\step{解得 $a=2$ 或 $0\leqslant a\leqslant1$，故 $a$ 的取值范围是 $[0,1]\cup\{2\}$.}
\solend
\fi

\ansspace[8]

\item 对平面直角坐标系的第一象限内的任意两点作如下定义：若 $\dfrac ab>\dfrac cd$，那么称点 $(a,b)$ 是点 $(c,d)$ 的“上位点”，同时点 $(c,d)$ 是点 $(a,b)$ 的“下位点”.

（1）试写出点 $(2,4)$ 的一个“上位点”坐标；

（2）已知点 $(a,b)$ 是点 $(c,d)$ 的“上位点”，判断点 $P(a+c,b+d)$ 是点 $(c,d)$ 的“上位点”还是“下位点”，并证明你的结论；

（3）设正整数 $n$ 满足以下条件：对集合 $\{t\mid0<t<2020,t\in\Z\}$ 内的任意元素 $m$，总存在正整数 $k=2m+1$，使得点 $(n,k)$ 既是点 $(2025,m)$ 的“下位点”，又是点 $(2026,m+1)$ 的“上位点”，问这样的正整数 $n$ 是否存在？若存在，求正整数 $n$ 的值；若不存在，请说明理由.

\ifanswer
\solbegin
\stepm{(1)}{由定义，点 $(a,b)$ 是点 $(2,4)$ 的“上位点”，}
\step{当且仅当 $\dfrac ab>\dfrac24=\dfrac12$ 且 $a>0,b>0$.}
\step{取 $(a,b)=(1,1)$，因为 $\dfrac11>\dfrac12$，}
\step{所以点 $(1,1)$ 是点 $(2,4)$ 的一个“上位点”.（答案不唯一）}
\stepm{(2)}{点 $P(a+c,b+d)$ 是点 $(c,d)$ 的“上位点”.}
\step{$(a,b)$ 是点 $(c,d)$ 的“上位点”，所以 $\dfrac ab>\dfrac cd$.}
\step{因为 $a,b,c,d>0$，两边同乘 $bd$ 得 $ad>bc$，即 $ad-bc>0$.}
\step{作差比较 $\dfrac{a+c}{b+d}$ 与 $\dfrac cd$：$\dfrac{a+c}{b+d}-\dfrac cd=\dfrac{d(a+c)-c(b+d)}{d(b+d)}=\dfrac{ad-bc}{d(b+d)}>0$，}
\step{所以点 $P(a+c,b+d)$ 是点 $(c,d)$ 的“上位点”.}
\stepm{(3)}{由定义，点 $(n,k)$ 是点 $(2025,m)$ 的下位点即 $\dfrac{2025}m>\dfrac nk$；}
\step{点 $(n,k)$ 是点 $(2026,m+1)$ 的上位点即 $\dfrac nk>\dfrac{2026}{m+1}$.}
\step{代入 $k=2m+1$，得对任意 $m\in\{1,2,\cdots,2019\}$，}
\step{都有 $4052-\dfrac{2026}{m+1}<n<4050+\dfrac{2025}m$.}
\step{右边 $4050+\dfrac{2025}m$ 随 $m$ 增大而减小，故最近处在 $m=2019$，}
\step{得 $n<4050+\dfrac{2025}{2019}<4051.01$，所以 $n\leqslant4051$.}
\step{左边 $4052-\dfrac{2026}{m+1}$ 随 $m$ 增大而增大，故最近处也在 $m=2019$，}
\step{得 $n>4052-\dfrac{2026}{2020}>4050.99$，所以 $n\geqslant4051$.}
\step{综上，$n=4051$.}
\step{检验：当 $n=4051$ 时，对任意 $1\leqslant m\leqslant2019,k=2m+1$ 是正整数，}
\step{且 $\dfrac{4051}{2m+1}<\dfrac{2025}m$,$4051m<4050m+2025,m<2025$，成立；}
\step{$\dfrac{4051}{2m+1}>\dfrac{2026}{m+1}$,$4051(m+1)>4052m+2026,m<2025$，成立.}
\step{又各点横纵坐标均为正，都在第一象限内，符合定义，}
\step{所以满足条件的正整数 $n$ 存在，其值为 $n=4051$.}
\solend
\fi

\ansspace*[10]

\end{enumerate}

\end{document}
"
%
% 原始材料：微信公众号“魔都数学加油站”文章，作者破晓，连载“27学年高一59”；
% 原文标题“27学年高一59——2026-2027年上海市华实高中高一上10月月考”；
% 原文链接：https://mp.weixin.qq.com/s/XzyBCmrvpRuMe-NuNk7gkA。
% 页面标注发布时间 2026-10-13 21:45（上海）。
% 正文为 9 张 PNG 页：第 01、02、03、05、06、07、08、09 页 4762×6735，第 04 页 2560×3620
% （**同一篇各页分辨率不同**），均按页序读取。9 页全为试卷正文页（卷首标题在第 1 页页顶，
% 卷末解答题第 21 题解析止于第 9 页），无封面页、目录页、二维码或宣传页。
% 粉色斜排“魔都数学加油站”为卷面水印，与公众号页脚一并未录入。
%
% 卷首标题：照卷面印作“2026--2027 年华实高中高一上 10 月月考”（公众号标题中的“上海市”
% 卷面标题行没有，本稿照卷面），年份区间按全稿惯例排 en dash。原卷只有一行标题、无副标题。
%
% 篇幅与题号：原卷自第 3 题起编号（第 1、2 题不在本卷内）。填空题 3—12（10 题）、
% 选择题 13—16（4 题）、解答题 17—21（5 题），共 19 题。本稿沿用原节与原题号，
% 只改排为全稿统一的大节样式（原卷小节标题为“一、填空题”“二、选择题”“三、解答题”）。
%
% 讲评版：原卷每题下直接印【答案】或【解析】。**只印【答案】的 1 题**（第 3 题），本稿只给 \ans{}；
% 其余 18 题都印【解析】，按全稿口径这类题不另提 \ans{}——第 13、14、15、16 题的结论写在解析末句
% （“故选 D”等），亦照录在解析内。本卷**没有**印在题干括号内的答案。
%
% 副标题：原卷未印副标题。本稿按“知识模块”口径标作“集合与常用逻辑用语、不等式”：
% 19 题中集合与常用逻辑用语 13 题（第 3、9、10、11、12、14、15、16、17、18、19、20、21 题），
% 不等式 10 题（含等式与不等式性质、一元二次不等式、恒成立、基本应用：第 4、5、6、7、8、
% 13、17、18、19、20 题），两类均远超三分之一，故并标。口径同上。
%
% 与原卷的排版差异：
%   ① 原卷每题下直接印【答案】/【解析】，本稿按全稿惯例将答案与解析置于答案版，试卷版保持干净；
%   ② 原卷未印副标题，本稿按“知识模块”口径补出；
%   ③ 原卷填空横线约两个字宽，本稿按全稿惯例统一排 \blankline[4.4em]；
%   ④ 原卷实数集排粗体 $\mathbf{R}$（第 12、13、20 题等），整数集排粗体 $\mathbf{Z}$（第 11、12 题），
%      自然数集排斜体 $N$ 或粗体 $\mathbf{Z}$ 处按本稿惯例改用 $\R$、$\Z$、$\N$；
%      空集记号统一为 \varnothing；
%   ⑤ 原卷在数学式内多处用半角逗号（如 $x_1,x_2$、$a,b\in R$），均照录未强改；
%      句末句点按全稿惯例排西文句点；“漂亮集”“上位点”“下位点”等自定义名词照录；
%   ⑥ 第 17—21 题的子问原卷各自另起一行，本稿按全稿惯例用 enumerate 的 (1)(2)(3)；
%      解析中的子问标号一律用 \stepm{(1)}{…}；第 9、11 题解析中的①、②（③④）分情形或分结论，
%      按全稿惯例作 \stepm{①}{…}、\stepm{②}{…}（第 11 题解析末“故正确的是②④”照录为行内）；
%   ⑦ 第 12 题题干里的“$A_p$”下标原卷作斜体 $p$，照录；
%   ⑧ 每道解答题原卷未留作答空白，本稿按全稿惯例加 \ansspace。
%
% 原卷笔误/存疑（均照录未改，正文对应处另有 % 注）：本卷首轮逐页转录，暂未发现影响答案的笔误。
%
% 插图：无（全卷 9 页无任何图形、无“如图”字样）。
%
\documentclass[a4paper,11pt]{article}
\usepackage{common}

\begin{document}

\examtitlest[3]{2026--2027 年华实高中高一上 10 月月考}{集合与常用逻辑用语、不等式}

\bigsec{一、填空题}

\begin{enumerate}
\setcounter{enumi}{2}

\item “$x\geqslant1$ 且 $y\geqslant1$”的否定形式为\blankline[4.4em].

\ifanswer
\ans{$x<1$ 或 $y<1$}
\fi

\item 已知方程 $2x^2+4x-3=0$ 的两个根为 $x_1,x_2$，则 $\dfrac1{x_1}+\dfrac1{x_2}=$\blankline[4.4em].

\ifanswer
\solbegin
\step{利用韦达定理得 $\begin{cases}x_1+x_2=-2\\ x_1\cdot x_2=-\dfrac32\end{cases}$，则 $\dfrac1{x_1}+\dfrac1{x_2}=\dfrac{x_1+x_2}{x_1\cdot x_2}=\dfrac43$.}
\solend
\fi

\item 设 $a>b>0$，比较大小：$\dfrac1{a-b}$\blankline[2.4em]$\dfrac1a$.

\ifanswer
\solbegin
\step{由 $a>b>0$ 得 $a>0$ 且 $a-b>0$，所以 $\dfrac1{a-b}-\dfrac1a=\dfrac b{a(a-b)}>0$，}
\step{故 $\dfrac1{a-b}>\dfrac1a$.}
\solend
\fi

\item 若不等式 $2ax^2-ax+1>0$ 的解集为 $R$，则实数 $a$ 的取值范围是\blankline[4.4em].

\ifanswer
\solbegin
\step{当 $a\neq0$ 时，由题意得 $\begin{cases}a>0\\ \Delta<0\end{cases}$，即 $\begin{cases}a>0\\ a^2-8a<0\end{cases}$，解得 $0<a<8$.}
\step{当 $a=0$ 时，恒有 $1>0$，不等式也成立.}
\step{综上所述，$a$ 的取值范围是 $[0,8)$.}
\solend
\fi

\item 集合 $A=\{x\mid m-1\leqslant x\leqslant m+7\}$，集合 $B=\{x\mid x^2-3x-10\geqslant0\}$，若 $A\cup B=R$，则实数 $m$ 的取值范围为\blankline[4.4em].

\ifanswer
\solbegin
\step{$A=\{x\mid m-1\leqslant x\leqslant m+7\}$，$B=\{x\mid x^2-3x-10\geqslant0\}=\{x\mid x\leqslant-2\text{ 或 }x\geqslant5\}$，}
\step{若 $A\cup B=R$，则 $\begin{cases}m-1\leqslant-2\\ m+7\geqslant5\end{cases}$，解得 $-2\leqslant m\leqslant-1$，}
\step{所以实数 $m$ 的取值范围为 $[-2,-1]$.}
\solend
\fi

\item 若关于 $x$ 的不等式 $ax^2+bx+c>0$ 的解集为 $(-1,3)$，则关于 $x$ 的不等式 $cx^2+bx+a>0$ 的解集为\blankline[4.4em].

\ifanswer
\solbegin
\step{因为不等式 $ax^2+bx+c>0$ 的解集为 $(-1,3)$，}
\step{所以 $-1$ 和 $3$ 是方程 $ax^2+bx+c=0$ 的实数解，}
\step{所以 $\begin{cases}-1+3=-\dfrac ba\\ -1\times3=\dfrac ca\\ a<0\end{cases}$，解得 $b=-2a$，$c=-3a$，且 $a<0$，}
\step{所以不等式 $cx^2+bx+a>0$ 可化为 $-3ax^2-2ax+a>0$，即 $3x^2+2x-1>0$，}
\step{解得 $x<-1$ 或 $x>\dfrac13$，所以该不等式的解集为 $(-\infty,-1)\cup\left(\dfrac13,+\infty\right)$.}
\solend
\fi

\item 若集合 $A=\{x\mid(a-1)x^2-2x+1=0\}$ 有且仅有 $2$ 个子集，则实数 $a$ 的值为\blankline[4.4em].

\ifanswer
\solbegin
\step{由题意得方程 $(a-1)x^2-2x+1=0$ 只有一个根，}
\stepm{①}{当 $a-1=0$ 时，$a=1$，此时方程化为 $-2x+1=0$，解得 $x=\dfrac12$，符合题意，}
\stepm{②}{当 $a-1\neq0$，即 $a\neq1$ 时，$\Delta=4-4(a-1)=0$，解得 $a=2$，}
\step{综上所述，实数 $a$ 的值为 $1$ 或 $2$.}
\solend
\fi

\item “$x+y\neq5$”是“$x\neq3$ 或 $y\neq2$”的\blankline[4.4em]条件.

\ifanswer
\solbegin
\step{分析充分性：}
\step{若 $x+y\neq5$ 成立，若 $x=3$，则 $y\neq2$，若 $y=2$，则 $x\neq3$，}
\step{所以一定能说明 $x\neq3$ 和 $y\neq2$ 至少有一个成立，即 $x\neq3$ 或 $y\neq2$ 成立，}
\step{充分性成立；}
\step{分析必要性：}
\step{若 $x\neq3$ 或 $y\neq2$，即 $x\neq3$ 和 $y\neq2$ 二者至少一个成立，}
\step{举例：$x=1,y=4$，此时 $x+y=5$，所以必要性不成立.}
\step{所以“$x+y\neq5$”是“$x\neq3$ 或 $y\neq2$”的充分不必要条件.}
\solend
\fi

\item 设 $P$ 为非空实数集满足：对任意给定的 $x,y\in P$（$x,y$ 可以相同），都有 $x+y\in P,x-y\in P,xy\in P$，则称 $P$ 为漂亮集. 关于漂亮集有下列结论：

①集合 $P=\{-2,-1,0,1,2\}$ 为漂亮集；②集合 $P=\{x\mid x=2n,n\in\Z\}$ 为漂亮集；

③若集合 $P_1,P_2$ 为漂亮集，则 $P_1\cup P_2$ 为漂亮集；④若集合 $P$ 为漂亮集，则一定有 $0\in P$.

其中所有正确结论的序号是\blankline[4.4em].

\ifanswer
\solbegin
\step{对于集合 $P=\{-2,-1,0,1,2\}$，取 $x=2,y=-2$，则 $x-y=2-(-2)=4\notin P$，}
\step{不满足漂亮集的定义，所以①错误.}
\step{对于集合 $P=\{x\mid x=2n,n\in Z\}$，设 $x=2n_1,y=2n_2,n_1,n_2\in Z$.}
\step{$x+y=2n_1+2n_2=2(n_1+n_2)$，因为 $n_1+n_2\in Z$，所以 $x+y\in P$.}
\step{$x-y=2n_1-2n_2=2(n_1-n_2)$，因为 $n_1-n_2\in Z$，所以 $x-y\in P$.}
\step{$xy=2n_1\times2n_2=4n_1n_2=2(2n_1n_2)$，因为 $2n_1n_2\in Z$，所以 $xy\in P$.}
\step{满足漂亮集的定义，所以②正确.}
\step{令 $P_1=\{x\mid x=2n,n\in Z\}$,$P_2=\{x\mid x=3n,n\in Z\}$，都是漂亮集.}
\step{$P_1\cup P_2$ 中，取 $x=2(x\in P_1)$,$y=3(y\in P_2)$,$x+y=5\notin P_1\cup P_2$，}
\step{不满足漂亮集的定义，所以③错误.}
\step{因为对任意 $x\in P$，有 $x-x=0\in P$，所以若集合 $P$ 为漂亮集，}
\step{则一定有 $0\in P$，所以④正确.}
\step{故正确的是②④.}
\solend
\fi

\item 已知：集合 $A_p=\{(x,y)\mid y=3x^2-px+5p-65,x\in\R\}$,$M=\{x\mid1\leqslant x\leqslant100,x\in\Z\}$，若 $p\in M$，则所有 $A_p$ 的交集\blankline[4.4em].

\ifanswer
\solbegin
\step{由 $M=\{1,2,\cdots,100\}$，所以所有 $A_p$ 的交集就是这些集合中共同拥有的点.}
\step{若 $(x,y)$ 是所有集合的一个公共点，}
\step{那么它特别要同时属于 $p=1$ 和 $p=2$ 时对应的两个集合.}
\step{当 $p=1$ 时，$y=3x^2-x-60$；}
\step{当 $p=2$ 时，$y=3x^2-2x-55$；联立得 $x=5,y=10$.}
\step{因此，如果公共点存在，只可能是 $(5,10)$.}
\step{接下来验证这个点是否属于所有集合，}
\step{对任意 $p\in M$，把 $x=5$ 代入集合的关系式，}
\step{得 $3x^2-px+5p-65=75-5p+5p-65=10$，与 $y=10$ 相符，}
\step{所以 $(5,10)\in A_p$ 对每一个 $p$ 都成立，因此所有 $A_p$ 的交集为 $\{(5,10)\}$.}
\solend
\fi

\end{enumerate}

% 守卫须放在 \bigsec 之前，否则标题仍会落页末（见 高一40 的教训）
\needvspace{5cm}

\bigsec{二、选择题}

\begin{enumerate}
\setcounter{enumi}{12}

\item 已知 $a,b,c\in\R$ 且 $a>b$，则下列不等式一定成立的是\mbox{（\quad）}

\keepwithprev
A. $\dfrac1a<\dfrac1b$\qquad B. $a^2>b^2$

C. $a|c|>b|c|$\qquad D. $\dfrac a{c^2+1}>\dfrac b{c^2+1}$

\ifanswer
\solbegin
\step{$A$：当 $a>0>b$ 时，$\dfrac1a>\dfrac1b$，故 $A$ 错误；}
\step{$B$：当 $a=-1,b=-2$ 时，满足 $a>b$，但 $a^2>b^2$ 不成立，故 $B$ 错误；}
\step{$C$：当 $c=0$ 时，$a|c|=b|c|$，故 $C$ 错误；}
\step{$D$：由 $a>b,c^2+1>0$，得 $\dfrac a{c^2+1}>\dfrac b{c^2+1}$，故 $D$ 正确；}
\step{故选 $D$.}
\solend
\fi

\item 已知集合 $A=\{a^2,0,-1\}$，$B=\{a,b,0\}$，若 $A=B$，则 $(ab)^{2026}$ 的值为\mbox{（\quad）}

\keepwithprev
A. $0$\qquad B. $-1$

C. $1$\qquad D. $\pm1$

\ifanswer
\solbegin
\step{因为集合 $A=\{a^2,0,-1\}$，$B=\{a,b,0\}$，$A=B$，所以 $\begin{cases}a=-1\\ a^2=b\end{cases}$，}
\step{所以 $a=-1$，$b=1$，所以 $(ab)^{2026}=1$，故选 $C$.}
\solend
\fi

\item 设全集 $U=\R$，$A=\{x\mid x<-1\text{ 或 }x>2\}$,$B=\{y\mid y=x^2+a\}$，若 $\overline A\subseteq B$，则实数 $a$ 的取值范围是\mbox{（\quad）}

\keepwithprev
A. $a\leqslant-1$\qquad B. $a>2$

C. $a<-1$\qquad D. $a\geqslant-2$

\ifanswer
\solbegin
\step{因为 $A=\{x\mid x<-1\text{ 或 }x>2\}$，则 $\overline A=\{x\mid-1\leqslant x\leqslant2\}$，}
\step{又 $B=\{y\mid y=x^2+a\}=\{y\mid y\geqslant a\}$，且 $\overline A\subseteq B$，所以 $a\leqslant-1$，故选 $A$.}
\solend
\fi

\item 设 $P$ 和 $Q$ 是两个集合，定义集合 $P-Q=\{x\mid x\in P$，且 $x\notin Q\}$，如果 $Q=\{x\mid(1-x)(x-3)>0\}$，$P=\left\{x\left|\dfrac x{2-x}>0\right.\right\}$，那么 $P-Q$ 等于\mbox{（\quad）}

\keepwithprev
A. $\{x\mid0<x<1\}$\qquad B. $\{x\mid1\leqslant x<2\}$

C. $\{x\mid0<x\leqslant1\}$\qquad D. $\{x\mid2\leqslant x<3\}$

\ifanswer
\solbegin
\step{因为 $P=\left\{x\left|\dfrac x{2-x}>0\right.\right\}=\{x\mid x(x-2)<0\}=\{x\mid0<x<2\}$，}
\step{$Q=\{x\mid(1-x)(x-3)>0\}=\{x\mid(x-3)(x-1)<0\}=\{x\mid1<x<3\}$，}
\step{所以 $P-Q=\{x\mid0<x\leqslant1\}$，故选 $C$.}
\solend
\fi

\end{enumerate}

\bigsec{三、解答题}

\begin{enumerate}
\setcounter{enumi}{16}

\item 已知集合 $A=\{x\mid x^2-9\geqslant0\}$，$B=\{x\mid|x-4|<2\}$，$C=\left\{x\left|\dfrac{x-8}{x+2}<0\right.\right\}$.

（1）求 $A\cup C$；

（2）若全集 $U=R$，求 $\overline A\cap B$.

\ifanswer
\solbegin
\stepm{(1)}{集合 $A=\{x\mid x^2-9\geqslant0\}=\{x\mid x\geqslant3\text{ 或 }x\leqslant-3\}$，}
\step{$B=\{x\mid|x-4|<2\}=\{x\mid2<x<6\}$，$C=\left\{x\left|\dfrac{x-8}{x+2}<0\right.\right\}=\{x\mid-2<x<8\}$，}
\step{所以 $A\cup C=\{x\mid x\leqslant-3\text{ 或 }x>-2\}$；}
\stepm{(2)}{全集 $U=R$，集合 $A=\{x\mid x^2-9\geqslant0\}=\{x\mid x\geqslant3\text{ 或 }x\leqslant-3\}$，}
\step{$B=\{x\mid|x-4|<2\}=\{x\mid2<x<6\}$，}
\step{所以 $\overline A=\{x\mid-3<x<3\}$，所以 $\overline A\cap B=\{x\mid2<x<3\}$.}
\solend
\fi

\ansspace[6]

\item 已知集合 $A=\left\{x\left|\dfrac{2x-1}{x+1}<1,x\in\R\right.\right\}$，集合 $B=\{x\mid(x-a-1)(x-a+1)\leqslant0\}$.

（1）求集合 $A$；

（2）若 $B\subseteq A$，求实数 $a$ 的取值范围.

\ifanswer
\solbegin
\stepm{(1)}{由 $\dfrac{2x-1}{x+1}<1$，得 $\dfrac{x-2}{x+1}<0$ 所以 $(x-2)(x+1)<0$，解得 $-1<x<2$，}
\step{即 $A=(-1,2)$；}
\stepm{(2)}{由 $(x-a-1)(x-a+1)\leqslant0$，得 $a-1\leqslant x\leqslant a+1$，}
\step{所以 $B=\{x\mid a-1\leqslant x\leqslant a+1\}$，}
\step{又 $B\subseteq A$，且 $a+1>a-1$ 即 $B\neq\varnothing$，所以 $\begin{cases}a-1>-1\\ a+1<2\end{cases}$，所以 $0<a<1$，}
\step{所以实数 $a$ 的取值范围是 $\{a\mid0<a<1\}$.}
\solend
\fi

\ansspace[6]

\item 已知集合 $A=\{x\mid x^2-4x+3=0\}$，$B=\{x\mid x^2-ax+3=0\}$.

（1）若 $A\cup B=B$，求实数 $a$ 的值；

（2）若 $A\cap B=B$，求实数 $a$ 的取值范围.

\ifanswer
\solbegin
\stepm{(1)}{集合 $A=\{x\mid x^2-4x+3=0\}=\{1,3\}$，$B=\{x\mid x^2-ax+3=0\}$.}
\step{$A\cup B=B$，所以 $1\in B$，$3\in B$，所以 $\begin{cases}1-a+3=0\\ 9-3a+3=0\end{cases}$，解得实数 $a=4$；}
\stepm{(2)}{因为 $A\cap B=B$，所以 $B\subseteq A$，}
\step{当 $B=\varnothing$ 时，$\Delta=(-a)^2-12<0$，解得 $-2\sqrt3<a<2\sqrt3$，}
\step{当 $B$ 是单元素集时，$\Delta=a^2-12=0$，解得 $a=2\sqrt3$，此时 $B=\{\sqrt3\}$，}
\step{不合题意；}
\step{当 $B=\{1,3\}$ 时，由 (1) 得 $a=4$，}
\step{综上，实数 $a$ 的取值范围是 $(-2\sqrt3,2\sqrt3)\cup\{4\}$.}
\solend
\fi

\ansspace[6]

\item 已知全集为 $R$，实数 $a$、$b$ 满足 $a\neq b$，集合 $M=\{x\mid x^2-3x-4\leqslant0\}$，集合 $A=\{x\mid(x+a)(x+b)>0\}$，集合 $B=\{x\mid x^2+(a-2)x-2a>0\}$.

（1）若 $\overline A=M$，求 $a$、$b$ 的值；

（2）若 $a>b>-1$，求 $A\cap B$；

（3）若 $a^2+a\in\overline B$，求 $a$ 的取值范围.

\ifanswer
\solbegin
\stepm{(1)}{$M=\{x\mid x^2-3x-4\leqslant0\}=\{x\mid-1\leqslant x\leqslant4\}$，$A=\{x\mid(x+a)(x+b)\leqslant0\}$，}
\step{因为 $\overline A=M$，所以 $a=1$，$b=-4$ 或 $a=-4$，$b=1$.}
\stepm{(2)}{因为 $a>b>-1$，所以 $-a<-b<1$，所以 $A=\{x\mid x<-a\text{ 或 }x>-b\}$，}
\step{因为 $B=\{x\mid x^2+(a-2)x-2a>0\}=\{x\mid x<-a\text{ 或 }x>2\}$，}
\step{所以 $A\cap B=\{x\mid x<-a\text{ 或 }x>2\}$，}
\stepm{(3)}{$\overline B=\{x\mid(x-2)(x+a)\leqslant0\}$，}
\step{由 $a^2+a\in\overline B$，得 $(a^2+a-2)(a^2+2a)\leqslant0\Leftrightarrow a(a-1)(a-2)^2\leqslant0$，}
\step{解得 $a=2$ 或 $0\leqslant a\leqslant1$，故 $a$ 的取值范围是 $[0,1]\cup\{2\}$.}
\solend
\fi

\ansspace[8]

\item 对平面直角坐标系的第一象限内的任意两点作如下定义：若 $\dfrac ab>\dfrac cd$，那么称点 $(a,b)$ 是点 $(c,d)$ 的“上位点”，同时点 $(c,d)$ 是点 $(a,b)$ 的“下位点”.

（1）试写出点 $(2,4)$ 的一个“上位点”坐标；

（2）已知点 $(a,b)$ 是点 $(c,d)$ 的“上位点”，判断点 $P(a+c,b+d)$ 是点 $(c,d)$ 的“上位点”还是“下位点”，并证明你的结论；

（3）设正整数 $n$ 满足以下条件：对集合 $\{t\mid0<t<2020,t\in\Z\}$ 内的任意元素 $m$，总存在正整数 $k=2m+1$，使得点 $(n,k)$ 既是点 $(2025,m)$ 的“下位点”，又是点 $(2026,m+1)$ 的“上位点”，问这样的正整数 $n$ 是否存在？若存在，求正整数 $n$ 的值；若不存在，请说明理由.

\ifanswer
\solbegin
\stepm{(1)}{由定义，点 $(a,b)$ 是点 $(2,4)$ 的“上位点”，}
\step{当且仅当 $\dfrac ab>\dfrac24=\dfrac12$ 且 $a>0,b>0$.}
\step{取 $(a,b)=(1,1)$，因为 $\dfrac11>\dfrac12$，}
\step{所以点 $(1,1)$ 是点 $(2,4)$ 的一个“上位点”.（答案不唯一）}
\stepm{(2)}{点 $P(a+c,b+d)$ 是点 $(c,d)$ 的“上位点”.}
\step{$(a,b)$ 是点 $(c,d)$ 的“上位点”，所以 $\dfrac ab>\dfrac cd$.}
\step{因为 $a,b,c,d>0$，两边同乘 $bd$ 得 $ad>bc$，即 $ad-bc>0$.}
\step{作差比较 $\dfrac{a+c}{b+d}$ 与 $\dfrac cd$：$\dfrac{a+c}{b+d}-\dfrac cd=\dfrac{d(a+c)-c(b+d)}{d(b+d)}=\dfrac{ad-bc}{d(b+d)}>0$，}
\step{所以点 $P(a+c,b+d)$ 是点 $(c,d)$ 的“上位点”.}
\stepm{(3)}{由定义，点 $(n,k)$ 是点 $(2025,m)$ 的下位点即 $\dfrac{2025}m>\dfrac nk$；}
\step{点 $(n,k)$ 是点 $(2026,m+1)$ 的上位点即 $\dfrac nk>\dfrac{2026}{m+1}$.}
\step{代入 $k=2m+1$，得对任意 $m\in\{1,2,\cdots,2019\}$，}
\step{都有 $4052-\dfrac{2026}{m+1}<n<4050+\dfrac{2025}m$.}
\step{右边 $4050+\dfrac{2025}m$ 随 $m$ 增大而减小，故最近处在 $m=2019$，}
\step{得 $n<4050+\dfrac{2025}{2019}<4051.01$，所以 $n\leqslant4051$.}
\step{左边 $4052-\dfrac{2026}{m+1}$ 随 $m$ 增大而增大，故最近处也在 $m=2019$，}
\step{得 $n>4052-\dfrac{2026}{2020}>4050.99$，所以 $n\geqslant4051$.}
\step{综上，$n=4051$.}
\step{检验：当 $n=4051$ 时，对任意 $1\leqslant m\leqslant2019,k=2m+1$ 是正整数，}
\step{且 $\dfrac{4051}{2m+1}<\dfrac{2025}m$,$4051m<4050m+2025,m<2025$，成立；}
\step{$\dfrac{4051}{2m+1}>\dfrac{2026}{m+1}$,$4051(m+1)>4052m+2026,m<2025$，成立.}
\step{又各点横纵坐标均为正，都在第一象限内，符合定义，}
\step{所以满足条件的正整数 $n$ 存在，其值为 $n=4051$.}
\solend
\fi

\ansspace*[10]

\end{enumerate}

\end{document}
"
% 紧凑版：xelatex "\def\iscompact{1}% 高一59_华实高中_10月月考.tex
%   —— 高一上数学“2026-2027 年华实高中高一上 10 月月考”（讲评版）
% 试卷版：xelatex 高一59_华实高中_10月月考.tex
% 答案版：xelatex "\def\isanswer{1}% 高一59_华实高中_10月月考.tex
%   —— 高一上数学“2026-2027 年华实高中高一上 10 月月考”（讲评版）
% 试卷版：xelatex 高一59_华实高中_10月月考.tex
% 答案版：xelatex "\def\isanswer{1}% 高一59_华实高中_10月月考.tex
%   —— 高一上数学“2026-2027 年华实高中高一上 10 月月考”（讲评版）
% 试卷版：xelatex 高一59_华实高中_10月月考.tex
% 答案版：xelatex "\def\isanswer{1}\input{高一59_华实高中_10月月考.tex}"
% 紧凑版：xelatex "\def\iscompact{1}\input{高一59_华实高中_10月月考.tex}"
%
% 原始材料：微信公众号“魔都数学加油站”文章，作者破晓，连载“27学年高一59”；
% 原文标题“27学年高一59——2026-2027年上海市华实高中高一上10月月考”；
% 原文链接：https://mp.weixin.qq.com/s/XzyBCmrvpRuMe-NuNk7gkA。
% 页面标注发布时间 2026-10-13 21:45（上海）。
% 正文为 9 张 PNG 页：第 01、02、03、05、06、07、08、09 页 4762×6735，第 04 页 2560×3620
% （**同一篇各页分辨率不同**），均按页序读取。9 页全为试卷正文页（卷首标题在第 1 页页顶，
% 卷末解答题第 21 题解析止于第 9 页），无封面页、目录页、二维码或宣传页。
% 粉色斜排“魔都数学加油站”为卷面水印，与公众号页脚一并未录入。
%
% 卷首标题：照卷面印作“2026--2027 年华实高中高一上 10 月月考”（公众号标题中的“上海市”
% 卷面标题行没有，本稿照卷面），年份区间按全稿惯例排 en dash。原卷只有一行标题、无副标题。
%
% 篇幅与题号：原卷自第 3 题起编号（第 1、2 题不在本卷内）。填空题 3—12（10 题）、
% 选择题 13—16（4 题）、解答题 17—21（5 题），共 19 题。本稿沿用原节与原题号，
% 只改排为全稿统一的大节样式（原卷小节标题为“一、填空题”“二、选择题”“三、解答题”）。
%
% 讲评版：原卷每题下直接印【答案】或【解析】。**只印【答案】的 1 题**（第 3 题），本稿只给 \ans{}；
% 其余 18 题都印【解析】，按全稿口径这类题不另提 \ans{}——第 13、14、15、16 题的结论写在解析末句
% （“故选 D”等），亦照录在解析内。本卷**没有**印在题干括号内的答案。
%
% 副标题：原卷未印副标题。本稿按“知识模块”口径标作“集合与常用逻辑用语、不等式”：
% 19 题中集合与常用逻辑用语 13 题（第 3、9、10、11、12、14、15、16、17、18、19、20、21 题），
% 不等式 10 题（含等式与不等式性质、一元二次不等式、恒成立、基本应用：第 4、5、6、7、8、
% 13、17、18、19、20 题），两类均远超三分之一，故并标。口径同上。
%
% 与原卷的排版差异：
%   ① 原卷每题下直接印【答案】/【解析】，本稿按全稿惯例将答案与解析置于答案版，试卷版保持干净；
%   ② 原卷未印副标题，本稿按“知识模块”口径补出；
%   ③ 原卷填空横线约两个字宽，本稿按全稿惯例统一排 \blankline[4.4em]；
%   ④ 原卷实数集排粗体 $\mathbf{R}$（第 12、13、20 题等），整数集排粗体 $\mathbf{Z}$（第 11、12 题），
%      自然数集排斜体 $N$ 或粗体 $\mathbf{Z}$ 处按本稿惯例改用 $\R$、$\Z$、$\N$；
%      空集记号统一为 \varnothing；
%   ⑤ 原卷在数学式内多处用半角逗号（如 $x_1,x_2$、$a,b\in R$），均照录未强改；
%      句末句点按全稿惯例排西文句点；“漂亮集”“上位点”“下位点”等自定义名词照录；
%   ⑥ 第 17—21 题的子问原卷各自另起一行，本稿按全稿惯例用 enumerate 的 (1)(2)(3)；
%      解析中的子问标号一律用 \stepm{(1)}{…}；第 9、11 题解析中的①、②（③④）分情形或分结论，
%      按全稿惯例作 \stepm{①}{…}、\stepm{②}{…}（第 11 题解析末“故正确的是②④”照录为行内）；
%   ⑦ 第 12 题题干里的“$A_p$”下标原卷作斜体 $p$，照录；
%   ⑧ 每道解答题原卷未留作答空白，本稿按全稿惯例加 \ansspace。
%
% 原卷笔误/存疑（均照录未改，正文对应处另有 % 注）：本卷首轮逐页转录，暂未发现影响答案的笔误。
%
% 插图：无（全卷 9 页无任何图形、无“如图”字样）。
%
\documentclass[a4paper,11pt]{article}
\usepackage{common}

\begin{document}

\examtitlest[3]{2026--2027 年华实高中高一上 10 月月考}{集合与常用逻辑用语、不等式}

\bigsec{一、填空题}

\begin{enumerate}
\setcounter{enumi}{2}

\item “$x\geqslant1$ 且 $y\geqslant1$”的否定形式为\blankline[4.4em].

\ifanswer
\ans{$x<1$ 或 $y<1$}
\fi

\item 已知方程 $2x^2+4x-3=0$ 的两个根为 $x_1,x_2$，则 $\dfrac1{x_1}+\dfrac1{x_2}=$\blankline[4.4em].

\ifanswer
\solbegin
\step{利用韦达定理得 $\begin{cases}x_1+x_2=-2\\ x_1\cdot x_2=-\dfrac32\end{cases}$，则 $\dfrac1{x_1}+\dfrac1{x_2}=\dfrac{x_1+x_2}{x_1\cdot x_2}=\dfrac43$.}
\solend
\fi

\item 设 $a>b>0$，比较大小：$\dfrac1{a-b}$\blankline[2.4em]$\dfrac1a$.

\ifanswer
\solbegin
\step{由 $a>b>0$ 得 $a>0$ 且 $a-b>0$，所以 $\dfrac1{a-b}-\dfrac1a=\dfrac b{a(a-b)}>0$，}
\step{故 $\dfrac1{a-b}>\dfrac1a$.}
\solend
\fi

\item 若不等式 $2ax^2-ax+1>0$ 的解集为 $R$，则实数 $a$ 的取值范围是\blankline[4.4em].

\ifanswer
\solbegin
\step{当 $a\neq0$ 时，由题意得 $\begin{cases}a>0\\ \Delta<0\end{cases}$，即 $\begin{cases}a>0\\ a^2-8a<0\end{cases}$，解得 $0<a<8$.}
\step{当 $a=0$ 时，恒有 $1>0$，不等式也成立.}
\step{综上所述，$a$ 的取值范围是 $[0,8)$.}
\solend
\fi

\item 集合 $A=\{x\mid m-1\leqslant x\leqslant m+7\}$，集合 $B=\{x\mid x^2-3x-10\geqslant0\}$，若 $A\cup B=R$，则实数 $m$ 的取值范围为\blankline[4.4em].

\ifanswer
\solbegin
\step{$A=\{x\mid m-1\leqslant x\leqslant m+7\}$，$B=\{x\mid x^2-3x-10\geqslant0\}=\{x\mid x\leqslant-2\text{ 或 }x\geqslant5\}$，}
\step{若 $A\cup B=R$，则 $\begin{cases}m-1\leqslant-2\\ m+7\geqslant5\end{cases}$，解得 $-2\leqslant m\leqslant-1$，}
\step{所以实数 $m$ 的取值范围为 $[-2,-1]$.}
\solend
\fi

\item 若关于 $x$ 的不等式 $ax^2+bx+c>0$ 的解集为 $(-1,3)$，则关于 $x$ 的不等式 $cx^2+bx+a>0$ 的解集为\blankline[4.4em].

\ifanswer
\solbegin
\step{因为不等式 $ax^2+bx+c>0$ 的解集为 $(-1,3)$，}
\step{所以 $-1$ 和 $3$ 是方程 $ax^2+bx+c=0$ 的实数解，}
\step{所以 $\begin{cases}-1+3=-\dfrac ba\\ -1\times3=\dfrac ca\\ a<0\end{cases}$，解得 $b=-2a$，$c=-3a$，且 $a<0$，}
\step{所以不等式 $cx^2+bx+a>0$ 可化为 $-3ax^2-2ax+a>0$，即 $3x^2+2x-1>0$，}
\step{解得 $x<-1$ 或 $x>\dfrac13$，所以该不等式的解集为 $(-\infty,-1)\cup\left(\dfrac13,+\infty\right)$.}
\solend
\fi

\item 若集合 $A=\{x\mid(a-1)x^2-2x+1=0\}$ 有且仅有 $2$ 个子集，则实数 $a$ 的值为\blankline[4.4em].

\ifanswer
\solbegin
\step{由题意得方程 $(a-1)x^2-2x+1=0$ 只有一个根，}
\stepm{①}{当 $a-1=0$ 时，$a=1$，此时方程化为 $-2x+1=0$，解得 $x=\dfrac12$，符合题意，}
\stepm{②}{当 $a-1\neq0$，即 $a\neq1$ 时，$\Delta=4-4(a-1)=0$，解得 $a=2$，}
\step{综上所述，实数 $a$ 的值为 $1$ 或 $2$.}
\solend
\fi

\item “$x+y\neq5$”是“$x\neq3$ 或 $y\neq2$”的\blankline[4.4em]条件.

\ifanswer
\solbegin
\step{分析充分性：}
\step{若 $x+y\neq5$ 成立，若 $x=3$，则 $y\neq2$，若 $y=2$，则 $x\neq3$，}
\step{所以一定能说明 $x\neq3$ 和 $y\neq2$ 至少有一个成立，即 $x\neq3$ 或 $y\neq2$ 成立，}
\step{充分性成立；}
\step{分析必要性：}
\step{若 $x\neq3$ 或 $y\neq2$，即 $x\neq3$ 和 $y\neq2$ 二者至少一个成立，}
\step{举例：$x=1,y=4$，此时 $x+y=5$，所以必要性不成立.}
\step{所以“$x+y\neq5$”是“$x\neq3$ 或 $y\neq2$”的充分不必要条件.}
\solend
\fi

\item 设 $P$ 为非空实数集满足：对任意给定的 $x,y\in P$（$x,y$ 可以相同），都有 $x+y\in P,x-y\in P,xy\in P$，则称 $P$ 为漂亮集. 关于漂亮集有下列结论：

①集合 $P=\{-2,-1,0,1,2\}$ 为漂亮集；②集合 $P=\{x\mid x=2n,n\in\Z\}$ 为漂亮集；

③若集合 $P_1,P_2$ 为漂亮集，则 $P_1\cup P_2$ 为漂亮集；④若集合 $P$ 为漂亮集，则一定有 $0\in P$.

其中所有正确结论的序号是\blankline[4.4em].

\ifanswer
\solbegin
\step{对于集合 $P=\{-2,-1,0,1,2\}$，取 $x=2,y=-2$，则 $x-y=2-(-2)=4\notin P$，}
\step{不满足漂亮集的定义，所以①错误.}
\step{对于集合 $P=\{x\mid x=2n,n\in Z\}$，设 $x=2n_1,y=2n_2,n_1,n_2\in Z$.}
\step{$x+y=2n_1+2n_2=2(n_1+n_2)$，因为 $n_1+n_2\in Z$，所以 $x+y\in P$.}
\step{$x-y=2n_1-2n_2=2(n_1-n_2)$，因为 $n_1-n_2\in Z$，所以 $x-y\in P$.}
\step{$xy=2n_1\times2n_2=4n_1n_2=2(2n_1n_2)$，因为 $2n_1n_2\in Z$，所以 $xy\in P$.}
\step{满足漂亮集的定义，所以②正确.}
\step{令 $P_1=\{x\mid x=2n,n\in Z\}$,$P_2=\{x\mid x=3n,n\in Z\}$，都是漂亮集.}
\step{$P_1\cup P_2$ 中，取 $x=2(x\in P_1)$,$y=3(y\in P_2)$,$x+y=5\notin P_1\cup P_2$，}
\step{不满足漂亮集的定义，所以③错误.}
\step{因为对任意 $x\in P$，有 $x-x=0\in P$，所以若集合 $P$ 为漂亮集，}
\step{则一定有 $0\in P$，所以④正确.}
\step{故正确的是②④.}
\solend
\fi

\item 已知：集合 $A_p=\{(x,y)\mid y=3x^2-px+5p-65,x\in\R\}$,$M=\{x\mid1\leqslant x\leqslant100,x\in\Z\}$，若 $p\in M$，则所有 $A_p$ 的交集\blankline[4.4em].

\ifanswer
\solbegin
\step{由 $M=\{1,2,\cdots,100\}$，所以所有 $A_p$ 的交集就是这些集合中共同拥有的点.}
\step{若 $(x,y)$ 是所有集合的一个公共点，}
\step{那么它特别要同时属于 $p=1$ 和 $p=2$ 时对应的两个集合.}
\step{当 $p=1$ 时，$y=3x^2-x-60$；}
\step{当 $p=2$ 时，$y=3x^2-2x-55$；联立得 $x=5,y=10$.}
\step{因此，如果公共点存在，只可能是 $(5,10)$.}
\step{接下来验证这个点是否属于所有集合，}
\step{对任意 $p\in M$，把 $x=5$ 代入集合的关系式，}
\step{得 $3x^2-px+5p-65=75-5p+5p-65=10$，与 $y=10$ 相符，}
\step{所以 $(5,10)\in A_p$ 对每一个 $p$ 都成立，因此所有 $A_p$ 的交集为 $\{(5,10)\}$.}
\solend
\fi

\end{enumerate}

% 守卫须放在 \bigsec 之前，否则标题仍会落页末（见 高一40 的教训）
\needvspace{5cm}

\bigsec{二、选择题}

\begin{enumerate}
\setcounter{enumi}{12}

\item 已知 $a,b,c\in\R$ 且 $a>b$，则下列不等式一定成立的是\mbox{（\quad）}

\keepwithprev
A. $\dfrac1a<\dfrac1b$\qquad B. $a^2>b^2$

C. $a|c|>b|c|$\qquad D. $\dfrac a{c^2+1}>\dfrac b{c^2+1}$

\ifanswer
\solbegin
\step{$A$：当 $a>0>b$ 时，$\dfrac1a>\dfrac1b$，故 $A$ 错误；}
\step{$B$：当 $a=-1,b=-2$ 时，满足 $a>b$，但 $a^2>b^2$ 不成立，故 $B$ 错误；}
\step{$C$：当 $c=0$ 时，$a|c|=b|c|$，故 $C$ 错误；}
\step{$D$：由 $a>b,c^2+1>0$，得 $\dfrac a{c^2+1}>\dfrac b{c^2+1}$，故 $D$ 正确；}
\step{故选 $D$.}
\solend
\fi

\item 已知集合 $A=\{a^2,0,-1\}$，$B=\{a,b,0\}$，若 $A=B$，则 $(ab)^{2026}$ 的值为\mbox{（\quad）}

\keepwithprev
A. $0$\qquad B. $-1$

C. $1$\qquad D. $\pm1$

\ifanswer
\solbegin
\step{因为集合 $A=\{a^2,0,-1\}$，$B=\{a,b,0\}$，$A=B$，所以 $\begin{cases}a=-1\\ a^2=b\end{cases}$，}
\step{所以 $a=-1$，$b=1$，所以 $(ab)^{2026}=1$，故选 $C$.}
\solend
\fi

\item 设全集 $U=\R$，$A=\{x\mid x<-1\text{ 或 }x>2\}$,$B=\{y\mid y=x^2+a\}$，若 $\overline A\subseteq B$，则实数 $a$ 的取值范围是\mbox{（\quad）}

\keepwithprev
A. $a\leqslant-1$\qquad B. $a>2$

C. $a<-1$\qquad D. $a\geqslant-2$

\ifanswer
\solbegin
\step{因为 $A=\{x\mid x<-1\text{ 或 }x>2\}$，则 $\overline A=\{x\mid-1\leqslant x\leqslant2\}$，}
\step{又 $B=\{y\mid y=x^2+a\}=\{y\mid y\geqslant a\}$，且 $\overline A\subseteq B$，所以 $a\leqslant-1$，故选 $A$.}
\solend
\fi

\item 设 $P$ 和 $Q$ 是两个集合，定义集合 $P-Q=\{x\mid x\in P$，且 $x\notin Q\}$，如果 $Q=\{x\mid(1-x)(x-3)>0\}$，$P=\left\{x\left|\dfrac x{2-x}>0\right.\right\}$，那么 $P-Q$ 等于\mbox{（\quad）}

\keepwithprev
A. $\{x\mid0<x<1\}$\qquad B. $\{x\mid1\leqslant x<2\}$

C. $\{x\mid0<x\leqslant1\}$\qquad D. $\{x\mid2\leqslant x<3\}$

\ifanswer
\solbegin
\step{因为 $P=\left\{x\left|\dfrac x{2-x}>0\right.\right\}=\{x\mid x(x-2)<0\}=\{x\mid0<x<2\}$，}
\step{$Q=\{x\mid(1-x)(x-3)>0\}=\{x\mid(x-3)(x-1)<0\}=\{x\mid1<x<3\}$，}
\step{所以 $P-Q=\{x\mid0<x\leqslant1\}$，故选 $C$.}
\solend
\fi

\end{enumerate}

\bigsec{三、解答题}

\begin{enumerate}
\setcounter{enumi}{16}

\item 已知集合 $A=\{x\mid x^2-9\geqslant0\}$，$B=\{x\mid|x-4|<2\}$，$C=\left\{x\left|\dfrac{x-8}{x+2}<0\right.\right\}$.

（1）求 $A\cup C$；

（2）若全集 $U=R$，求 $\overline A\cap B$.

\ifanswer
\solbegin
\stepm{(1)}{集合 $A=\{x\mid x^2-9\geqslant0\}=\{x\mid x\geqslant3\text{ 或 }x\leqslant-3\}$，}
\step{$B=\{x\mid|x-4|<2\}=\{x\mid2<x<6\}$，$C=\left\{x\left|\dfrac{x-8}{x+2}<0\right.\right\}=\{x\mid-2<x<8\}$，}
\step{所以 $A\cup C=\{x\mid x\leqslant-3\text{ 或 }x>-2\}$；}
\stepm{(2)}{全集 $U=R$，集合 $A=\{x\mid x^2-9\geqslant0\}=\{x\mid x\geqslant3\text{ 或 }x\leqslant-3\}$，}
\step{$B=\{x\mid|x-4|<2\}=\{x\mid2<x<6\}$，}
\step{所以 $\overline A=\{x\mid-3<x<3\}$，所以 $\overline A\cap B=\{x\mid2<x<3\}$.}
\solend
\fi

\ansspace[6]

\item 已知集合 $A=\left\{x\left|\dfrac{2x-1}{x+1}<1,x\in\R\right.\right\}$，集合 $B=\{x\mid(x-a-1)(x-a+1)\leqslant0\}$.

（1）求集合 $A$；

（2）若 $B\subseteq A$，求实数 $a$ 的取值范围.

\ifanswer
\solbegin
\stepm{(1)}{由 $\dfrac{2x-1}{x+1}<1$，得 $\dfrac{x-2}{x+1}<0$ 所以 $(x-2)(x+1)<0$，解得 $-1<x<2$，}
\step{即 $A=(-1,2)$；}
\stepm{(2)}{由 $(x-a-1)(x-a+1)\leqslant0$，得 $a-1\leqslant x\leqslant a+1$，}
\step{所以 $B=\{x\mid a-1\leqslant x\leqslant a+1\}$，}
\step{又 $B\subseteq A$，且 $a+1>a-1$ 即 $B\neq\varnothing$，所以 $\begin{cases}a-1>-1\\ a+1<2\end{cases}$，所以 $0<a<1$，}
\step{所以实数 $a$ 的取值范围是 $\{a\mid0<a<1\}$.}
\solend
\fi

\ansspace[6]

\item 已知集合 $A=\{x\mid x^2-4x+3=0\}$，$B=\{x\mid x^2-ax+3=0\}$.

（1）若 $A\cup B=B$，求实数 $a$ 的值；

（2）若 $A\cap B=B$，求实数 $a$ 的取值范围.

\ifanswer
\solbegin
\stepm{(1)}{集合 $A=\{x\mid x^2-4x+3=0\}=\{1,3\}$，$B=\{x\mid x^2-ax+3=0\}$.}
\step{$A\cup B=B$，所以 $1\in B$，$3\in B$，所以 $\begin{cases}1-a+3=0\\ 9-3a+3=0\end{cases}$，解得实数 $a=4$；}
\stepm{(2)}{因为 $A\cap B=B$，所以 $B\subseteq A$，}
\step{当 $B=\varnothing$ 时，$\Delta=(-a)^2-12<0$，解得 $-2\sqrt3<a<2\sqrt3$，}
\step{当 $B$ 是单元素集时，$\Delta=a^2-12=0$，解得 $a=2\sqrt3$，此时 $B=\{\sqrt3\}$，}
\step{不合题意；}
\step{当 $B=\{1,3\}$ 时，由 (1) 得 $a=4$，}
\step{综上，实数 $a$ 的取值范围是 $(-2\sqrt3,2\sqrt3)\cup\{4\}$.}
\solend
\fi

\ansspace[6]

\item 已知全集为 $R$，实数 $a$、$b$ 满足 $a\neq b$，集合 $M=\{x\mid x^2-3x-4\leqslant0\}$，集合 $A=\{x\mid(x+a)(x+b)>0\}$，集合 $B=\{x\mid x^2+(a-2)x-2a>0\}$.

（1）若 $\overline A=M$，求 $a$、$b$ 的值；

（2）若 $a>b>-1$，求 $A\cap B$；

（3）若 $a^2+a\in\overline B$，求 $a$ 的取值范围.

\ifanswer
\solbegin
\stepm{(1)}{$M=\{x\mid x^2-3x-4\leqslant0\}=\{x\mid-1\leqslant x\leqslant4\}$，$A=\{x\mid(x+a)(x+b)\leqslant0\}$，}
\step{因为 $\overline A=M$，所以 $a=1$，$b=-4$ 或 $a=-4$，$b=1$.}
\stepm{(2)}{因为 $a>b>-1$，所以 $-a<-b<1$，所以 $A=\{x\mid x<-a\text{ 或 }x>-b\}$，}
\step{因为 $B=\{x\mid x^2+(a-2)x-2a>0\}=\{x\mid x<-a\text{ 或 }x>2\}$，}
\step{所以 $A\cap B=\{x\mid x<-a\text{ 或 }x>2\}$，}
\stepm{(3)}{$\overline B=\{x\mid(x-2)(x+a)\leqslant0\}$，}
\step{由 $a^2+a\in\overline B$，得 $(a^2+a-2)(a^2+2a)\leqslant0\Leftrightarrow a(a-1)(a-2)^2\leqslant0$，}
\step{解得 $a=2$ 或 $0\leqslant a\leqslant1$，故 $a$ 的取值范围是 $[0,1]\cup\{2\}$.}
\solend
\fi

\ansspace[8]

\item 对平面直角坐标系的第一象限内的任意两点作如下定义：若 $\dfrac ab>\dfrac cd$，那么称点 $(a,b)$ 是点 $(c,d)$ 的“上位点”，同时点 $(c,d)$ 是点 $(a,b)$ 的“下位点”.

（1）试写出点 $(2,4)$ 的一个“上位点”坐标；

（2）已知点 $(a,b)$ 是点 $(c,d)$ 的“上位点”，判断点 $P(a+c,b+d)$ 是点 $(c,d)$ 的“上位点”还是“下位点”，并证明你的结论；

（3）设正整数 $n$ 满足以下条件：对集合 $\{t\mid0<t<2020,t\in\Z\}$ 内的任意元素 $m$，总存在正整数 $k=2m+1$，使得点 $(n,k)$ 既是点 $(2025,m)$ 的“下位点”，又是点 $(2026,m+1)$ 的“上位点”，问这样的正整数 $n$ 是否存在？若存在，求正整数 $n$ 的值；若不存在，请说明理由.

\ifanswer
\solbegin
\stepm{(1)}{由定义，点 $(a,b)$ 是点 $(2,4)$ 的“上位点”，}
\step{当且仅当 $\dfrac ab>\dfrac24=\dfrac12$ 且 $a>0,b>0$.}
\step{取 $(a,b)=(1,1)$，因为 $\dfrac11>\dfrac12$，}
\step{所以点 $(1,1)$ 是点 $(2,4)$ 的一个“上位点”.（答案不唯一）}
\stepm{(2)}{点 $P(a+c,b+d)$ 是点 $(c,d)$ 的“上位点”.}
\step{$(a,b)$ 是点 $(c,d)$ 的“上位点”，所以 $\dfrac ab>\dfrac cd$.}
\step{因为 $a,b,c,d>0$，两边同乘 $bd$ 得 $ad>bc$，即 $ad-bc>0$.}
\step{作差比较 $\dfrac{a+c}{b+d}$ 与 $\dfrac cd$：$\dfrac{a+c}{b+d}-\dfrac cd=\dfrac{d(a+c)-c(b+d)}{d(b+d)}=\dfrac{ad-bc}{d(b+d)}>0$，}
\step{所以点 $P(a+c,b+d)$ 是点 $(c,d)$ 的“上位点”.}
\stepm{(3)}{由定义，点 $(n,k)$ 是点 $(2025,m)$ 的下位点即 $\dfrac{2025}m>\dfrac nk$；}
\step{点 $(n,k)$ 是点 $(2026,m+1)$ 的上位点即 $\dfrac nk>\dfrac{2026}{m+1}$.}
\step{代入 $k=2m+1$，得对任意 $m\in\{1,2,\cdots,2019\}$，}
\step{都有 $4052-\dfrac{2026}{m+1}<n<4050+\dfrac{2025}m$.}
\step{右边 $4050+\dfrac{2025}m$ 随 $m$ 增大而减小，故最近处在 $m=2019$，}
\step{得 $n<4050+\dfrac{2025}{2019}<4051.01$，所以 $n\leqslant4051$.}
\step{左边 $4052-\dfrac{2026}{m+1}$ 随 $m$ 增大而增大，故最近处也在 $m=2019$，}
\step{得 $n>4052-\dfrac{2026}{2020}>4050.99$，所以 $n\geqslant4051$.}
\step{综上，$n=4051$.}
\step{检验：当 $n=4051$ 时，对任意 $1\leqslant m\leqslant2019,k=2m+1$ 是正整数，}
\step{且 $\dfrac{4051}{2m+1}<\dfrac{2025}m$,$4051m<4050m+2025,m<2025$，成立；}
\step{$\dfrac{4051}{2m+1}>\dfrac{2026}{m+1}$,$4051(m+1)>4052m+2026,m<2025$，成立.}
\step{又各点横纵坐标均为正，都在第一象限内，符合定义，}
\step{所以满足条件的正整数 $n$ 存在，其值为 $n=4051$.}
\solend
\fi

\ansspace*[10]

\end{enumerate}

\end{document}
"
% 紧凑版：xelatex "\def\iscompact{1}% 高一59_华实高中_10月月考.tex
%   —— 高一上数学“2026-2027 年华实高中高一上 10 月月考”（讲评版）
% 试卷版：xelatex 高一59_华实高中_10月月考.tex
% 答案版：xelatex "\def\isanswer{1}\input{高一59_华实高中_10月月考.tex}"
% 紧凑版：xelatex "\def\iscompact{1}\input{高一59_华实高中_10月月考.tex}"
%
% 原始材料：微信公众号“魔都数学加油站”文章，作者破晓，连载“27学年高一59”；
% 原文标题“27学年高一59——2026-2027年上海市华实高中高一上10月月考”；
% 原文链接：https://mp.weixin.qq.com/s/XzyBCmrvpRuMe-NuNk7gkA。
% 页面标注发布时间 2026-10-13 21:45（上海）。
% 正文为 9 张 PNG 页：第 01、02、03、05、06、07、08、09 页 4762×6735，第 04 页 2560×3620
% （**同一篇各页分辨率不同**），均按页序读取。9 页全为试卷正文页（卷首标题在第 1 页页顶，
% 卷末解答题第 21 题解析止于第 9 页），无封面页、目录页、二维码或宣传页。
% 粉色斜排“魔都数学加油站”为卷面水印，与公众号页脚一并未录入。
%
% 卷首标题：照卷面印作“2026--2027 年华实高中高一上 10 月月考”（公众号标题中的“上海市”
% 卷面标题行没有，本稿照卷面），年份区间按全稿惯例排 en dash。原卷只有一行标题、无副标题。
%
% 篇幅与题号：原卷自第 3 题起编号（第 1、2 题不在本卷内）。填空题 3—12（10 题）、
% 选择题 13—16（4 题）、解答题 17—21（5 题），共 19 题。本稿沿用原节与原题号，
% 只改排为全稿统一的大节样式（原卷小节标题为“一、填空题”“二、选择题”“三、解答题”）。
%
% 讲评版：原卷每题下直接印【答案】或【解析】。**只印【答案】的 1 题**（第 3 题），本稿只给 \ans{}；
% 其余 18 题都印【解析】，按全稿口径这类题不另提 \ans{}——第 13、14、15、16 题的结论写在解析末句
% （“故选 D”等），亦照录在解析内。本卷**没有**印在题干括号内的答案。
%
% 副标题：原卷未印副标题。本稿按“知识模块”口径标作“集合与常用逻辑用语、不等式”：
% 19 题中集合与常用逻辑用语 13 题（第 3、9、10、11、12、14、15、16、17、18、19、20、21 题），
% 不等式 10 题（含等式与不等式性质、一元二次不等式、恒成立、基本应用：第 4、5、6、7、8、
% 13、17、18、19、20 题），两类均远超三分之一，故并标。口径同上。
%
% 与原卷的排版差异：
%   ① 原卷每题下直接印【答案】/【解析】，本稿按全稿惯例将答案与解析置于答案版，试卷版保持干净；
%   ② 原卷未印副标题，本稿按“知识模块”口径补出；
%   ③ 原卷填空横线约两个字宽，本稿按全稿惯例统一排 \blankline[4.4em]；
%   ④ 原卷实数集排粗体 $\mathbf{R}$（第 12、13、20 题等），整数集排粗体 $\mathbf{Z}$（第 11、12 题），
%      自然数集排斜体 $N$ 或粗体 $\mathbf{Z}$ 处按本稿惯例改用 $\R$、$\Z$、$\N$；
%      空集记号统一为 \varnothing；
%   ⑤ 原卷在数学式内多处用半角逗号（如 $x_1,x_2$、$a,b\in R$），均照录未强改；
%      句末句点按全稿惯例排西文句点；“漂亮集”“上位点”“下位点”等自定义名词照录；
%   ⑥ 第 17—21 题的子问原卷各自另起一行，本稿按全稿惯例用 enumerate 的 (1)(2)(3)；
%      解析中的子问标号一律用 \stepm{(1)}{…}；第 9、11 题解析中的①、②（③④）分情形或分结论，
%      按全稿惯例作 \stepm{①}{…}、\stepm{②}{…}（第 11 题解析末“故正确的是②④”照录为行内）；
%   ⑦ 第 12 题题干里的“$A_p$”下标原卷作斜体 $p$，照录；
%   ⑧ 每道解答题原卷未留作答空白，本稿按全稿惯例加 \ansspace。
%
% 原卷笔误/存疑（均照录未改，正文对应处另有 % 注）：本卷首轮逐页转录，暂未发现影响答案的笔误。
%
% 插图：无（全卷 9 页无任何图形、无“如图”字样）。
%
\documentclass[a4paper,11pt]{article}
\usepackage{common}

\begin{document}

\examtitlest[3]{2026--2027 年华实高中高一上 10 月月考}{集合与常用逻辑用语、不等式}

\bigsec{一、填空题}

\begin{enumerate}
\setcounter{enumi}{2}

\item “$x\geqslant1$ 且 $y\geqslant1$”的否定形式为\blankline[4.4em].

\ifanswer
\ans{$x<1$ 或 $y<1$}
\fi

\item 已知方程 $2x^2+4x-3=0$ 的两个根为 $x_1,x_2$，则 $\dfrac1{x_1}+\dfrac1{x_2}=$\blankline[4.4em].

\ifanswer
\solbegin
\step{利用韦达定理得 $\begin{cases}x_1+x_2=-2\\ x_1\cdot x_2=-\dfrac32\end{cases}$，则 $\dfrac1{x_1}+\dfrac1{x_2}=\dfrac{x_1+x_2}{x_1\cdot x_2}=\dfrac43$.}
\solend
\fi

\item 设 $a>b>0$，比较大小：$\dfrac1{a-b}$\blankline[2.4em]$\dfrac1a$.

\ifanswer
\solbegin
\step{由 $a>b>0$ 得 $a>0$ 且 $a-b>0$，所以 $\dfrac1{a-b}-\dfrac1a=\dfrac b{a(a-b)}>0$，}
\step{故 $\dfrac1{a-b}>\dfrac1a$.}
\solend
\fi

\item 若不等式 $2ax^2-ax+1>0$ 的解集为 $R$，则实数 $a$ 的取值范围是\blankline[4.4em].

\ifanswer
\solbegin
\step{当 $a\neq0$ 时，由题意得 $\begin{cases}a>0\\ \Delta<0\end{cases}$，即 $\begin{cases}a>0\\ a^2-8a<0\end{cases}$，解得 $0<a<8$.}
\step{当 $a=0$ 时，恒有 $1>0$，不等式也成立.}
\step{综上所述，$a$ 的取值范围是 $[0,8)$.}
\solend
\fi

\item 集合 $A=\{x\mid m-1\leqslant x\leqslant m+7\}$，集合 $B=\{x\mid x^2-3x-10\geqslant0\}$，若 $A\cup B=R$，则实数 $m$ 的取值范围为\blankline[4.4em].

\ifanswer
\solbegin
\step{$A=\{x\mid m-1\leqslant x\leqslant m+7\}$，$B=\{x\mid x^2-3x-10\geqslant0\}=\{x\mid x\leqslant-2\text{ 或 }x\geqslant5\}$，}
\step{若 $A\cup B=R$，则 $\begin{cases}m-1\leqslant-2\\ m+7\geqslant5\end{cases}$，解得 $-2\leqslant m\leqslant-1$，}
\step{所以实数 $m$ 的取值范围为 $[-2,-1]$.}
\solend
\fi

\item 若关于 $x$ 的不等式 $ax^2+bx+c>0$ 的解集为 $(-1,3)$，则关于 $x$ 的不等式 $cx^2+bx+a>0$ 的解集为\blankline[4.4em].

\ifanswer
\solbegin
\step{因为不等式 $ax^2+bx+c>0$ 的解集为 $(-1,3)$，}
\step{所以 $-1$ 和 $3$ 是方程 $ax^2+bx+c=0$ 的实数解，}
\step{所以 $\begin{cases}-1+3=-\dfrac ba\\ -1\times3=\dfrac ca\\ a<0\end{cases}$，解得 $b=-2a$，$c=-3a$，且 $a<0$，}
\step{所以不等式 $cx^2+bx+a>0$ 可化为 $-3ax^2-2ax+a>0$，即 $3x^2+2x-1>0$，}
\step{解得 $x<-1$ 或 $x>\dfrac13$，所以该不等式的解集为 $(-\infty,-1)\cup\left(\dfrac13,+\infty\right)$.}
\solend
\fi

\item 若集合 $A=\{x\mid(a-1)x^2-2x+1=0\}$ 有且仅有 $2$ 个子集，则实数 $a$ 的值为\blankline[4.4em].

\ifanswer
\solbegin
\step{由题意得方程 $(a-1)x^2-2x+1=0$ 只有一个根，}
\stepm{①}{当 $a-1=0$ 时，$a=1$，此时方程化为 $-2x+1=0$，解得 $x=\dfrac12$，符合题意，}
\stepm{②}{当 $a-1\neq0$，即 $a\neq1$ 时，$\Delta=4-4(a-1)=0$，解得 $a=2$，}
\step{综上所述，实数 $a$ 的值为 $1$ 或 $2$.}
\solend
\fi

\item “$x+y\neq5$”是“$x\neq3$ 或 $y\neq2$”的\blankline[4.4em]条件.

\ifanswer
\solbegin
\step{分析充分性：}
\step{若 $x+y\neq5$ 成立，若 $x=3$，则 $y\neq2$，若 $y=2$，则 $x\neq3$，}
\step{所以一定能说明 $x\neq3$ 和 $y\neq2$ 至少有一个成立，即 $x\neq3$ 或 $y\neq2$ 成立，}
\step{充分性成立；}
\step{分析必要性：}
\step{若 $x\neq3$ 或 $y\neq2$，即 $x\neq3$ 和 $y\neq2$ 二者至少一个成立，}
\step{举例：$x=1,y=4$，此时 $x+y=5$，所以必要性不成立.}
\step{所以“$x+y\neq5$”是“$x\neq3$ 或 $y\neq2$”的充分不必要条件.}
\solend
\fi

\item 设 $P$ 为非空实数集满足：对任意给定的 $x,y\in P$（$x,y$ 可以相同），都有 $x+y\in P,x-y\in P,xy\in P$，则称 $P$ 为漂亮集. 关于漂亮集有下列结论：

①集合 $P=\{-2,-1,0,1,2\}$ 为漂亮集；②集合 $P=\{x\mid x=2n,n\in\Z\}$ 为漂亮集；

③若集合 $P_1,P_2$ 为漂亮集，则 $P_1\cup P_2$ 为漂亮集；④若集合 $P$ 为漂亮集，则一定有 $0\in P$.

其中所有正确结论的序号是\blankline[4.4em].

\ifanswer
\solbegin
\step{对于集合 $P=\{-2,-1,0,1,2\}$，取 $x=2,y=-2$，则 $x-y=2-(-2)=4\notin P$，}
\step{不满足漂亮集的定义，所以①错误.}
\step{对于集合 $P=\{x\mid x=2n,n\in Z\}$，设 $x=2n_1,y=2n_2,n_1,n_2\in Z$.}
\step{$x+y=2n_1+2n_2=2(n_1+n_2)$，因为 $n_1+n_2\in Z$，所以 $x+y\in P$.}
\step{$x-y=2n_1-2n_2=2(n_1-n_2)$，因为 $n_1-n_2\in Z$，所以 $x-y\in P$.}
\step{$xy=2n_1\times2n_2=4n_1n_2=2(2n_1n_2)$，因为 $2n_1n_2\in Z$，所以 $xy\in P$.}
\step{满足漂亮集的定义，所以②正确.}
\step{令 $P_1=\{x\mid x=2n,n\in Z\}$,$P_2=\{x\mid x=3n,n\in Z\}$，都是漂亮集.}
\step{$P_1\cup P_2$ 中，取 $x=2(x\in P_1)$,$y=3(y\in P_2)$,$x+y=5\notin P_1\cup P_2$，}
\step{不满足漂亮集的定义，所以③错误.}
\step{因为对任意 $x\in P$，有 $x-x=0\in P$，所以若集合 $P$ 为漂亮集，}
\step{则一定有 $0\in P$，所以④正确.}
\step{故正确的是②④.}
\solend
\fi

\item 已知：集合 $A_p=\{(x,y)\mid y=3x^2-px+5p-65,x\in\R\}$,$M=\{x\mid1\leqslant x\leqslant100,x\in\Z\}$，若 $p\in M$，则所有 $A_p$ 的交集\blankline[4.4em].

\ifanswer
\solbegin
\step{由 $M=\{1,2,\cdots,100\}$，所以所有 $A_p$ 的交集就是这些集合中共同拥有的点.}
\step{若 $(x,y)$ 是所有集合的一个公共点，}
\step{那么它特别要同时属于 $p=1$ 和 $p=2$ 时对应的两个集合.}
\step{当 $p=1$ 时，$y=3x^2-x-60$；}
\step{当 $p=2$ 时，$y=3x^2-2x-55$；联立得 $x=5,y=10$.}
\step{因此，如果公共点存在，只可能是 $(5,10)$.}
\step{接下来验证这个点是否属于所有集合，}
\step{对任意 $p\in M$，把 $x=5$ 代入集合的关系式，}
\step{得 $3x^2-px+5p-65=75-5p+5p-65=10$，与 $y=10$ 相符，}
\step{所以 $(5,10)\in A_p$ 对每一个 $p$ 都成立，因此所有 $A_p$ 的交集为 $\{(5,10)\}$.}
\solend
\fi

\end{enumerate}

% 守卫须放在 \bigsec 之前，否则标题仍会落页末（见 高一40 的教训）
\needvspace{5cm}

\bigsec{二、选择题}

\begin{enumerate}
\setcounter{enumi}{12}

\item 已知 $a,b,c\in\R$ 且 $a>b$，则下列不等式一定成立的是\mbox{（\quad）}

\keepwithprev
A. $\dfrac1a<\dfrac1b$\qquad B. $a^2>b^2$

C. $a|c|>b|c|$\qquad D. $\dfrac a{c^2+1}>\dfrac b{c^2+1}$

\ifanswer
\solbegin
\step{$A$：当 $a>0>b$ 时，$\dfrac1a>\dfrac1b$，故 $A$ 错误；}
\step{$B$：当 $a=-1,b=-2$ 时，满足 $a>b$，但 $a^2>b^2$ 不成立，故 $B$ 错误；}
\step{$C$：当 $c=0$ 时，$a|c|=b|c|$，故 $C$ 错误；}
\step{$D$：由 $a>b,c^2+1>0$，得 $\dfrac a{c^2+1}>\dfrac b{c^2+1}$，故 $D$ 正确；}
\step{故选 $D$.}
\solend
\fi

\item 已知集合 $A=\{a^2,0,-1\}$，$B=\{a,b,0\}$，若 $A=B$，则 $(ab)^{2026}$ 的值为\mbox{（\quad）}

\keepwithprev
A. $0$\qquad B. $-1$

C. $1$\qquad D. $\pm1$

\ifanswer
\solbegin
\step{因为集合 $A=\{a^2,0,-1\}$，$B=\{a,b,0\}$，$A=B$，所以 $\begin{cases}a=-1\\ a^2=b\end{cases}$，}
\step{所以 $a=-1$，$b=1$，所以 $(ab)^{2026}=1$，故选 $C$.}
\solend
\fi

\item 设全集 $U=\R$，$A=\{x\mid x<-1\text{ 或 }x>2\}$,$B=\{y\mid y=x^2+a\}$，若 $\overline A\subseteq B$，则实数 $a$ 的取值范围是\mbox{（\quad）}

\keepwithprev
A. $a\leqslant-1$\qquad B. $a>2$

C. $a<-1$\qquad D. $a\geqslant-2$

\ifanswer
\solbegin
\step{因为 $A=\{x\mid x<-1\text{ 或 }x>2\}$，则 $\overline A=\{x\mid-1\leqslant x\leqslant2\}$，}
\step{又 $B=\{y\mid y=x^2+a\}=\{y\mid y\geqslant a\}$，且 $\overline A\subseteq B$，所以 $a\leqslant-1$，故选 $A$.}
\solend
\fi

\item 设 $P$ 和 $Q$ 是两个集合，定义集合 $P-Q=\{x\mid x\in P$，且 $x\notin Q\}$，如果 $Q=\{x\mid(1-x)(x-3)>0\}$，$P=\left\{x\left|\dfrac x{2-x}>0\right.\right\}$，那么 $P-Q$ 等于\mbox{（\quad）}

\keepwithprev
A. $\{x\mid0<x<1\}$\qquad B. $\{x\mid1\leqslant x<2\}$

C. $\{x\mid0<x\leqslant1\}$\qquad D. $\{x\mid2\leqslant x<3\}$

\ifanswer
\solbegin
\step{因为 $P=\left\{x\left|\dfrac x{2-x}>0\right.\right\}=\{x\mid x(x-2)<0\}=\{x\mid0<x<2\}$，}
\step{$Q=\{x\mid(1-x)(x-3)>0\}=\{x\mid(x-3)(x-1)<0\}=\{x\mid1<x<3\}$，}
\step{所以 $P-Q=\{x\mid0<x\leqslant1\}$，故选 $C$.}
\solend
\fi

\end{enumerate}

\bigsec{三、解答题}

\begin{enumerate}
\setcounter{enumi}{16}

\item 已知集合 $A=\{x\mid x^2-9\geqslant0\}$，$B=\{x\mid|x-4|<2\}$，$C=\left\{x\left|\dfrac{x-8}{x+2}<0\right.\right\}$.

（1）求 $A\cup C$；

（2）若全集 $U=R$，求 $\overline A\cap B$.

\ifanswer
\solbegin
\stepm{(1)}{集合 $A=\{x\mid x^2-9\geqslant0\}=\{x\mid x\geqslant3\text{ 或 }x\leqslant-3\}$，}
\step{$B=\{x\mid|x-4|<2\}=\{x\mid2<x<6\}$，$C=\left\{x\left|\dfrac{x-8}{x+2}<0\right.\right\}=\{x\mid-2<x<8\}$，}
\step{所以 $A\cup C=\{x\mid x\leqslant-3\text{ 或 }x>-2\}$；}
\stepm{(2)}{全集 $U=R$，集合 $A=\{x\mid x^2-9\geqslant0\}=\{x\mid x\geqslant3\text{ 或 }x\leqslant-3\}$，}
\step{$B=\{x\mid|x-4|<2\}=\{x\mid2<x<6\}$，}
\step{所以 $\overline A=\{x\mid-3<x<3\}$，所以 $\overline A\cap B=\{x\mid2<x<3\}$.}
\solend
\fi

\ansspace[6]

\item 已知集合 $A=\left\{x\left|\dfrac{2x-1}{x+1}<1,x\in\R\right.\right\}$，集合 $B=\{x\mid(x-a-1)(x-a+1)\leqslant0\}$.

（1）求集合 $A$；

（2）若 $B\subseteq A$，求实数 $a$ 的取值范围.

\ifanswer
\solbegin
\stepm{(1)}{由 $\dfrac{2x-1}{x+1}<1$，得 $\dfrac{x-2}{x+1}<0$ 所以 $(x-2)(x+1)<0$，解得 $-1<x<2$，}
\step{即 $A=(-1,2)$；}
\stepm{(2)}{由 $(x-a-1)(x-a+1)\leqslant0$，得 $a-1\leqslant x\leqslant a+1$，}
\step{所以 $B=\{x\mid a-1\leqslant x\leqslant a+1\}$，}
\step{又 $B\subseteq A$，且 $a+1>a-1$ 即 $B\neq\varnothing$，所以 $\begin{cases}a-1>-1\\ a+1<2\end{cases}$，所以 $0<a<1$，}
\step{所以实数 $a$ 的取值范围是 $\{a\mid0<a<1\}$.}
\solend
\fi

\ansspace[6]

\item 已知集合 $A=\{x\mid x^2-4x+3=0\}$，$B=\{x\mid x^2-ax+3=0\}$.

（1）若 $A\cup B=B$，求实数 $a$ 的值；

（2）若 $A\cap B=B$，求实数 $a$ 的取值范围.

\ifanswer
\solbegin
\stepm{(1)}{集合 $A=\{x\mid x^2-4x+3=0\}=\{1,3\}$，$B=\{x\mid x^2-ax+3=0\}$.}
\step{$A\cup B=B$，所以 $1\in B$，$3\in B$，所以 $\begin{cases}1-a+3=0\\ 9-3a+3=0\end{cases}$，解得实数 $a=4$；}
\stepm{(2)}{因为 $A\cap B=B$，所以 $B\subseteq A$，}
\step{当 $B=\varnothing$ 时，$\Delta=(-a)^2-12<0$，解得 $-2\sqrt3<a<2\sqrt3$，}
\step{当 $B$ 是单元素集时，$\Delta=a^2-12=0$，解得 $a=2\sqrt3$，此时 $B=\{\sqrt3\}$，}
\step{不合题意；}
\step{当 $B=\{1,3\}$ 时，由 (1) 得 $a=4$，}
\step{综上，实数 $a$ 的取值范围是 $(-2\sqrt3,2\sqrt3)\cup\{4\}$.}
\solend
\fi

\ansspace[6]

\item 已知全集为 $R$，实数 $a$、$b$ 满足 $a\neq b$，集合 $M=\{x\mid x^2-3x-4\leqslant0\}$，集合 $A=\{x\mid(x+a)(x+b)>0\}$，集合 $B=\{x\mid x^2+(a-2)x-2a>0\}$.

（1）若 $\overline A=M$，求 $a$、$b$ 的值；

（2）若 $a>b>-1$，求 $A\cap B$；

（3）若 $a^2+a\in\overline B$，求 $a$ 的取值范围.

\ifanswer
\solbegin
\stepm{(1)}{$M=\{x\mid x^2-3x-4\leqslant0\}=\{x\mid-1\leqslant x\leqslant4\}$，$A=\{x\mid(x+a)(x+b)\leqslant0\}$，}
\step{因为 $\overline A=M$，所以 $a=1$，$b=-4$ 或 $a=-4$，$b=1$.}
\stepm{(2)}{因为 $a>b>-1$，所以 $-a<-b<1$，所以 $A=\{x\mid x<-a\text{ 或 }x>-b\}$，}
\step{因为 $B=\{x\mid x^2+(a-2)x-2a>0\}=\{x\mid x<-a\text{ 或 }x>2\}$，}
\step{所以 $A\cap B=\{x\mid x<-a\text{ 或 }x>2\}$，}
\stepm{(3)}{$\overline B=\{x\mid(x-2)(x+a)\leqslant0\}$，}
\step{由 $a^2+a\in\overline B$，得 $(a^2+a-2)(a^2+2a)\leqslant0\Leftrightarrow a(a-1)(a-2)^2\leqslant0$，}
\step{解得 $a=2$ 或 $0\leqslant a\leqslant1$，故 $a$ 的取值范围是 $[0,1]\cup\{2\}$.}
\solend
\fi

\ansspace[8]

\item 对平面直角坐标系的第一象限内的任意两点作如下定义：若 $\dfrac ab>\dfrac cd$，那么称点 $(a,b)$ 是点 $(c,d)$ 的“上位点”，同时点 $(c,d)$ 是点 $(a,b)$ 的“下位点”.

（1）试写出点 $(2,4)$ 的一个“上位点”坐标；

（2）已知点 $(a,b)$ 是点 $(c,d)$ 的“上位点”，判断点 $P(a+c,b+d)$ 是点 $(c,d)$ 的“上位点”还是“下位点”，并证明你的结论；

（3）设正整数 $n$ 满足以下条件：对集合 $\{t\mid0<t<2020,t\in\Z\}$ 内的任意元素 $m$，总存在正整数 $k=2m+1$，使得点 $(n,k)$ 既是点 $(2025,m)$ 的“下位点”，又是点 $(2026,m+1)$ 的“上位点”，问这样的正整数 $n$ 是否存在？若存在，求正整数 $n$ 的值；若不存在，请说明理由.

\ifanswer
\solbegin
\stepm{(1)}{由定义，点 $(a,b)$ 是点 $(2,4)$ 的“上位点”，}
\step{当且仅当 $\dfrac ab>\dfrac24=\dfrac12$ 且 $a>0,b>0$.}
\step{取 $(a,b)=(1,1)$，因为 $\dfrac11>\dfrac12$，}
\step{所以点 $(1,1)$ 是点 $(2,4)$ 的一个“上位点”.（答案不唯一）}
\stepm{(2)}{点 $P(a+c,b+d)$ 是点 $(c,d)$ 的“上位点”.}
\step{$(a,b)$ 是点 $(c,d)$ 的“上位点”，所以 $\dfrac ab>\dfrac cd$.}
\step{因为 $a,b,c,d>0$，两边同乘 $bd$ 得 $ad>bc$，即 $ad-bc>0$.}
\step{作差比较 $\dfrac{a+c}{b+d}$ 与 $\dfrac cd$：$\dfrac{a+c}{b+d}-\dfrac cd=\dfrac{d(a+c)-c(b+d)}{d(b+d)}=\dfrac{ad-bc}{d(b+d)}>0$，}
\step{所以点 $P(a+c,b+d)$ 是点 $(c,d)$ 的“上位点”.}
\stepm{(3)}{由定义，点 $(n,k)$ 是点 $(2025,m)$ 的下位点即 $\dfrac{2025}m>\dfrac nk$；}
\step{点 $(n,k)$ 是点 $(2026,m+1)$ 的上位点即 $\dfrac nk>\dfrac{2026}{m+1}$.}
\step{代入 $k=2m+1$，得对任意 $m\in\{1,2,\cdots,2019\}$，}
\step{都有 $4052-\dfrac{2026}{m+1}<n<4050+\dfrac{2025}m$.}
\step{右边 $4050+\dfrac{2025}m$ 随 $m$ 增大而减小，故最近处在 $m=2019$，}
\step{得 $n<4050+\dfrac{2025}{2019}<4051.01$，所以 $n\leqslant4051$.}
\step{左边 $4052-\dfrac{2026}{m+1}$ 随 $m$ 增大而增大，故最近处也在 $m=2019$，}
\step{得 $n>4052-\dfrac{2026}{2020}>4050.99$，所以 $n\geqslant4051$.}
\step{综上，$n=4051$.}
\step{检验：当 $n=4051$ 时，对任意 $1\leqslant m\leqslant2019,k=2m+1$ 是正整数，}
\step{且 $\dfrac{4051}{2m+1}<\dfrac{2025}m$,$4051m<4050m+2025,m<2025$，成立；}
\step{$\dfrac{4051}{2m+1}>\dfrac{2026}{m+1}$,$4051(m+1)>4052m+2026,m<2025$，成立.}
\step{又各点横纵坐标均为正，都在第一象限内，符合定义，}
\step{所以满足条件的正整数 $n$ 存在，其值为 $n=4051$.}
\solend
\fi

\ansspace*[10]

\end{enumerate}

\end{document}
"
%
% 原始材料：微信公众号“魔都数学加油站”文章，作者破晓，连载“27学年高一59”；
% 原文标题“27学年高一59——2026-2027年上海市华实高中高一上10月月考”；
% 原文链接：https://mp.weixin.qq.com/s/XzyBCmrvpRuMe-NuNk7gkA。
% 页面标注发布时间 2026-10-13 21:45（上海）。
% 正文为 9 张 PNG 页：第 01、02、03、05、06、07、08、09 页 4762×6735，第 04 页 2560×3620
% （**同一篇各页分辨率不同**），均按页序读取。9 页全为试卷正文页（卷首标题在第 1 页页顶，
% 卷末解答题第 21 题解析止于第 9 页），无封面页、目录页、二维码或宣传页。
% 粉色斜排“魔都数学加油站”为卷面水印，与公众号页脚一并未录入。
%
% 卷首标题：照卷面印作“2026--2027 年华实高中高一上 10 月月考”（公众号标题中的“上海市”
% 卷面标题行没有，本稿照卷面），年份区间按全稿惯例排 en dash。原卷只有一行标题、无副标题。
%
% 篇幅与题号：原卷自第 3 题起编号（第 1、2 题不在本卷内）。填空题 3—12（10 题）、
% 选择题 13—16（4 题）、解答题 17—21（5 题），共 19 题。本稿沿用原节与原题号，
% 只改排为全稿统一的大节样式（原卷小节标题为“一、填空题”“二、选择题”“三、解答题”）。
%
% 讲评版：原卷每题下直接印【答案】或【解析】。**只印【答案】的 1 题**（第 3 题），本稿只给 \ans{}；
% 其余 18 题都印【解析】，按全稿口径这类题不另提 \ans{}——第 13、14、15、16 题的结论写在解析末句
% （“故选 D”等），亦照录在解析内。本卷**没有**印在题干括号内的答案。
%
% 副标题：原卷未印副标题。本稿按“知识模块”口径标作“集合与常用逻辑用语、不等式”：
% 19 题中集合与常用逻辑用语 13 题（第 3、9、10、11、12、14、15、16、17、18、19、20、21 题），
% 不等式 10 题（含等式与不等式性质、一元二次不等式、恒成立、基本应用：第 4、5、6、7、8、
% 13、17、18、19、20 题），两类均远超三分之一，故并标。口径同上。
%
% 与原卷的排版差异：
%   ① 原卷每题下直接印【答案】/【解析】，本稿按全稿惯例将答案与解析置于答案版，试卷版保持干净；
%   ② 原卷未印副标题，本稿按“知识模块”口径补出；
%   ③ 原卷填空横线约两个字宽，本稿按全稿惯例统一排 \blankline[4.4em]；
%   ④ 原卷实数集排粗体 $\mathbf{R}$（第 12、13、20 题等），整数集排粗体 $\mathbf{Z}$（第 11、12 题），
%      自然数集排斜体 $N$ 或粗体 $\mathbf{Z}$ 处按本稿惯例改用 $\R$、$\Z$、$\N$；
%      空集记号统一为 \varnothing；
%   ⑤ 原卷在数学式内多处用半角逗号（如 $x_1,x_2$、$a,b\in R$），均照录未强改；
%      句末句点按全稿惯例排西文句点；“漂亮集”“上位点”“下位点”等自定义名词照录；
%   ⑥ 第 17—21 题的子问原卷各自另起一行，本稿按全稿惯例用 enumerate 的 (1)(2)(3)；
%      解析中的子问标号一律用 \stepm{(1)}{…}；第 9、11 题解析中的①、②（③④）分情形或分结论，
%      按全稿惯例作 \stepm{①}{…}、\stepm{②}{…}（第 11 题解析末“故正确的是②④”照录为行内）；
%   ⑦ 第 12 题题干里的“$A_p$”下标原卷作斜体 $p$，照录；
%   ⑧ 每道解答题原卷未留作答空白，本稿按全稿惯例加 \ansspace。
%
% 原卷笔误/存疑（均照录未改，正文对应处另有 % 注）：本卷首轮逐页转录，暂未发现影响答案的笔误。
%
% 插图：无（全卷 9 页无任何图形、无“如图”字样）。
%
\documentclass[a4paper,11pt]{article}
\usepackage{common}

\begin{document}

\examtitlest[3]{2026--2027 年华实高中高一上 10 月月考}{集合与常用逻辑用语、不等式}

\bigsec{一、填空题}

\begin{enumerate}
\setcounter{enumi}{2}

\item “$x\geqslant1$ 且 $y\geqslant1$”的否定形式为\blankline[4.4em].

\ifanswer
\ans{$x<1$ 或 $y<1$}
\fi

\item 已知方程 $2x^2+4x-3=0$ 的两个根为 $x_1,x_2$，则 $\dfrac1{x_1}+\dfrac1{x_2}=$\blankline[4.4em].

\ifanswer
\solbegin
\step{利用韦达定理得 $\begin{cases}x_1+x_2=-2\\ x_1\cdot x_2=-\dfrac32\end{cases}$，则 $\dfrac1{x_1}+\dfrac1{x_2}=\dfrac{x_1+x_2}{x_1\cdot x_2}=\dfrac43$.}
\solend
\fi

\item 设 $a>b>0$，比较大小：$\dfrac1{a-b}$\blankline[2.4em]$\dfrac1a$.

\ifanswer
\solbegin
\step{由 $a>b>0$ 得 $a>0$ 且 $a-b>0$，所以 $\dfrac1{a-b}-\dfrac1a=\dfrac b{a(a-b)}>0$，}
\step{故 $\dfrac1{a-b}>\dfrac1a$.}
\solend
\fi

\item 若不等式 $2ax^2-ax+1>0$ 的解集为 $R$，则实数 $a$ 的取值范围是\blankline[4.4em].

\ifanswer
\solbegin
\step{当 $a\neq0$ 时，由题意得 $\begin{cases}a>0\\ \Delta<0\end{cases}$，即 $\begin{cases}a>0\\ a^2-8a<0\end{cases}$，解得 $0<a<8$.}
\step{当 $a=0$ 时，恒有 $1>0$，不等式也成立.}
\step{综上所述，$a$ 的取值范围是 $[0,8)$.}
\solend
\fi

\item 集合 $A=\{x\mid m-1\leqslant x\leqslant m+7\}$，集合 $B=\{x\mid x^2-3x-10\geqslant0\}$，若 $A\cup B=R$，则实数 $m$ 的取值范围为\blankline[4.4em].

\ifanswer
\solbegin
\step{$A=\{x\mid m-1\leqslant x\leqslant m+7\}$，$B=\{x\mid x^2-3x-10\geqslant0\}=\{x\mid x\leqslant-2\text{ 或 }x\geqslant5\}$，}
\step{若 $A\cup B=R$，则 $\begin{cases}m-1\leqslant-2\\ m+7\geqslant5\end{cases}$，解得 $-2\leqslant m\leqslant-1$，}
\step{所以实数 $m$ 的取值范围为 $[-2,-1]$.}
\solend
\fi

\item 若关于 $x$ 的不等式 $ax^2+bx+c>0$ 的解集为 $(-1,3)$，则关于 $x$ 的不等式 $cx^2+bx+a>0$ 的解集为\blankline[4.4em].

\ifanswer
\solbegin
\step{因为不等式 $ax^2+bx+c>0$ 的解集为 $(-1,3)$，}
\step{所以 $-1$ 和 $3$ 是方程 $ax^2+bx+c=0$ 的实数解，}
\step{所以 $\begin{cases}-1+3=-\dfrac ba\\ -1\times3=\dfrac ca\\ a<0\end{cases}$，解得 $b=-2a$，$c=-3a$，且 $a<0$，}
\step{所以不等式 $cx^2+bx+a>0$ 可化为 $-3ax^2-2ax+a>0$，即 $3x^2+2x-1>0$，}
\step{解得 $x<-1$ 或 $x>\dfrac13$，所以该不等式的解集为 $(-\infty,-1)\cup\left(\dfrac13,+\infty\right)$.}
\solend
\fi

\item 若集合 $A=\{x\mid(a-1)x^2-2x+1=0\}$ 有且仅有 $2$ 个子集，则实数 $a$ 的值为\blankline[4.4em].

\ifanswer
\solbegin
\step{由题意得方程 $(a-1)x^2-2x+1=0$ 只有一个根，}
\stepm{①}{当 $a-1=0$ 时，$a=1$，此时方程化为 $-2x+1=0$，解得 $x=\dfrac12$，符合题意，}
\stepm{②}{当 $a-1\neq0$，即 $a\neq1$ 时，$\Delta=4-4(a-1)=0$，解得 $a=2$，}
\step{综上所述，实数 $a$ 的值为 $1$ 或 $2$.}
\solend
\fi

\item “$x+y\neq5$”是“$x\neq3$ 或 $y\neq2$”的\blankline[4.4em]条件.

\ifanswer
\solbegin
\step{分析充分性：}
\step{若 $x+y\neq5$ 成立，若 $x=3$，则 $y\neq2$，若 $y=2$，则 $x\neq3$，}
\step{所以一定能说明 $x\neq3$ 和 $y\neq2$ 至少有一个成立，即 $x\neq3$ 或 $y\neq2$ 成立，}
\step{充分性成立；}
\step{分析必要性：}
\step{若 $x\neq3$ 或 $y\neq2$，即 $x\neq3$ 和 $y\neq2$ 二者至少一个成立，}
\step{举例：$x=1,y=4$，此时 $x+y=5$，所以必要性不成立.}
\step{所以“$x+y\neq5$”是“$x\neq3$ 或 $y\neq2$”的充分不必要条件.}
\solend
\fi

\item 设 $P$ 为非空实数集满足：对任意给定的 $x,y\in P$（$x,y$ 可以相同），都有 $x+y\in P,x-y\in P,xy\in P$，则称 $P$ 为漂亮集. 关于漂亮集有下列结论：

①集合 $P=\{-2,-1,0,1,2\}$ 为漂亮集；②集合 $P=\{x\mid x=2n,n\in\Z\}$ 为漂亮集；

③若集合 $P_1,P_2$ 为漂亮集，则 $P_1\cup P_2$ 为漂亮集；④若集合 $P$ 为漂亮集，则一定有 $0\in P$.

其中所有正确结论的序号是\blankline[4.4em].

\ifanswer
\solbegin
\step{对于集合 $P=\{-2,-1,0,1,2\}$，取 $x=2,y=-2$，则 $x-y=2-(-2)=4\notin P$，}
\step{不满足漂亮集的定义，所以①错误.}
\step{对于集合 $P=\{x\mid x=2n,n\in Z\}$，设 $x=2n_1,y=2n_2,n_1,n_2\in Z$.}
\step{$x+y=2n_1+2n_2=2(n_1+n_2)$，因为 $n_1+n_2\in Z$，所以 $x+y\in P$.}
\step{$x-y=2n_1-2n_2=2(n_1-n_2)$，因为 $n_1-n_2\in Z$，所以 $x-y\in P$.}
\step{$xy=2n_1\times2n_2=4n_1n_2=2(2n_1n_2)$，因为 $2n_1n_2\in Z$，所以 $xy\in P$.}
\step{满足漂亮集的定义，所以②正确.}
\step{令 $P_1=\{x\mid x=2n,n\in Z\}$,$P_2=\{x\mid x=3n,n\in Z\}$，都是漂亮集.}
\step{$P_1\cup P_2$ 中，取 $x=2(x\in P_1)$,$y=3(y\in P_2)$,$x+y=5\notin P_1\cup P_2$，}
\step{不满足漂亮集的定义，所以③错误.}
\step{因为对任意 $x\in P$，有 $x-x=0\in P$，所以若集合 $P$ 为漂亮集，}
\step{则一定有 $0\in P$，所以④正确.}
\step{故正确的是②④.}
\solend
\fi

\item 已知：集合 $A_p=\{(x,y)\mid y=3x^2-px+5p-65,x\in\R\}$,$M=\{x\mid1\leqslant x\leqslant100,x\in\Z\}$，若 $p\in M$，则所有 $A_p$ 的交集\blankline[4.4em].

\ifanswer
\solbegin
\step{由 $M=\{1,2,\cdots,100\}$，所以所有 $A_p$ 的交集就是这些集合中共同拥有的点.}
\step{若 $(x,y)$ 是所有集合的一个公共点，}
\step{那么它特别要同时属于 $p=1$ 和 $p=2$ 时对应的两个集合.}
\step{当 $p=1$ 时，$y=3x^2-x-60$；}
\step{当 $p=2$ 时，$y=3x^2-2x-55$；联立得 $x=5,y=10$.}
\step{因此，如果公共点存在，只可能是 $(5,10)$.}
\step{接下来验证这个点是否属于所有集合，}
\step{对任意 $p\in M$，把 $x=5$ 代入集合的关系式，}
\step{得 $3x^2-px+5p-65=75-5p+5p-65=10$，与 $y=10$ 相符，}
\step{所以 $(5,10)\in A_p$ 对每一个 $p$ 都成立，因此所有 $A_p$ 的交集为 $\{(5,10)\}$.}
\solend
\fi

\end{enumerate}

% 守卫须放在 \bigsec 之前，否则标题仍会落页末（见 高一40 的教训）
\needvspace{5cm}

\bigsec{二、选择题}

\begin{enumerate}
\setcounter{enumi}{12}

\item 已知 $a,b,c\in\R$ 且 $a>b$，则下列不等式一定成立的是\mbox{（\quad）}

\keepwithprev
A. $\dfrac1a<\dfrac1b$\qquad B. $a^2>b^2$

C. $a|c|>b|c|$\qquad D. $\dfrac a{c^2+1}>\dfrac b{c^2+1}$

\ifanswer
\solbegin
\step{$A$：当 $a>0>b$ 时，$\dfrac1a>\dfrac1b$，故 $A$ 错误；}
\step{$B$：当 $a=-1,b=-2$ 时，满足 $a>b$，但 $a^2>b^2$ 不成立，故 $B$ 错误；}
\step{$C$：当 $c=0$ 时，$a|c|=b|c|$，故 $C$ 错误；}
\step{$D$：由 $a>b,c^2+1>0$，得 $\dfrac a{c^2+1}>\dfrac b{c^2+1}$，故 $D$ 正确；}
\step{故选 $D$.}
\solend
\fi

\item 已知集合 $A=\{a^2,0,-1\}$，$B=\{a,b,0\}$，若 $A=B$，则 $(ab)^{2026}$ 的值为\mbox{（\quad）}

\keepwithprev
A. $0$\qquad B. $-1$

C. $1$\qquad D. $\pm1$

\ifanswer
\solbegin
\step{因为集合 $A=\{a^2,0,-1\}$，$B=\{a,b,0\}$，$A=B$，所以 $\begin{cases}a=-1\\ a^2=b\end{cases}$，}
\step{所以 $a=-1$，$b=1$，所以 $(ab)^{2026}=1$，故选 $C$.}
\solend
\fi

\item 设全集 $U=\R$，$A=\{x\mid x<-1\text{ 或 }x>2\}$,$B=\{y\mid y=x^2+a\}$，若 $\overline A\subseteq B$，则实数 $a$ 的取值范围是\mbox{（\quad）}

\keepwithprev
A. $a\leqslant-1$\qquad B. $a>2$

C. $a<-1$\qquad D. $a\geqslant-2$

\ifanswer
\solbegin
\step{因为 $A=\{x\mid x<-1\text{ 或 }x>2\}$，则 $\overline A=\{x\mid-1\leqslant x\leqslant2\}$，}
\step{又 $B=\{y\mid y=x^2+a\}=\{y\mid y\geqslant a\}$，且 $\overline A\subseteq B$，所以 $a\leqslant-1$，故选 $A$.}
\solend
\fi

\item 设 $P$ 和 $Q$ 是两个集合，定义集合 $P-Q=\{x\mid x\in P$，且 $x\notin Q\}$，如果 $Q=\{x\mid(1-x)(x-3)>0\}$，$P=\left\{x\left|\dfrac x{2-x}>0\right.\right\}$，那么 $P-Q$ 等于\mbox{（\quad）}

\keepwithprev
A. $\{x\mid0<x<1\}$\qquad B. $\{x\mid1\leqslant x<2\}$

C. $\{x\mid0<x\leqslant1\}$\qquad D. $\{x\mid2\leqslant x<3\}$

\ifanswer
\solbegin
\step{因为 $P=\left\{x\left|\dfrac x{2-x}>0\right.\right\}=\{x\mid x(x-2)<0\}=\{x\mid0<x<2\}$，}
\step{$Q=\{x\mid(1-x)(x-3)>0\}=\{x\mid(x-3)(x-1)<0\}=\{x\mid1<x<3\}$，}
\step{所以 $P-Q=\{x\mid0<x\leqslant1\}$，故选 $C$.}
\solend
\fi

\end{enumerate}

\bigsec{三、解答题}

\begin{enumerate}
\setcounter{enumi}{16}

\item 已知集合 $A=\{x\mid x^2-9\geqslant0\}$，$B=\{x\mid|x-4|<2\}$，$C=\left\{x\left|\dfrac{x-8}{x+2}<0\right.\right\}$.

（1）求 $A\cup C$；

（2）若全集 $U=R$，求 $\overline A\cap B$.

\ifanswer
\solbegin
\stepm{(1)}{集合 $A=\{x\mid x^2-9\geqslant0\}=\{x\mid x\geqslant3\text{ 或 }x\leqslant-3\}$，}
\step{$B=\{x\mid|x-4|<2\}=\{x\mid2<x<6\}$，$C=\left\{x\left|\dfrac{x-8}{x+2}<0\right.\right\}=\{x\mid-2<x<8\}$，}
\step{所以 $A\cup C=\{x\mid x\leqslant-3\text{ 或 }x>-2\}$；}
\stepm{(2)}{全集 $U=R$，集合 $A=\{x\mid x^2-9\geqslant0\}=\{x\mid x\geqslant3\text{ 或 }x\leqslant-3\}$，}
\step{$B=\{x\mid|x-4|<2\}=\{x\mid2<x<6\}$，}
\step{所以 $\overline A=\{x\mid-3<x<3\}$，所以 $\overline A\cap B=\{x\mid2<x<3\}$.}
\solend
\fi

\ansspace[6]

\item 已知集合 $A=\left\{x\left|\dfrac{2x-1}{x+1}<1,x\in\R\right.\right\}$，集合 $B=\{x\mid(x-a-1)(x-a+1)\leqslant0\}$.

（1）求集合 $A$；

（2）若 $B\subseteq A$，求实数 $a$ 的取值范围.

\ifanswer
\solbegin
\stepm{(1)}{由 $\dfrac{2x-1}{x+1}<1$，得 $\dfrac{x-2}{x+1}<0$ 所以 $(x-2)(x+1)<0$，解得 $-1<x<2$，}
\step{即 $A=(-1,2)$；}
\stepm{(2)}{由 $(x-a-1)(x-a+1)\leqslant0$，得 $a-1\leqslant x\leqslant a+1$，}
\step{所以 $B=\{x\mid a-1\leqslant x\leqslant a+1\}$，}
\step{又 $B\subseteq A$，且 $a+1>a-1$ 即 $B\neq\varnothing$，所以 $\begin{cases}a-1>-1\\ a+1<2\end{cases}$，所以 $0<a<1$，}
\step{所以实数 $a$ 的取值范围是 $\{a\mid0<a<1\}$.}
\solend
\fi

\ansspace[6]

\item 已知集合 $A=\{x\mid x^2-4x+3=0\}$，$B=\{x\mid x^2-ax+3=0\}$.

（1）若 $A\cup B=B$，求实数 $a$ 的值；

（2）若 $A\cap B=B$，求实数 $a$ 的取值范围.

\ifanswer
\solbegin
\stepm{(1)}{集合 $A=\{x\mid x^2-4x+3=0\}=\{1,3\}$，$B=\{x\mid x^2-ax+3=0\}$.}
\step{$A\cup B=B$，所以 $1\in B$，$3\in B$，所以 $\begin{cases}1-a+3=0\\ 9-3a+3=0\end{cases}$，解得实数 $a=4$；}
\stepm{(2)}{因为 $A\cap B=B$，所以 $B\subseteq A$，}
\step{当 $B=\varnothing$ 时，$\Delta=(-a)^2-12<0$，解得 $-2\sqrt3<a<2\sqrt3$，}
\step{当 $B$ 是单元素集时，$\Delta=a^2-12=0$，解得 $a=2\sqrt3$，此时 $B=\{\sqrt3\}$，}
\step{不合题意；}
\step{当 $B=\{1,3\}$ 时，由 (1) 得 $a=4$，}
\step{综上，实数 $a$ 的取值范围是 $(-2\sqrt3,2\sqrt3)\cup\{4\}$.}
\solend
\fi

\ansspace[6]

\item 已知全集为 $R$，实数 $a$、$b$ 满足 $a\neq b$，集合 $M=\{x\mid x^2-3x-4\leqslant0\}$，集合 $A=\{x\mid(x+a)(x+b)>0\}$，集合 $B=\{x\mid x^2+(a-2)x-2a>0\}$.

（1）若 $\overline A=M$，求 $a$、$b$ 的值；

（2）若 $a>b>-1$，求 $A\cap B$；

（3）若 $a^2+a\in\overline B$，求 $a$ 的取值范围.

\ifanswer
\solbegin
\stepm{(1)}{$M=\{x\mid x^2-3x-4\leqslant0\}=\{x\mid-1\leqslant x\leqslant4\}$，$A=\{x\mid(x+a)(x+b)\leqslant0\}$，}
\step{因为 $\overline A=M$，所以 $a=1$，$b=-4$ 或 $a=-4$，$b=1$.}
\stepm{(2)}{因为 $a>b>-1$，所以 $-a<-b<1$，所以 $A=\{x\mid x<-a\text{ 或 }x>-b\}$，}
\step{因为 $B=\{x\mid x^2+(a-2)x-2a>0\}=\{x\mid x<-a\text{ 或 }x>2\}$，}
\step{所以 $A\cap B=\{x\mid x<-a\text{ 或 }x>2\}$，}
\stepm{(3)}{$\overline B=\{x\mid(x-2)(x+a)\leqslant0\}$，}
\step{由 $a^2+a\in\overline B$，得 $(a^2+a-2)(a^2+2a)\leqslant0\Leftrightarrow a(a-1)(a-2)^2\leqslant0$，}
\step{解得 $a=2$ 或 $0\leqslant a\leqslant1$，故 $a$ 的取值范围是 $[0,1]\cup\{2\}$.}
\solend
\fi

\ansspace[8]

\item 对平面直角坐标系的第一象限内的任意两点作如下定义：若 $\dfrac ab>\dfrac cd$，那么称点 $(a,b)$ 是点 $(c,d)$ 的“上位点”，同时点 $(c,d)$ 是点 $(a,b)$ 的“下位点”.

（1）试写出点 $(2,4)$ 的一个“上位点”坐标；

（2）已知点 $(a,b)$ 是点 $(c,d)$ 的“上位点”，判断点 $P(a+c,b+d)$ 是点 $(c,d)$ 的“上位点”还是“下位点”，并证明你的结论；

（3）设正整数 $n$ 满足以下条件：对集合 $\{t\mid0<t<2020,t\in\Z\}$ 内的任意元素 $m$，总存在正整数 $k=2m+1$，使得点 $(n,k)$ 既是点 $(2025,m)$ 的“下位点”，又是点 $(2026,m+1)$ 的“上位点”，问这样的正整数 $n$ 是否存在？若存在，求正整数 $n$ 的值；若不存在，请说明理由.

\ifanswer
\solbegin
\stepm{(1)}{由定义，点 $(a,b)$ 是点 $(2,4)$ 的“上位点”，}
\step{当且仅当 $\dfrac ab>\dfrac24=\dfrac12$ 且 $a>0,b>0$.}
\step{取 $(a,b)=(1,1)$，因为 $\dfrac11>\dfrac12$，}
\step{所以点 $(1,1)$ 是点 $(2,4)$ 的一个“上位点”.（答案不唯一）}
\stepm{(2)}{点 $P(a+c,b+d)$ 是点 $(c,d)$ 的“上位点”.}
\step{$(a,b)$ 是点 $(c,d)$ 的“上位点”，所以 $\dfrac ab>\dfrac cd$.}
\step{因为 $a,b,c,d>0$，两边同乘 $bd$ 得 $ad>bc$，即 $ad-bc>0$.}
\step{作差比较 $\dfrac{a+c}{b+d}$ 与 $\dfrac cd$：$\dfrac{a+c}{b+d}-\dfrac cd=\dfrac{d(a+c)-c(b+d)}{d(b+d)}=\dfrac{ad-bc}{d(b+d)}>0$，}
\step{所以点 $P(a+c,b+d)$ 是点 $(c,d)$ 的“上位点”.}
\stepm{(3)}{由定义，点 $(n,k)$ 是点 $(2025,m)$ 的下位点即 $\dfrac{2025}m>\dfrac nk$；}
\step{点 $(n,k)$ 是点 $(2026,m+1)$ 的上位点即 $\dfrac nk>\dfrac{2026}{m+1}$.}
\step{代入 $k=2m+1$，得对任意 $m\in\{1,2,\cdots,2019\}$，}
\step{都有 $4052-\dfrac{2026}{m+1}<n<4050+\dfrac{2025}m$.}
\step{右边 $4050+\dfrac{2025}m$ 随 $m$ 增大而减小，故最近处在 $m=2019$，}
\step{得 $n<4050+\dfrac{2025}{2019}<4051.01$，所以 $n\leqslant4051$.}
\step{左边 $4052-\dfrac{2026}{m+1}$ 随 $m$ 增大而增大，故最近处也在 $m=2019$，}
\step{得 $n>4052-\dfrac{2026}{2020}>4050.99$，所以 $n\geqslant4051$.}
\step{综上，$n=4051$.}
\step{检验：当 $n=4051$ 时，对任意 $1\leqslant m\leqslant2019,k=2m+1$ 是正整数，}
\step{且 $\dfrac{4051}{2m+1}<\dfrac{2025}m$,$4051m<4050m+2025,m<2025$，成立；}
\step{$\dfrac{4051}{2m+1}>\dfrac{2026}{m+1}$,$4051(m+1)>4052m+2026,m<2025$，成立.}
\step{又各点横纵坐标均为正，都在第一象限内，符合定义，}
\step{所以满足条件的正整数 $n$ 存在，其值为 $n=4051$.}
\solend
\fi

\ansspace*[10]

\end{enumerate}

\end{document}
"
% 紧凑版：xelatex "\def\iscompact{1}% 高一59_华实高中_10月月考.tex
%   —— 高一上数学“2026-2027 年华实高中高一上 10 月月考”（讲评版）
% 试卷版：xelatex 高一59_华实高中_10月月考.tex
% 答案版：xelatex "\def\isanswer{1}% 高一59_华实高中_10月月考.tex
%   —— 高一上数学“2026-2027 年华实高中高一上 10 月月考”（讲评版）
% 试卷版：xelatex 高一59_华实高中_10月月考.tex
% 答案版：xelatex "\def\isanswer{1}\input{高一59_华实高中_10月月考.tex}"
% 紧凑版：xelatex "\def\iscompact{1}\input{高一59_华实高中_10月月考.tex}"
%
% 原始材料：微信公众号“魔都数学加油站”文章，作者破晓，连载“27学年高一59”；
% 原文标题“27学年高一59——2026-2027年上海市华实高中高一上10月月考”；
% 原文链接：https://mp.weixin.qq.com/s/XzyBCmrvpRuMe-NuNk7gkA。
% 页面标注发布时间 2026-10-13 21:45（上海）。
% 正文为 9 张 PNG 页：第 01、02、03、05、06、07、08、09 页 4762×6735，第 04 页 2560×3620
% （**同一篇各页分辨率不同**），均按页序读取。9 页全为试卷正文页（卷首标题在第 1 页页顶，
% 卷末解答题第 21 题解析止于第 9 页），无封面页、目录页、二维码或宣传页。
% 粉色斜排“魔都数学加油站”为卷面水印，与公众号页脚一并未录入。
%
% 卷首标题：照卷面印作“2026--2027 年华实高中高一上 10 月月考”（公众号标题中的“上海市”
% 卷面标题行没有，本稿照卷面），年份区间按全稿惯例排 en dash。原卷只有一行标题、无副标题。
%
% 篇幅与题号：原卷自第 3 题起编号（第 1、2 题不在本卷内）。填空题 3—12（10 题）、
% 选择题 13—16（4 题）、解答题 17—21（5 题），共 19 题。本稿沿用原节与原题号，
% 只改排为全稿统一的大节样式（原卷小节标题为“一、填空题”“二、选择题”“三、解答题”）。
%
% 讲评版：原卷每题下直接印【答案】或【解析】。**只印【答案】的 1 题**（第 3 题），本稿只给 \ans{}；
% 其余 18 题都印【解析】，按全稿口径这类题不另提 \ans{}——第 13、14、15、16 题的结论写在解析末句
% （“故选 D”等），亦照录在解析内。本卷**没有**印在题干括号内的答案。
%
% 副标题：原卷未印副标题。本稿按“知识模块”口径标作“集合与常用逻辑用语、不等式”：
% 19 题中集合与常用逻辑用语 13 题（第 3、9、10、11、12、14、15、16、17、18、19、20、21 题），
% 不等式 10 题（含等式与不等式性质、一元二次不等式、恒成立、基本应用：第 4、5、6、7、8、
% 13、17、18、19、20 题），两类均远超三分之一，故并标。口径同上。
%
% 与原卷的排版差异：
%   ① 原卷每题下直接印【答案】/【解析】，本稿按全稿惯例将答案与解析置于答案版，试卷版保持干净；
%   ② 原卷未印副标题，本稿按“知识模块”口径补出；
%   ③ 原卷填空横线约两个字宽，本稿按全稿惯例统一排 \blankline[4.4em]；
%   ④ 原卷实数集排粗体 $\mathbf{R}$（第 12、13、20 题等），整数集排粗体 $\mathbf{Z}$（第 11、12 题），
%      自然数集排斜体 $N$ 或粗体 $\mathbf{Z}$ 处按本稿惯例改用 $\R$、$\Z$、$\N$；
%      空集记号统一为 \varnothing；
%   ⑤ 原卷在数学式内多处用半角逗号（如 $x_1,x_2$、$a,b\in R$），均照录未强改；
%      句末句点按全稿惯例排西文句点；“漂亮集”“上位点”“下位点”等自定义名词照录；
%   ⑥ 第 17—21 题的子问原卷各自另起一行，本稿按全稿惯例用 enumerate 的 (1)(2)(3)；
%      解析中的子问标号一律用 \stepm{(1)}{…}；第 9、11 题解析中的①、②（③④）分情形或分结论，
%      按全稿惯例作 \stepm{①}{…}、\stepm{②}{…}（第 11 题解析末“故正确的是②④”照录为行内）；
%   ⑦ 第 12 题题干里的“$A_p$”下标原卷作斜体 $p$，照录；
%   ⑧ 每道解答题原卷未留作答空白，本稿按全稿惯例加 \ansspace。
%
% 原卷笔误/存疑（均照录未改，正文对应处另有 % 注）：本卷首轮逐页转录，暂未发现影响答案的笔误。
%
% 插图：无（全卷 9 页无任何图形、无“如图”字样）。
%
\documentclass[a4paper,11pt]{article}
\usepackage{common}

\begin{document}

\examtitlest[3]{2026--2027 年华实高中高一上 10 月月考}{集合与常用逻辑用语、不等式}

\bigsec{一、填空题}

\begin{enumerate}
\setcounter{enumi}{2}

\item “$x\geqslant1$ 且 $y\geqslant1$”的否定形式为\blankline[4.4em].

\ifanswer
\ans{$x<1$ 或 $y<1$}
\fi

\item 已知方程 $2x^2+4x-3=0$ 的两个根为 $x_1,x_2$，则 $\dfrac1{x_1}+\dfrac1{x_2}=$\blankline[4.4em].

\ifanswer
\solbegin
\step{利用韦达定理得 $\begin{cases}x_1+x_2=-2\\ x_1\cdot x_2=-\dfrac32\end{cases}$，则 $\dfrac1{x_1}+\dfrac1{x_2}=\dfrac{x_1+x_2}{x_1\cdot x_2}=\dfrac43$.}
\solend
\fi

\item 设 $a>b>0$，比较大小：$\dfrac1{a-b}$\blankline[2.4em]$\dfrac1a$.

\ifanswer
\solbegin
\step{由 $a>b>0$ 得 $a>0$ 且 $a-b>0$，所以 $\dfrac1{a-b}-\dfrac1a=\dfrac b{a(a-b)}>0$，}
\step{故 $\dfrac1{a-b}>\dfrac1a$.}
\solend
\fi

\item 若不等式 $2ax^2-ax+1>0$ 的解集为 $R$，则实数 $a$ 的取值范围是\blankline[4.4em].

\ifanswer
\solbegin
\step{当 $a\neq0$ 时，由题意得 $\begin{cases}a>0\\ \Delta<0\end{cases}$，即 $\begin{cases}a>0\\ a^2-8a<0\end{cases}$，解得 $0<a<8$.}
\step{当 $a=0$ 时，恒有 $1>0$，不等式也成立.}
\step{综上所述，$a$ 的取值范围是 $[0,8)$.}
\solend
\fi

\item 集合 $A=\{x\mid m-1\leqslant x\leqslant m+7\}$，集合 $B=\{x\mid x^2-3x-10\geqslant0\}$，若 $A\cup B=R$，则实数 $m$ 的取值范围为\blankline[4.4em].

\ifanswer
\solbegin
\step{$A=\{x\mid m-1\leqslant x\leqslant m+7\}$，$B=\{x\mid x^2-3x-10\geqslant0\}=\{x\mid x\leqslant-2\text{ 或 }x\geqslant5\}$，}
\step{若 $A\cup B=R$，则 $\begin{cases}m-1\leqslant-2\\ m+7\geqslant5\end{cases}$，解得 $-2\leqslant m\leqslant-1$，}
\step{所以实数 $m$ 的取值范围为 $[-2,-1]$.}
\solend
\fi

\item 若关于 $x$ 的不等式 $ax^2+bx+c>0$ 的解集为 $(-1,3)$，则关于 $x$ 的不等式 $cx^2+bx+a>0$ 的解集为\blankline[4.4em].

\ifanswer
\solbegin
\step{因为不等式 $ax^2+bx+c>0$ 的解集为 $(-1,3)$，}
\step{所以 $-1$ 和 $3$ 是方程 $ax^2+bx+c=0$ 的实数解，}
\step{所以 $\begin{cases}-1+3=-\dfrac ba\\ -1\times3=\dfrac ca\\ a<0\end{cases}$，解得 $b=-2a$，$c=-3a$，且 $a<0$，}
\step{所以不等式 $cx^2+bx+a>0$ 可化为 $-3ax^2-2ax+a>0$，即 $3x^2+2x-1>0$，}
\step{解得 $x<-1$ 或 $x>\dfrac13$，所以该不等式的解集为 $(-\infty,-1)\cup\left(\dfrac13,+\infty\right)$.}
\solend
\fi

\item 若集合 $A=\{x\mid(a-1)x^2-2x+1=0\}$ 有且仅有 $2$ 个子集，则实数 $a$ 的值为\blankline[4.4em].

\ifanswer
\solbegin
\step{由题意得方程 $(a-1)x^2-2x+1=0$ 只有一个根，}
\stepm{①}{当 $a-1=0$ 时，$a=1$，此时方程化为 $-2x+1=0$，解得 $x=\dfrac12$，符合题意，}
\stepm{②}{当 $a-1\neq0$，即 $a\neq1$ 时，$\Delta=4-4(a-1)=0$，解得 $a=2$，}
\step{综上所述，实数 $a$ 的值为 $1$ 或 $2$.}
\solend
\fi

\item “$x+y\neq5$”是“$x\neq3$ 或 $y\neq2$”的\blankline[4.4em]条件.

\ifanswer
\solbegin
\step{分析充分性：}
\step{若 $x+y\neq5$ 成立，若 $x=3$，则 $y\neq2$，若 $y=2$，则 $x\neq3$，}
\step{所以一定能说明 $x\neq3$ 和 $y\neq2$ 至少有一个成立，即 $x\neq3$ 或 $y\neq2$ 成立，}
\step{充分性成立；}
\step{分析必要性：}
\step{若 $x\neq3$ 或 $y\neq2$，即 $x\neq3$ 和 $y\neq2$ 二者至少一个成立，}
\step{举例：$x=1,y=4$，此时 $x+y=5$，所以必要性不成立.}
\step{所以“$x+y\neq5$”是“$x\neq3$ 或 $y\neq2$”的充分不必要条件.}
\solend
\fi

\item 设 $P$ 为非空实数集满足：对任意给定的 $x,y\in P$（$x,y$ 可以相同），都有 $x+y\in P,x-y\in P,xy\in P$，则称 $P$ 为漂亮集. 关于漂亮集有下列结论：

①集合 $P=\{-2,-1,0,1,2\}$ 为漂亮集；②集合 $P=\{x\mid x=2n,n\in\Z\}$ 为漂亮集；

③若集合 $P_1,P_2$ 为漂亮集，则 $P_1\cup P_2$ 为漂亮集；④若集合 $P$ 为漂亮集，则一定有 $0\in P$.

其中所有正确结论的序号是\blankline[4.4em].

\ifanswer
\solbegin
\step{对于集合 $P=\{-2,-1,0,1,2\}$，取 $x=2,y=-2$，则 $x-y=2-(-2)=4\notin P$，}
\step{不满足漂亮集的定义，所以①错误.}
\step{对于集合 $P=\{x\mid x=2n,n\in Z\}$，设 $x=2n_1,y=2n_2,n_1,n_2\in Z$.}
\step{$x+y=2n_1+2n_2=2(n_1+n_2)$，因为 $n_1+n_2\in Z$，所以 $x+y\in P$.}
\step{$x-y=2n_1-2n_2=2(n_1-n_2)$，因为 $n_1-n_2\in Z$，所以 $x-y\in P$.}
\step{$xy=2n_1\times2n_2=4n_1n_2=2(2n_1n_2)$，因为 $2n_1n_2\in Z$，所以 $xy\in P$.}
\step{满足漂亮集的定义，所以②正确.}
\step{令 $P_1=\{x\mid x=2n,n\in Z\}$,$P_2=\{x\mid x=3n,n\in Z\}$，都是漂亮集.}
\step{$P_1\cup P_2$ 中，取 $x=2(x\in P_1)$,$y=3(y\in P_2)$,$x+y=5\notin P_1\cup P_2$，}
\step{不满足漂亮集的定义，所以③错误.}
\step{因为对任意 $x\in P$，有 $x-x=0\in P$，所以若集合 $P$ 为漂亮集，}
\step{则一定有 $0\in P$，所以④正确.}
\step{故正确的是②④.}
\solend
\fi

\item 已知：集合 $A_p=\{(x,y)\mid y=3x^2-px+5p-65,x\in\R\}$,$M=\{x\mid1\leqslant x\leqslant100,x\in\Z\}$，若 $p\in M$，则所有 $A_p$ 的交集\blankline[4.4em].

\ifanswer
\solbegin
\step{由 $M=\{1,2,\cdots,100\}$，所以所有 $A_p$ 的交集就是这些集合中共同拥有的点.}
\step{若 $(x,y)$ 是所有集合的一个公共点，}
\step{那么它特别要同时属于 $p=1$ 和 $p=2$ 时对应的两个集合.}
\step{当 $p=1$ 时，$y=3x^2-x-60$；}
\step{当 $p=2$ 时，$y=3x^2-2x-55$；联立得 $x=5,y=10$.}
\step{因此，如果公共点存在，只可能是 $(5,10)$.}
\step{接下来验证这个点是否属于所有集合，}
\step{对任意 $p\in M$，把 $x=5$ 代入集合的关系式，}
\step{得 $3x^2-px+5p-65=75-5p+5p-65=10$，与 $y=10$ 相符，}
\step{所以 $(5,10)\in A_p$ 对每一个 $p$ 都成立，因此所有 $A_p$ 的交集为 $\{(5,10)\}$.}
\solend
\fi

\end{enumerate}

% 守卫须放在 \bigsec 之前，否则标题仍会落页末（见 高一40 的教训）
\needvspace{5cm}

\bigsec{二、选择题}

\begin{enumerate}
\setcounter{enumi}{12}

\item 已知 $a,b,c\in\R$ 且 $a>b$，则下列不等式一定成立的是\mbox{（\quad）}

\keepwithprev
A. $\dfrac1a<\dfrac1b$\qquad B. $a^2>b^2$

C. $a|c|>b|c|$\qquad D. $\dfrac a{c^2+1}>\dfrac b{c^2+1}$

\ifanswer
\solbegin
\step{$A$：当 $a>0>b$ 时，$\dfrac1a>\dfrac1b$，故 $A$ 错误；}
\step{$B$：当 $a=-1,b=-2$ 时，满足 $a>b$，但 $a^2>b^2$ 不成立，故 $B$ 错误；}
\step{$C$：当 $c=0$ 时，$a|c|=b|c|$，故 $C$ 错误；}
\step{$D$：由 $a>b,c^2+1>0$，得 $\dfrac a{c^2+1}>\dfrac b{c^2+1}$，故 $D$ 正确；}
\step{故选 $D$.}
\solend
\fi

\item 已知集合 $A=\{a^2,0,-1\}$，$B=\{a,b,0\}$，若 $A=B$，则 $(ab)^{2026}$ 的值为\mbox{（\quad）}

\keepwithprev
A. $0$\qquad B. $-1$

C. $1$\qquad D. $\pm1$

\ifanswer
\solbegin
\step{因为集合 $A=\{a^2,0,-1\}$，$B=\{a,b,0\}$，$A=B$，所以 $\begin{cases}a=-1\\ a^2=b\end{cases}$，}
\step{所以 $a=-1$，$b=1$，所以 $(ab)^{2026}=1$，故选 $C$.}
\solend
\fi

\item 设全集 $U=\R$，$A=\{x\mid x<-1\text{ 或 }x>2\}$,$B=\{y\mid y=x^2+a\}$，若 $\overline A\subseteq B$，则实数 $a$ 的取值范围是\mbox{（\quad）}

\keepwithprev
A. $a\leqslant-1$\qquad B. $a>2$

C. $a<-1$\qquad D. $a\geqslant-2$

\ifanswer
\solbegin
\step{因为 $A=\{x\mid x<-1\text{ 或 }x>2\}$，则 $\overline A=\{x\mid-1\leqslant x\leqslant2\}$，}
\step{又 $B=\{y\mid y=x^2+a\}=\{y\mid y\geqslant a\}$，且 $\overline A\subseteq B$，所以 $a\leqslant-1$，故选 $A$.}
\solend
\fi

\item 设 $P$ 和 $Q$ 是两个集合，定义集合 $P-Q=\{x\mid x\in P$，且 $x\notin Q\}$，如果 $Q=\{x\mid(1-x)(x-3)>0\}$，$P=\left\{x\left|\dfrac x{2-x}>0\right.\right\}$，那么 $P-Q$ 等于\mbox{（\quad）}

\keepwithprev
A. $\{x\mid0<x<1\}$\qquad B. $\{x\mid1\leqslant x<2\}$

C. $\{x\mid0<x\leqslant1\}$\qquad D. $\{x\mid2\leqslant x<3\}$

\ifanswer
\solbegin
\step{因为 $P=\left\{x\left|\dfrac x{2-x}>0\right.\right\}=\{x\mid x(x-2)<0\}=\{x\mid0<x<2\}$，}
\step{$Q=\{x\mid(1-x)(x-3)>0\}=\{x\mid(x-3)(x-1)<0\}=\{x\mid1<x<3\}$，}
\step{所以 $P-Q=\{x\mid0<x\leqslant1\}$，故选 $C$.}
\solend
\fi

\end{enumerate}

\bigsec{三、解答题}

\begin{enumerate}
\setcounter{enumi}{16}

\item 已知集合 $A=\{x\mid x^2-9\geqslant0\}$，$B=\{x\mid|x-4|<2\}$，$C=\left\{x\left|\dfrac{x-8}{x+2}<0\right.\right\}$.

（1）求 $A\cup C$；

（2）若全集 $U=R$，求 $\overline A\cap B$.

\ifanswer
\solbegin
\stepm{(1)}{集合 $A=\{x\mid x^2-9\geqslant0\}=\{x\mid x\geqslant3\text{ 或 }x\leqslant-3\}$，}
\step{$B=\{x\mid|x-4|<2\}=\{x\mid2<x<6\}$，$C=\left\{x\left|\dfrac{x-8}{x+2}<0\right.\right\}=\{x\mid-2<x<8\}$，}
\step{所以 $A\cup C=\{x\mid x\leqslant-3\text{ 或 }x>-2\}$；}
\stepm{(2)}{全集 $U=R$，集合 $A=\{x\mid x^2-9\geqslant0\}=\{x\mid x\geqslant3\text{ 或 }x\leqslant-3\}$，}
\step{$B=\{x\mid|x-4|<2\}=\{x\mid2<x<6\}$，}
\step{所以 $\overline A=\{x\mid-3<x<3\}$，所以 $\overline A\cap B=\{x\mid2<x<3\}$.}
\solend
\fi

\ansspace[6]

\item 已知集合 $A=\left\{x\left|\dfrac{2x-1}{x+1}<1,x\in\R\right.\right\}$，集合 $B=\{x\mid(x-a-1)(x-a+1)\leqslant0\}$.

（1）求集合 $A$；

（2）若 $B\subseteq A$，求实数 $a$ 的取值范围.

\ifanswer
\solbegin
\stepm{(1)}{由 $\dfrac{2x-1}{x+1}<1$，得 $\dfrac{x-2}{x+1}<0$ 所以 $(x-2)(x+1)<0$，解得 $-1<x<2$，}
\step{即 $A=(-1,2)$；}
\stepm{(2)}{由 $(x-a-1)(x-a+1)\leqslant0$，得 $a-1\leqslant x\leqslant a+1$，}
\step{所以 $B=\{x\mid a-1\leqslant x\leqslant a+1\}$，}
\step{又 $B\subseteq A$，且 $a+1>a-1$ 即 $B\neq\varnothing$，所以 $\begin{cases}a-1>-1\\ a+1<2\end{cases}$，所以 $0<a<1$，}
\step{所以实数 $a$ 的取值范围是 $\{a\mid0<a<1\}$.}
\solend
\fi

\ansspace[6]

\item 已知集合 $A=\{x\mid x^2-4x+3=0\}$，$B=\{x\mid x^2-ax+3=0\}$.

（1）若 $A\cup B=B$，求实数 $a$ 的值；

（2）若 $A\cap B=B$，求实数 $a$ 的取值范围.

\ifanswer
\solbegin
\stepm{(1)}{集合 $A=\{x\mid x^2-4x+3=0\}=\{1,3\}$，$B=\{x\mid x^2-ax+3=0\}$.}
\step{$A\cup B=B$，所以 $1\in B$，$3\in B$，所以 $\begin{cases}1-a+3=0\\ 9-3a+3=0\end{cases}$，解得实数 $a=4$；}
\stepm{(2)}{因为 $A\cap B=B$，所以 $B\subseteq A$，}
\step{当 $B=\varnothing$ 时，$\Delta=(-a)^2-12<0$，解得 $-2\sqrt3<a<2\sqrt3$，}
\step{当 $B$ 是单元素集时，$\Delta=a^2-12=0$，解得 $a=2\sqrt3$，此时 $B=\{\sqrt3\}$，}
\step{不合题意；}
\step{当 $B=\{1,3\}$ 时，由 (1) 得 $a=4$，}
\step{综上，实数 $a$ 的取值范围是 $(-2\sqrt3,2\sqrt3)\cup\{4\}$.}
\solend
\fi

\ansspace[6]

\item 已知全集为 $R$，实数 $a$、$b$ 满足 $a\neq b$，集合 $M=\{x\mid x^2-3x-4\leqslant0\}$，集合 $A=\{x\mid(x+a)(x+b)>0\}$，集合 $B=\{x\mid x^2+(a-2)x-2a>0\}$.

（1）若 $\overline A=M$，求 $a$、$b$ 的值；

（2）若 $a>b>-1$，求 $A\cap B$；

（3）若 $a^2+a\in\overline B$，求 $a$ 的取值范围.

\ifanswer
\solbegin
\stepm{(1)}{$M=\{x\mid x^2-3x-4\leqslant0\}=\{x\mid-1\leqslant x\leqslant4\}$，$A=\{x\mid(x+a)(x+b)\leqslant0\}$，}
\step{因为 $\overline A=M$，所以 $a=1$，$b=-4$ 或 $a=-4$，$b=1$.}
\stepm{(2)}{因为 $a>b>-1$，所以 $-a<-b<1$，所以 $A=\{x\mid x<-a\text{ 或 }x>-b\}$，}
\step{因为 $B=\{x\mid x^2+(a-2)x-2a>0\}=\{x\mid x<-a\text{ 或 }x>2\}$，}
\step{所以 $A\cap B=\{x\mid x<-a\text{ 或 }x>2\}$，}
\stepm{(3)}{$\overline B=\{x\mid(x-2)(x+a)\leqslant0\}$，}
\step{由 $a^2+a\in\overline B$，得 $(a^2+a-2)(a^2+2a)\leqslant0\Leftrightarrow a(a-1)(a-2)^2\leqslant0$，}
\step{解得 $a=2$ 或 $0\leqslant a\leqslant1$，故 $a$ 的取值范围是 $[0,1]\cup\{2\}$.}
\solend
\fi

\ansspace[8]

\item 对平面直角坐标系的第一象限内的任意两点作如下定义：若 $\dfrac ab>\dfrac cd$，那么称点 $(a,b)$ 是点 $(c,d)$ 的“上位点”，同时点 $(c,d)$ 是点 $(a,b)$ 的“下位点”.

（1）试写出点 $(2,4)$ 的一个“上位点”坐标；

（2）已知点 $(a,b)$ 是点 $(c,d)$ 的“上位点”，判断点 $P(a+c,b+d)$ 是点 $(c,d)$ 的“上位点”还是“下位点”，并证明你的结论；

（3）设正整数 $n$ 满足以下条件：对集合 $\{t\mid0<t<2020,t\in\Z\}$ 内的任意元素 $m$，总存在正整数 $k=2m+1$，使得点 $(n,k)$ 既是点 $(2025,m)$ 的“下位点”，又是点 $(2026,m+1)$ 的“上位点”，问这样的正整数 $n$ 是否存在？若存在，求正整数 $n$ 的值；若不存在，请说明理由.

\ifanswer
\solbegin
\stepm{(1)}{由定义，点 $(a,b)$ 是点 $(2,4)$ 的“上位点”，}
\step{当且仅当 $\dfrac ab>\dfrac24=\dfrac12$ 且 $a>0,b>0$.}
\step{取 $(a,b)=(1,1)$，因为 $\dfrac11>\dfrac12$，}
\step{所以点 $(1,1)$ 是点 $(2,4)$ 的一个“上位点”.（答案不唯一）}
\stepm{(2)}{点 $P(a+c,b+d)$ 是点 $(c,d)$ 的“上位点”.}
\step{$(a,b)$ 是点 $(c,d)$ 的“上位点”，所以 $\dfrac ab>\dfrac cd$.}
\step{因为 $a,b,c,d>0$，两边同乘 $bd$ 得 $ad>bc$，即 $ad-bc>0$.}
\step{作差比较 $\dfrac{a+c}{b+d}$ 与 $\dfrac cd$：$\dfrac{a+c}{b+d}-\dfrac cd=\dfrac{d(a+c)-c(b+d)}{d(b+d)}=\dfrac{ad-bc}{d(b+d)}>0$，}
\step{所以点 $P(a+c,b+d)$ 是点 $(c,d)$ 的“上位点”.}
\stepm{(3)}{由定义，点 $(n,k)$ 是点 $(2025,m)$ 的下位点即 $\dfrac{2025}m>\dfrac nk$；}
\step{点 $(n,k)$ 是点 $(2026,m+1)$ 的上位点即 $\dfrac nk>\dfrac{2026}{m+1}$.}
\step{代入 $k=2m+1$，得对任意 $m\in\{1,2,\cdots,2019\}$，}
\step{都有 $4052-\dfrac{2026}{m+1}<n<4050+\dfrac{2025}m$.}
\step{右边 $4050+\dfrac{2025}m$ 随 $m$ 增大而减小，故最近处在 $m=2019$，}
\step{得 $n<4050+\dfrac{2025}{2019}<4051.01$，所以 $n\leqslant4051$.}
\step{左边 $4052-\dfrac{2026}{m+1}$ 随 $m$ 增大而增大，故最近处也在 $m=2019$，}
\step{得 $n>4052-\dfrac{2026}{2020}>4050.99$，所以 $n\geqslant4051$.}
\step{综上，$n=4051$.}
\step{检验：当 $n=4051$ 时，对任意 $1\leqslant m\leqslant2019,k=2m+1$ 是正整数，}
\step{且 $\dfrac{4051}{2m+1}<\dfrac{2025}m$,$4051m<4050m+2025,m<2025$，成立；}
\step{$\dfrac{4051}{2m+1}>\dfrac{2026}{m+1}$,$4051(m+1)>4052m+2026,m<2025$，成立.}
\step{又各点横纵坐标均为正，都在第一象限内，符合定义，}
\step{所以满足条件的正整数 $n$ 存在，其值为 $n=4051$.}
\solend
\fi

\ansspace*[10]

\end{enumerate}

\end{document}
"
% 紧凑版：xelatex "\def\iscompact{1}% 高一59_华实高中_10月月考.tex
%   —— 高一上数学“2026-2027 年华实高中高一上 10 月月考”（讲评版）
% 试卷版：xelatex 高一59_华实高中_10月月考.tex
% 答案版：xelatex "\def\isanswer{1}\input{高一59_华实高中_10月月考.tex}"
% 紧凑版：xelatex "\def\iscompact{1}\input{高一59_华实高中_10月月考.tex}"
%
% 原始材料：微信公众号“魔都数学加油站”文章，作者破晓，连载“27学年高一59”；
% 原文标题“27学年高一59——2026-2027年上海市华实高中高一上10月月考”；
% 原文链接：https://mp.weixin.qq.com/s/XzyBCmrvpRuMe-NuNk7gkA。
% 页面标注发布时间 2026-10-13 21:45（上海）。
% 正文为 9 张 PNG 页：第 01、02、03、05、06、07、08、09 页 4762×6735，第 04 页 2560×3620
% （**同一篇各页分辨率不同**），均按页序读取。9 页全为试卷正文页（卷首标题在第 1 页页顶，
% 卷末解答题第 21 题解析止于第 9 页），无封面页、目录页、二维码或宣传页。
% 粉色斜排“魔都数学加油站”为卷面水印，与公众号页脚一并未录入。
%
% 卷首标题：照卷面印作“2026--2027 年华实高中高一上 10 月月考”（公众号标题中的“上海市”
% 卷面标题行没有，本稿照卷面），年份区间按全稿惯例排 en dash。原卷只有一行标题、无副标题。
%
% 篇幅与题号：原卷自第 3 题起编号（第 1、2 题不在本卷内）。填空题 3—12（10 题）、
% 选择题 13—16（4 题）、解答题 17—21（5 题），共 19 题。本稿沿用原节与原题号，
% 只改排为全稿统一的大节样式（原卷小节标题为“一、填空题”“二、选择题”“三、解答题”）。
%
% 讲评版：原卷每题下直接印【答案】或【解析】。**只印【答案】的 1 题**（第 3 题），本稿只给 \ans{}；
% 其余 18 题都印【解析】，按全稿口径这类题不另提 \ans{}——第 13、14、15、16 题的结论写在解析末句
% （“故选 D”等），亦照录在解析内。本卷**没有**印在题干括号内的答案。
%
% 副标题：原卷未印副标题。本稿按“知识模块”口径标作“集合与常用逻辑用语、不等式”：
% 19 题中集合与常用逻辑用语 13 题（第 3、9、10、11、12、14、15、16、17、18、19、20、21 题），
% 不等式 10 题（含等式与不等式性质、一元二次不等式、恒成立、基本应用：第 4、5、6、7、8、
% 13、17、18、19、20 题），两类均远超三分之一，故并标。口径同上。
%
% 与原卷的排版差异：
%   ① 原卷每题下直接印【答案】/【解析】，本稿按全稿惯例将答案与解析置于答案版，试卷版保持干净；
%   ② 原卷未印副标题，本稿按“知识模块”口径补出；
%   ③ 原卷填空横线约两个字宽，本稿按全稿惯例统一排 \blankline[4.4em]；
%   ④ 原卷实数集排粗体 $\mathbf{R}$（第 12、13、20 题等），整数集排粗体 $\mathbf{Z}$（第 11、12 题），
%      自然数集排斜体 $N$ 或粗体 $\mathbf{Z}$ 处按本稿惯例改用 $\R$、$\Z$、$\N$；
%      空集记号统一为 \varnothing；
%   ⑤ 原卷在数学式内多处用半角逗号（如 $x_1,x_2$、$a,b\in R$），均照录未强改；
%      句末句点按全稿惯例排西文句点；“漂亮集”“上位点”“下位点”等自定义名词照录；
%   ⑥ 第 17—21 题的子问原卷各自另起一行，本稿按全稿惯例用 enumerate 的 (1)(2)(3)；
%      解析中的子问标号一律用 \stepm{(1)}{…}；第 9、11 题解析中的①、②（③④）分情形或分结论，
%      按全稿惯例作 \stepm{①}{…}、\stepm{②}{…}（第 11 题解析末“故正确的是②④”照录为行内）；
%   ⑦ 第 12 题题干里的“$A_p$”下标原卷作斜体 $p$，照录；
%   ⑧ 每道解答题原卷未留作答空白，本稿按全稿惯例加 \ansspace。
%
% 原卷笔误/存疑（均照录未改，正文对应处另有 % 注）：本卷首轮逐页转录，暂未发现影响答案的笔误。
%
% 插图：无（全卷 9 页无任何图形、无“如图”字样）。
%
\documentclass[a4paper,11pt]{article}
\usepackage{common}

\begin{document}

\examtitlest[3]{2026--2027 年华实高中高一上 10 月月考}{集合与常用逻辑用语、不等式}

\bigsec{一、填空题}

\begin{enumerate}
\setcounter{enumi}{2}

\item “$x\geqslant1$ 且 $y\geqslant1$”的否定形式为\blankline[4.4em].

\ifanswer
\ans{$x<1$ 或 $y<1$}
\fi

\item 已知方程 $2x^2+4x-3=0$ 的两个根为 $x_1,x_2$，则 $\dfrac1{x_1}+\dfrac1{x_2}=$\blankline[4.4em].

\ifanswer
\solbegin
\step{利用韦达定理得 $\begin{cases}x_1+x_2=-2\\ x_1\cdot x_2=-\dfrac32\end{cases}$，则 $\dfrac1{x_1}+\dfrac1{x_2}=\dfrac{x_1+x_2}{x_1\cdot x_2}=\dfrac43$.}
\solend
\fi

\item 设 $a>b>0$，比较大小：$\dfrac1{a-b}$\blankline[2.4em]$\dfrac1a$.

\ifanswer
\solbegin
\step{由 $a>b>0$ 得 $a>0$ 且 $a-b>0$，所以 $\dfrac1{a-b}-\dfrac1a=\dfrac b{a(a-b)}>0$，}
\step{故 $\dfrac1{a-b}>\dfrac1a$.}
\solend
\fi

\item 若不等式 $2ax^2-ax+1>0$ 的解集为 $R$，则实数 $a$ 的取值范围是\blankline[4.4em].

\ifanswer
\solbegin
\step{当 $a\neq0$ 时，由题意得 $\begin{cases}a>0\\ \Delta<0\end{cases}$，即 $\begin{cases}a>0\\ a^2-8a<0\end{cases}$，解得 $0<a<8$.}
\step{当 $a=0$ 时，恒有 $1>0$，不等式也成立.}
\step{综上所述，$a$ 的取值范围是 $[0,8)$.}
\solend
\fi

\item 集合 $A=\{x\mid m-1\leqslant x\leqslant m+7\}$，集合 $B=\{x\mid x^2-3x-10\geqslant0\}$，若 $A\cup B=R$，则实数 $m$ 的取值范围为\blankline[4.4em].

\ifanswer
\solbegin
\step{$A=\{x\mid m-1\leqslant x\leqslant m+7\}$，$B=\{x\mid x^2-3x-10\geqslant0\}=\{x\mid x\leqslant-2\text{ 或 }x\geqslant5\}$，}
\step{若 $A\cup B=R$，则 $\begin{cases}m-1\leqslant-2\\ m+7\geqslant5\end{cases}$，解得 $-2\leqslant m\leqslant-1$，}
\step{所以实数 $m$ 的取值范围为 $[-2,-1]$.}
\solend
\fi

\item 若关于 $x$ 的不等式 $ax^2+bx+c>0$ 的解集为 $(-1,3)$，则关于 $x$ 的不等式 $cx^2+bx+a>0$ 的解集为\blankline[4.4em].

\ifanswer
\solbegin
\step{因为不等式 $ax^2+bx+c>0$ 的解集为 $(-1,3)$，}
\step{所以 $-1$ 和 $3$ 是方程 $ax^2+bx+c=0$ 的实数解，}
\step{所以 $\begin{cases}-1+3=-\dfrac ba\\ -1\times3=\dfrac ca\\ a<0\end{cases}$，解得 $b=-2a$，$c=-3a$，且 $a<0$，}
\step{所以不等式 $cx^2+bx+a>0$ 可化为 $-3ax^2-2ax+a>0$，即 $3x^2+2x-1>0$，}
\step{解得 $x<-1$ 或 $x>\dfrac13$，所以该不等式的解集为 $(-\infty,-1)\cup\left(\dfrac13,+\infty\right)$.}
\solend
\fi

\item 若集合 $A=\{x\mid(a-1)x^2-2x+1=0\}$ 有且仅有 $2$ 个子集，则实数 $a$ 的值为\blankline[4.4em].

\ifanswer
\solbegin
\step{由题意得方程 $(a-1)x^2-2x+1=0$ 只有一个根，}
\stepm{①}{当 $a-1=0$ 时，$a=1$，此时方程化为 $-2x+1=0$，解得 $x=\dfrac12$，符合题意，}
\stepm{②}{当 $a-1\neq0$，即 $a\neq1$ 时，$\Delta=4-4(a-1)=0$，解得 $a=2$，}
\step{综上所述，实数 $a$ 的值为 $1$ 或 $2$.}
\solend
\fi

\item “$x+y\neq5$”是“$x\neq3$ 或 $y\neq2$”的\blankline[4.4em]条件.

\ifanswer
\solbegin
\step{分析充分性：}
\step{若 $x+y\neq5$ 成立，若 $x=3$，则 $y\neq2$，若 $y=2$，则 $x\neq3$，}
\step{所以一定能说明 $x\neq3$ 和 $y\neq2$ 至少有一个成立，即 $x\neq3$ 或 $y\neq2$ 成立，}
\step{充分性成立；}
\step{分析必要性：}
\step{若 $x\neq3$ 或 $y\neq2$，即 $x\neq3$ 和 $y\neq2$ 二者至少一个成立，}
\step{举例：$x=1,y=4$，此时 $x+y=5$，所以必要性不成立.}
\step{所以“$x+y\neq5$”是“$x\neq3$ 或 $y\neq2$”的充分不必要条件.}
\solend
\fi

\item 设 $P$ 为非空实数集满足：对任意给定的 $x,y\in P$（$x,y$ 可以相同），都有 $x+y\in P,x-y\in P,xy\in P$，则称 $P$ 为漂亮集. 关于漂亮集有下列结论：

①集合 $P=\{-2,-1,0,1,2\}$ 为漂亮集；②集合 $P=\{x\mid x=2n,n\in\Z\}$ 为漂亮集；

③若集合 $P_1,P_2$ 为漂亮集，则 $P_1\cup P_2$ 为漂亮集；④若集合 $P$ 为漂亮集，则一定有 $0\in P$.

其中所有正确结论的序号是\blankline[4.4em].

\ifanswer
\solbegin
\step{对于集合 $P=\{-2,-1,0,1,2\}$，取 $x=2,y=-2$，则 $x-y=2-(-2)=4\notin P$，}
\step{不满足漂亮集的定义，所以①错误.}
\step{对于集合 $P=\{x\mid x=2n,n\in Z\}$，设 $x=2n_1,y=2n_2,n_1,n_2\in Z$.}
\step{$x+y=2n_1+2n_2=2(n_1+n_2)$，因为 $n_1+n_2\in Z$，所以 $x+y\in P$.}
\step{$x-y=2n_1-2n_2=2(n_1-n_2)$，因为 $n_1-n_2\in Z$，所以 $x-y\in P$.}
\step{$xy=2n_1\times2n_2=4n_1n_2=2(2n_1n_2)$，因为 $2n_1n_2\in Z$，所以 $xy\in P$.}
\step{满足漂亮集的定义，所以②正确.}
\step{令 $P_1=\{x\mid x=2n,n\in Z\}$,$P_2=\{x\mid x=3n,n\in Z\}$，都是漂亮集.}
\step{$P_1\cup P_2$ 中，取 $x=2(x\in P_1)$,$y=3(y\in P_2)$,$x+y=5\notin P_1\cup P_2$，}
\step{不满足漂亮集的定义，所以③错误.}
\step{因为对任意 $x\in P$，有 $x-x=0\in P$，所以若集合 $P$ 为漂亮集，}
\step{则一定有 $0\in P$，所以④正确.}
\step{故正确的是②④.}
\solend
\fi

\item 已知：集合 $A_p=\{(x,y)\mid y=3x^2-px+5p-65,x\in\R\}$,$M=\{x\mid1\leqslant x\leqslant100,x\in\Z\}$，若 $p\in M$，则所有 $A_p$ 的交集\blankline[4.4em].

\ifanswer
\solbegin
\step{由 $M=\{1,2,\cdots,100\}$，所以所有 $A_p$ 的交集就是这些集合中共同拥有的点.}
\step{若 $(x,y)$ 是所有集合的一个公共点，}
\step{那么它特别要同时属于 $p=1$ 和 $p=2$ 时对应的两个集合.}
\step{当 $p=1$ 时，$y=3x^2-x-60$；}
\step{当 $p=2$ 时，$y=3x^2-2x-55$；联立得 $x=5,y=10$.}
\step{因此，如果公共点存在，只可能是 $(5,10)$.}
\step{接下来验证这个点是否属于所有集合，}
\step{对任意 $p\in M$，把 $x=5$ 代入集合的关系式，}
\step{得 $3x^2-px+5p-65=75-5p+5p-65=10$，与 $y=10$ 相符，}
\step{所以 $(5,10)\in A_p$ 对每一个 $p$ 都成立，因此所有 $A_p$ 的交集为 $\{(5,10)\}$.}
\solend
\fi

\end{enumerate}

% 守卫须放在 \bigsec 之前，否则标题仍会落页末（见 高一40 的教训）
\needvspace{5cm}

\bigsec{二、选择题}

\begin{enumerate}
\setcounter{enumi}{12}

\item 已知 $a,b,c\in\R$ 且 $a>b$，则下列不等式一定成立的是\mbox{（\quad）}

\keepwithprev
A. $\dfrac1a<\dfrac1b$\qquad B. $a^2>b^2$

C. $a|c|>b|c|$\qquad D. $\dfrac a{c^2+1}>\dfrac b{c^2+1}$

\ifanswer
\solbegin
\step{$A$：当 $a>0>b$ 时，$\dfrac1a>\dfrac1b$，故 $A$ 错误；}
\step{$B$：当 $a=-1,b=-2$ 时，满足 $a>b$，但 $a^2>b^2$ 不成立，故 $B$ 错误；}
\step{$C$：当 $c=0$ 时，$a|c|=b|c|$，故 $C$ 错误；}
\step{$D$：由 $a>b,c^2+1>0$，得 $\dfrac a{c^2+1}>\dfrac b{c^2+1}$，故 $D$ 正确；}
\step{故选 $D$.}
\solend
\fi

\item 已知集合 $A=\{a^2,0,-1\}$，$B=\{a,b,0\}$，若 $A=B$，则 $(ab)^{2026}$ 的值为\mbox{（\quad）}

\keepwithprev
A. $0$\qquad B. $-1$

C. $1$\qquad D. $\pm1$

\ifanswer
\solbegin
\step{因为集合 $A=\{a^2,0,-1\}$，$B=\{a,b,0\}$，$A=B$，所以 $\begin{cases}a=-1\\ a^2=b\end{cases}$，}
\step{所以 $a=-1$，$b=1$，所以 $(ab)^{2026}=1$，故选 $C$.}
\solend
\fi

\item 设全集 $U=\R$，$A=\{x\mid x<-1\text{ 或 }x>2\}$,$B=\{y\mid y=x^2+a\}$，若 $\overline A\subseteq B$，则实数 $a$ 的取值范围是\mbox{（\quad）}

\keepwithprev
A. $a\leqslant-1$\qquad B. $a>2$

C. $a<-1$\qquad D. $a\geqslant-2$

\ifanswer
\solbegin
\step{因为 $A=\{x\mid x<-1\text{ 或 }x>2\}$，则 $\overline A=\{x\mid-1\leqslant x\leqslant2\}$，}
\step{又 $B=\{y\mid y=x^2+a\}=\{y\mid y\geqslant a\}$，且 $\overline A\subseteq B$，所以 $a\leqslant-1$，故选 $A$.}
\solend
\fi

\item 设 $P$ 和 $Q$ 是两个集合，定义集合 $P-Q=\{x\mid x\in P$，且 $x\notin Q\}$，如果 $Q=\{x\mid(1-x)(x-3)>0\}$，$P=\left\{x\left|\dfrac x{2-x}>0\right.\right\}$，那么 $P-Q$ 等于\mbox{（\quad）}

\keepwithprev
A. $\{x\mid0<x<1\}$\qquad B. $\{x\mid1\leqslant x<2\}$

C. $\{x\mid0<x\leqslant1\}$\qquad D. $\{x\mid2\leqslant x<3\}$

\ifanswer
\solbegin
\step{因为 $P=\left\{x\left|\dfrac x{2-x}>0\right.\right\}=\{x\mid x(x-2)<0\}=\{x\mid0<x<2\}$，}
\step{$Q=\{x\mid(1-x)(x-3)>0\}=\{x\mid(x-3)(x-1)<0\}=\{x\mid1<x<3\}$，}
\step{所以 $P-Q=\{x\mid0<x\leqslant1\}$，故选 $C$.}
\solend
\fi

\end{enumerate}

\bigsec{三、解答题}

\begin{enumerate}
\setcounter{enumi}{16}

\item 已知集合 $A=\{x\mid x^2-9\geqslant0\}$，$B=\{x\mid|x-4|<2\}$，$C=\left\{x\left|\dfrac{x-8}{x+2}<0\right.\right\}$.

（1）求 $A\cup C$；

（2）若全集 $U=R$，求 $\overline A\cap B$.

\ifanswer
\solbegin
\stepm{(1)}{集合 $A=\{x\mid x^2-9\geqslant0\}=\{x\mid x\geqslant3\text{ 或 }x\leqslant-3\}$，}
\step{$B=\{x\mid|x-4|<2\}=\{x\mid2<x<6\}$，$C=\left\{x\left|\dfrac{x-8}{x+2}<0\right.\right\}=\{x\mid-2<x<8\}$，}
\step{所以 $A\cup C=\{x\mid x\leqslant-3\text{ 或 }x>-2\}$；}
\stepm{(2)}{全集 $U=R$，集合 $A=\{x\mid x^2-9\geqslant0\}=\{x\mid x\geqslant3\text{ 或 }x\leqslant-3\}$，}
\step{$B=\{x\mid|x-4|<2\}=\{x\mid2<x<6\}$，}
\step{所以 $\overline A=\{x\mid-3<x<3\}$，所以 $\overline A\cap B=\{x\mid2<x<3\}$.}
\solend
\fi

\ansspace[6]

\item 已知集合 $A=\left\{x\left|\dfrac{2x-1}{x+1}<1,x\in\R\right.\right\}$，集合 $B=\{x\mid(x-a-1)(x-a+1)\leqslant0\}$.

（1）求集合 $A$；

（2）若 $B\subseteq A$，求实数 $a$ 的取值范围.

\ifanswer
\solbegin
\stepm{(1)}{由 $\dfrac{2x-1}{x+1}<1$，得 $\dfrac{x-2}{x+1}<0$ 所以 $(x-2)(x+1)<0$，解得 $-1<x<2$，}
\step{即 $A=(-1,2)$；}
\stepm{(2)}{由 $(x-a-1)(x-a+1)\leqslant0$，得 $a-1\leqslant x\leqslant a+1$，}
\step{所以 $B=\{x\mid a-1\leqslant x\leqslant a+1\}$，}
\step{又 $B\subseteq A$，且 $a+1>a-1$ 即 $B\neq\varnothing$，所以 $\begin{cases}a-1>-1\\ a+1<2\end{cases}$，所以 $0<a<1$，}
\step{所以实数 $a$ 的取值范围是 $\{a\mid0<a<1\}$.}
\solend
\fi

\ansspace[6]

\item 已知集合 $A=\{x\mid x^2-4x+3=0\}$，$B=\{x\mid x^2-ax+3=0\}$.

（1）若 $A\cup B=B$，求实数 $a$ 的值；

（2）若 $A\cap B=B$，求实数 $a$ 的取值范围.

\ifanswer
\solbegin
\stepm{(1)}{集合 $A=\{x\mid x^2-4x+3=0\}=\{1,3\}$，$B=\{x\mid x^2-ax+3=0\}$.}
\step{$A\cup B=B$，所以 $1\in B$，$3\in B$，所以 $\begin{cases}1-a+3=0\\ 9-3a+3=0\end{cases}$，解得实数 $a=4$；}
\stepm{(2)}{因为 $A\cap B=B$，所以 $B\subseteq A$，}
\step{当 $B=\varnothing$ 时，$\Delta=(-a)^2-12<0$，解得 $-2\sqrt3<a<2\sqrt3$，}
\step{当 $B$ 是单元素集时，$\Delta=a^2-12=0$，解得 $a=2\sqrt3$，此时 $B=\{\sqrt3\}$，}
\step{不合题意；}
\step{当 $B=\{1,3\}$ 时，由 (1) 得 $a=4$，}
\step{综上，实数 $a$ 的取值范围是 $(-2\sqrt3,2\sqrt3)\cup\{4\}$.}
\solend
\fi

\ansspace[6]

\item 已知全集为 $R$，实数 $a$、$b$ 满足 $a\neq b$，集合 $M=\{x\mid x^2-3x-4\leqslant0\}$，集合 $A=\{x\mid(x+a)(x+b)>0\}$，集合 $B=\{x\mid x^2+(a-2)x-2a>0\}$.

（1）若 $\overline A=M$，求 $a$、$b$ 的值；

（2）若 $a>b>-1$，求 $A\cap B$；

（3）若 $a^2+a\in\overline B$，求 $a$ 的取值范围.

\ifanswer
\solbegin
\stepm{(1)}{$M=\{x\mid x^2-3x-4\leqslant0\}=\{x\mid-1\leqslant x\leqslant4\}$，$A=\{x\mid(x+a)(x+b)\leqslant0\}$，}
\step{因为 $\overline A=M$，所以 $a=1$，$b=-4$ 或 $a=-4$，$b=1$.}
\stepm{(2)}{因为 $a>b>-1$，所以 $-a<-b<1$，所以 $A=\{x\mid x<-a\text{ 或 }x>-b\}$，}
\step{因为 $B=\{x\mid x^2+(a-2)x-2a>0\}=\{x\mid x<-a\text{ 或 }x>2\}$，}
\step{所以 $A\cap B=\{x\mid x<-a\text{ 或 }x>2\}$，}
\stepm{(3)}{$\overline B=\{x\mid(x-2)(x+a)\leqslant0\}$，}
\step{由 $a^2+a\in\overline B$，得 $(a^2+a-2)(a^2+2a)\leqslant0\Leftrightarrow a(a-1)(a-2)^2\leqslant0$，}
\step{解得 $a=2$ 或 $0\leqslant a\leqslant1$，故 $a$ 的取值范围是 $[0,1]\cup\{2\}$.}
\solend
\fi

\ansspace[8]

\item 对平面直角坐标系的第一象限内的任意两点作如下定义：若 $\dfrac ab>\dfrac cd$，那么称点 $(a,b)$ 是点 $(c,d)$ 的“上位点”，同时点 $(c,d)$ 是点 $(a,b)$ 的“下位点”.

（1）试写出点 $(2,4)$ 的一个“上位点”坐标；

（2）已知点 $(a,b)$ 是点 $(c,d)$ 的“上位点”，判断点 $P(a+c,b+d)$ 是点 $(c,d)$ 的“上位点”还是“下位点”，并证明你的结论；

（3）设正整数 $n$ 满足以下条件：对集合 $\{t\mid0<t<2020,t\in\Z\}$ 内的任意元素 $m$，总存在正整数 $k=2m+1$，使得点 $(n,k)$ 既是点 $(2025,m)$ 的“下位点”，又是点 $(2026,m+1)$ 的“上位点”，问这样的正整数 $n$ 是否存在？若存在，求正整数 $n$ 的值；若不存在，请说明理由.

\ifanswer
\solbegin
\stepm{(1)}{由定义，点 $(a,b)$ 是点 $(2,4)$ 的“上位点”，}
\step{当且仅当 $\dfrac ab>\dfrac24=\dfrac12$ 且 $a>0,b>0$.}
\step{取 $(a,b)=(1,1)$，因为 $\dfrac11>\dfrac12$，}
\step{所以点 $(1,1)$ 是点 $(2,4)$ 的一个“上位点”.（答案不唯一）}
\stepm{(2)}{点 $P(a+c,b+d)$ 是点 $(c,d)$ 的“上位点”.}
\step{$(a,b)$ 是点 $(c,d)$ 的“上位点”，所以 $\dfrac ab>\dfrac cd$.}
\step{因为 $a,b,c,d>0$，两边同乘 $bd$ 得 $ad>bc$，即 $ad-bc>0$.}
\step{作差比较 $\dfrac{a+c}{b+d}$ 与 $\dfrac cd$：$\dfrac{a+c}{b+d}-\dfrac cd=\dfrac{d(a+c)-c(b+d)}{d(b+d)}=\dfrac{ad-bc}{d(b+d)}>0$，}
\step{所以点 $P(a+c,b+d)$ 是点 $(c,d)$ 的“上位点”.}
\stepm{(3)}{由定义，点 $(n,k)$ 是点 $(2025,m)$ 的下位点即 $\dfrac{2025}m>\dfrac nk$；}
\step{点 $(n,k)$ 是点 $(2026,m+1)$ 的上位点即 $\dfrac nk>\dfrac{2026}{m+1}$.}
\step{代入 $k=2m+1$，得对任意 $m\in\{1,2,\cdots,2019\}$，}
\step{都有 $4052-\dfrac{2026}{m+1}<n<4050+\dfrac{2025}m$.}
\step{右边 $4050+\dfrac{2025}m$ 随 $m$ 增大而减小，故最近处在 $m=2019$，}
\step{得 $n<4050+\dfrac{2025}{2019}<4051.01$，所以 $n\leqslant4051$.}
\step{左边 $4052-\dfrac{2026}{m+1}$ 随 $m$ 增大而增大，故最近处也在 $m=2019$，}
\step{得 $n>4052-\dfrac{2026}{2020}>4050.99$，所以 $n\geqslant4051$.}
\step{综上，$n=4051$.}
\step{检验：当 $n=4051$ 时，对任意 $1\leqslant m\leqslant2019,k=2m+1$ 是正整数，}
\step{且 $\dfrac{4051}{2m+1}<\dfrac{2025}m$,$4051m<4050m+2025,m<2025$，成立；}
\step{$\dfrac{4051}{2m+1}>\dfrac{2026}{m+1}$,$4051(m+1)>4052m+2026,m<2025$，成立.}
\step{又各点横纵坐标均为正，都在第一象限内，符合定义，}
\step{所以满足条件的正整数 $n$ 存在，其值为 $n=4051$.}
\solend
\fi

\ansspace*[10]

\end{enumerate}

\end{document}
"
%
% 原始材料：微信公众号“魔都数学加油站”文章，作者破晓，连载“27学年高一59”；
% 原文标题“27学年高一59——2026-2027年上海市华实高中高一上10月月考”；
% 原文链接：https://mp.weixin.qq.com/s/XzyBCmrvpRuMe-NuNk7gkA。
% 页面标注发布时间 2026-10-13 21:45（上海）。
% 正文为 9 张 PNG 页：第 01、02、03、05、06、07、08、09 页 4762×6735，第 04 页 2560×3620
% （**同一篇各页分辨率不同**），均按页序读取。9 页全为试卷正文页（卷首标题在第 1 页页顶，
% 卷末解答题第 21 题解析止于第 9 页），无封面页、目录页、二维码或宣传页。
% 粉色斜排“魔都数学加油站”为卷面水印，与公众号页脚一并未录入。
%
% 卷首标题：照卷面印作“2026--2027 年华实高中高一上 10 月月考”（公众号标题中的“上海市”
% 卷面标题行没有，本稿照卷面），年份区间按全稿惯例排 en dash。原卷只有一行标题、无副标题。
%
% 篇幅与题号：原卷自第 3 题起编号（第 1、2 题不在本卷内）。填空题 3—12（10 题）、
% 选择题 13—16（4 题）、解答题 17—21（5 题），共 19 题。本稿沿用原节与原题号，
% 只改排为全稿统一的大节样式（原卷小节标题为“一、填空题”“二、选择题”“三、解答题”）。
%
% 讲评版：原卷每题下直接印【答案】或【解析】。**只印【答案】的 1 题**（第 3 题），本稿只给 \ans{}；
% 其余 18 题都印【解析】，按全稿口径这类题不另提 \ans{}——第 13、14、15、16 题的结论写在解析末句
% （“故选 D”等），亦照录在解析内。本卷**没有**印在题干括号内的答案。
%
% 副标题：原卷未印副标题。本稿按“知识模块”口径标作“集合与常用逻辑用语、不等式”：
% 19 题中集合与常用逻辑用语 13 题（第 3、9、10、11、12、14、15、16、17、18、19、20、21 题），
% 不等式 10 题（含等式与不等式性质、一元二次不等式、恒成立、基本应用：第 4、5、6、7、8、
% 13、17、18、19、20 题），两类均远超三分之一，故并标。口径同上。
%
% 与原卷的排版差异：
%   ① 原卷每题下直接印【答案】/【解析】，本稿按全稿惯例将答案与解析置于答案版，试卷版保持干净；
%   ② 原卷未印副标题，本稿按“知识模块”口径补出；
%   ③ 原卷填空横线约两个字宽，本稿按全稿惯例统一排 \blankline[4.4em]；
%   ④ 原卷实数集排粗体 $\mathbf{R}$（第 12、13、20 题等），整数集排粗体 $\mathbf{Z}$（第 11、12 题），
%      自然数集排斜体 $N$ 或粗体 $\mathbf{Z}$ 处按本稿惯例改用 $\R$、$\Z$、$\N$；
%      空集记号统一为 \varnothing；
%   ⑤ 原卷在数学式内多处用半角逗号（如 $x_1,x_2$、$a,b\in R$），均照录未强改；
%      句末句点按全稿惯例排西文句点；“漂亮集”“上位点”“下位点”等自定义名词照录；
%   ⑥ 第 17—21 题的子问原卷各自另起一行，本稿按全稿惯例用 enumerate 的 (1)(2)(3)；
%      解析中的子问标号一律用 \stepm{(1)}{…}；第 9、11 题解析中的①、②（③④）分情形或分结论，
%      按全稿惯例作 \stepm{①}{…}、\stepm{②}{…}（第 11 题解析末“故正确的是②④”照录为行内）；
%   ⑦ 第 12 题题干里的“$A_p$”下标原卷作斜体 $p$，照录；
%   ⑧ 每道解答题原卷未留作答空白，本稿按全稿惯例加 \ansspace。
%
% 原卷笔误/存疑（均照录未改，正文对应处另有 % 注）：本卷首轮逐页转录，暂未发现影响答案的笔误。
%
% 插图：无（全卷 9 页无任何图形、无“如图”字样）。
%
\documentclass[a4paper,11pt]{article}
\usepackage{common}

\begin{document}

\examtitlest[3]{2026--2027 年华实高中高一上 10 月月考}{集合与常用逻辑用语、不等式}

\bigsec{一、填空题}

\begin{enumerate}
\setcounter{enumi}{2}

\item “$x\geqslant1$ 且 $y\geqslant1$”的否定形式为\blankline[4.4em].

\ifanswer
\ans{$x<1$ 或 $y<1$}
\fi

\item 已知方程 $2x^2+4x-3=0$ 的两个根为 $x_1,x_2$，则 $\dfrac1{x_1}+\dfrac1{x_2}=$\blankline[4.4em].

\ifanswer
\solbegin
\step{利用韦达定理得 $\begin{cases}x_1+x_2=-2\\ x_1\cdot x_2=-\dfrac32\end{cases}$，则 $\dfrac1{x_1}+\dfrac1{x_2}=\dfrac{x_1+x_2}{x_1\cdot x_2}=\dfrac43$.}
\solend
\fi

\item 设 $a>b>0$，比较大小：$\dfrac1{a-b}$\blankline[2.4em]$\dfrac1a$.

\ifanswer
\solbegin
\step{由 $a>b>0$ 得 $a>0$ 且 $a-b>0$，所以 $\dfrac1{a-b}-\dfrac1a=\dfrac b{a(a-b)}>0$，}
\step{故 $\dfrac1{a-b}>\dfrac1a$.}
\solend
\fi

\item 若不等式 $2ax^2-ax+1>0$ 的解集为 $R$，则实数 $a$ 的取值范围是\blankline[4.4em].

\ifanswer
\solbegin
\step{当 $a\neq0$ 时，由题意得 $\begin{cases}a>0\\ \Delta<0\end{cases}$，即 $\begin{cases}a>0\\ a^2-8a<0\end{cases}$，解得 $0<a<8$.}
\step{当 $a=0$ 时，恒有 $1>0$，不等式也成立.}
\step{综上所述，$a$ 的取值范围是 $[0,8)$.}
\solend
\fi

\item 集合 $A=\{x\mid m-1\leqslant x\leqslant m+7\}$，集合 $B=\{x\mid x^2-3x-10\geqslant0\}$，若 $A\cup B=R$，则实数 $m$ 的取值范围为\blankline[4.4em].

\ifanswer
\solbegin
\step{$A=\{x\mid m-1\leqslant x\leqslant m+7\}$，$B=\{x\mid x^2-3x-10\geqslant0\}=\{x\mid x\leqslant-2\text{ 或 }x\geqslant5\}$，}
\step{若 $A\cup B=R$，则 $\begin{cases}m-1\leqslant-2\\ m+7\geqslant5\end{cases}$，解得 $-2\leqslant m\leqslant-1$，}
\step{所以实数 $m$ 的取值范围为 $[-2,-1]$.}
\solend
\fi

\item 若关于 $x$ 的不等式 $ax^2+bx+c>0$ 的解集为 $(-1,3)$，则关于 $x$ 的不等式 $cx^2+bx+a>0$ 的解集为\blankline[4.4em].

\ifanswer
\solbegin
\step{因为不等式 $ax^2+bx+c>0$ 的解集为 $(-1,3)$，}
\step{所以 $-1$ 和 $3$ 是方程 $ax^2+bx+c=0$ 的实数解，}
\step{所以 $\begin{cases}-1+3=-\dfrac ba\\ -1\times3=\dfrac ca\\ a<0\end{cases}$，解得 $b=-2a$，$c=-3a$，且 $a<0$，}
\step{所以不等式 $cx^2+bx+a>0$ 可化为 $-3ax^2-2ax+a>0$，即 $3x^2+2x-1>0$，}
\step{解得 $x<-1$ 或 $x>\dfrac13$，所以该不等式的解集为 $(-\infty,-1)\cup\left(\dfrac13,+\infty\right)$.}
\solend
\fi

\item 若集合 $A=\{x\mid(a-1)x^2-2x+1=0\}$ 有且仅有 $2$ 个子集，则实数 $a$ 的值为\blankline[4.4em].

\ifanswer
\solbegin
\step{由题意得方程 $(a-1)x^2-2x+1=0$ 只有一个根，}
\stepm{①}{当 $a-1=0$ 时，$a=1$，此时方程化为 $-2x+1=0$，解得 $x=\dfrac12$，符合题意，}
\stepm{②}{当 $a-1\neq0$，即 $a\neq1$ 时，$\Delta=4-4(a-1)=0$，解得 $a=2$，}
\step{综上所述，实数 $a$ 的值为 $1$ 或 $2$.}
\solend
\fi

\item “$x+y\neq5$”是“$x\neq3$ 或 $y\neq2$”的\blankline[4.4em]条件.

\ifanswer
\solbegin
\step{分析充分性：}
\step{若 $x+y\neq5$ 成立，若 $x=3$，则 $y\neq2$，若 $y=2$，则 $x\neq3$，}
\step{所以一定能说明 $x\neq3$ 和 $y\neq2$ 至少有一个成立，即 $x\neq3$ 或 $y\neq2$ 成立，}
\step{充分性成立；}
\step{分析必要性：}
\step{若 $x\neq3$ 或 $y\neq2$，即 $x\neq3$ 和 $y\neq2$ 二者至少一个成立，}
\step{举例：$x=1,y=4$，此时 $x+y=5$，所以必要性不成立.}
\step{所以“$x+y\neq5$”是“$x\neq3$ 或 $y\neq2$”的充分不必要条件.}
\solend
\fi

\item 设 $P$ 为非空实数集满足：对任意给定的 $x,y\in P$（$x,y$ 可以相同），都有 $x+y\in P,x-y\in P,xy\in P$，则称 $P$ 为漂亮集. 关于漂亮集有下列结论：

①集合 $P=\{-2,-1,0,1,2\}$ 为漂亮集；②集合 $P=\{x\mid x=2n,n\in\Z\}$ 为漂亮集；

③若集合 $P_1,P_2$ 为漂亮集，则 $P_1\cup P_2$ 为漂亮集；④若集合 $P$ 为漂亮集，则一定有 $0\in P$.

其中所有正确结论的序号是\blankline[4.4em].

\ifanswer
\solbegin
\step{对于集合 $P=\{-2,-1,0,1,2\}$，取 $x=2,y=-2$，则 $x-y=2-(-2)=4\notin P$，}
\step{不满足漂亮集的定义，所以①错误.}
\step{对于集合 $P=\{x\mid x=2n,n\in Z\}$，设 $x=2n_1,y=2n_2,n_1,n_2\in Z$.}
\step{$x+y=2n_1+2n_2=2(n_1+n_2)$，因为 $n_1+n_2\in Z$，所以 $x+y\in P$.}
\step{$x-y=2n_1-2n_2=2(n_1-n_2)$，因为 $n_1-n_2\in Z$，所以 $x-y\in P$.}
\step{$xy=2n_1\times2n_2=4n_1n_2=2(2n_1n_2)$，因为 $2n_1n_2\in Z$，所以 $xy\in P$.}
\step{满足漂亮集的定义，所以②正确.}
\step{令 $P_1=\{x\mid x=2n,n\in Z\}$,$P_2=\{x\mid x=3n,n\in Z\}$，都是漂亮集.}
\step{$P_1\cup P_2$ 中，取 $x=2(x\in P_1)$,$y=3(y\in P_2)$,$x+y=5\notin P_1\cup P_2$，}
\step{不满足漂亮集的定义，所以③错误.}
\step{因为对任意 $x\in P$，有 $x-x=0\in P$，所以若集合 $P$ 为漂亮集，}
\step{则一定有 $0\in P$，所以④正确.}
\step{故正确的是②④.}
\solend
\fi

\item 已知：集合 $A_p=\{(x,y)\mid y=3x^2-px+5p-65,x\in\R\}$,$M=\{x\mid1\leqslant x\leqslant100,x\in\Z\}$，若 $p\in M$，则所有 $A_p$ 的交集\blankline[4.4em].

\ifanswer
\solbegin
\step{由 $M=\{1,2,\cdots,100\}$，所以所有 $A_p$ 的交集就是这些集合中共同拥有的点.}
\step{若 $(x,y)$ 是所有集合的一个公共点，}
\step{那么它特别要同时属于 $p=1$ 和 $p=2$ 时对应的两个集合.}
\step{当 $p=1$ 时，$y=3x^2-x-60$；}
\step{当 $p=2$ 时，$y=3x^2-2x-55$；联立得 $x=5,y=10$.}
\step{因此，如果公共点存在，只可能是 $(5,10)$.}
\step{接下来验证这个点是否属于所有集合，}
\step{对任意 $p\in M$，把 $x=5$ 代入集合的关系式，}
\step{得 $3x^2-px+5p-65=75-5p+5p-65=10$，与 $y=10$ 相符，}
\step{所以 $(5,10)\in A_p$ 对每一个 $p$ 都成立，因此所有 $A_p$ 的交集为 $\{(5,10)\}$.}
\solend
\fi

\end{enumerate}

% 守卫须放在 \bigsec 之前，否则标题仍会落页末（见 高一40 的教训）
\needvspace{5cm}

\bigsec{二、选择题}

\begin{enumerate}
\setcounter{enumi}{12}

\item 已知 $a,b,c\in\R$ 且 $a>b$，则下列不等式一定成立的是\mbox{（\quad）}

\keepwithprev
A. $\dfrac1a<\dfrac1b$\qquad B. $a^2>b^2$

C. $a|c|>b|c|$\qquad D. $\dfrac a{c^2+1}>\dfrac b{c^2+1}$

\ifanswer
\solbegin
\step{$A$：当 $a>0>b$ 时，$\dfrac1a>\dfrac1b$，故 $A$ 错误；}
\step{$B$：当 $a=-1,b=-2$ 时，满足 $a>b$，但 $a^2>b^2$ 不成立，故 $B$ 错误；}
\step{$C$：当 $c=0$ 时，$a|c|=b|c|$，故 $C$ 错误；}
\step{$D$：由 $a>b,c^2+1>0$，得 $\dfrac a{c^2+1}>\dfrac b{c^2+1}$，故 $D$ 正确；}
\step{故选 $D$.}
\solend
\fi

\item 已知集合 $A=\{a^2,0,-1\}$，$B=\{a,b,0\}$，若 $A=B$，则 $(ab)^{2026}$ 的值为\mbox{（\quad）}

\keepwithprev
A. $0$\qquad B. $-1$

C. $1$\qquad D. $\pm1$

\ifanswer
\solbegin
\step{因为集合 $A=\{a^2,0,-1\}$，$B=\{a,b,0\}$，$A=B$，所以 $\begin{cases}a=-1\\ a^2=b\end{cases}$，}
\step{所以 $a=-1$，$b=1$，所以 $(ab)^{2026}=1$，故选 $C$.}
\solend
\fi

\item 设全集 $U=\R$，$A=\{x\mid x<-1\text{ 或 }x>2\}$,$B=\{y\mid y=x^2+a\}$，若 $\overline A\subseteq B$，则实数 $a$ 的取值范围是\mbox{（\quad）}

\keepwithprev
A. $a\leqslant-1$\qquad B. $a>2$

C. $a<-1$\qquad D. $a\geqslant-2$

\ifanswer
\solbegin
\step{因为 $A=\{x\mid x<-1\text{ 或 }x>2\}$，则 $\overline A=\{x\mid-1\leqslant x\leqslant2\}$，}
\step{又 $B=\{y\mid y=x^2+a\}=\{y\mid y\geqslant a\}$，且 $\overline A\subseteq B$，所以 $a\leqslant-1$，故选 $A$.}
\solend
\fi

\item 设 $P$ 和 $Q$ 是两个集合，定义集合 $P-Q=\{x\mid x\in P$，且 $x\notin Q\}$，如果 $Q=\{x\mid(1-x)(x-3)>0\}$，$P=\left\{x\left|\dfrac x{2-x}>0\right.\right\}$，那么 $P-Q$ 等于\mbox{（\quad）}

\keepwithprev
A. $\{x\mid0<x<1\}$\qquad B. $\{x\mid1\leqslant x<2\}$

C. $\{x\mid0<x\leqslant1\}$\qquad D. $\{x\mid2\leqslant x<3\}$

\ifanswer
\solbegin
\step{因为 $P=\left\{x\left|\dfrac x{2-x}>0\right.\right\}=\{x\mid x(x-2)<0\}=\{x\mid0<x<2\}$，}
\step{$Q=\{x\mid(1-x)(x-3)>0\}=\{x\mid(x-3)(x-1)<0\}=\{x\mid1<x<3\}$，}
\step{所以 $P-Q=\{x\mid0<x\leqslant1\}$，故选 $C$.}
\solend
\fi

\end{enumerate}

\bigsec{三、解答题}

\begin{enumerate}
\setcounter{enumi}{16}

\item 已知集合 $A=\{x\mid x^2-9\geqslant0\}$，$B=\{x\mid|x-4|<2\}$，$C=\left\{x\left|\dfrac{x-8}{x+2}<0\right.\right\}$.

（1）求 $A\cup C$；

（2）若全集 $U=R$，求 $\overline A\cap B$.

\ifanswer
\solbegin
\stepm{(1)}{集合 $A=\{x\mid x^2-9\geqslant0\}=\{x\mid x\geqslant3\text{ 或 }x\leqslant-3\}$，}
\step{$B=\{x\mid|x-4|<2\}=\{x\mid2<x<6\}$，$C=\left\{x\left|\dfrac{x-8}{x+2}<0\right.\right\}=\{x\mid-2<x<8\}$，}
\step{所以 $A\cup C=\{x\mid x\leqslant-3\text{ 或 }x>-2\}$；}
\stepm{(2)}{全集 $U=R$，集合 $A=\{x\mid x^2-9\geqslant0\}=\{x\mid x\geqslant3\text{ 或 }x\leqslant-3\}$，}
\step{$B=\{x\mid|x-4|<2\}=\{x\mid2<x<6\}$，}
\step{所以 $\overline A=\{x\mid-3<x<3\}$，所以 $\overline A\cap B=\{x\mid2<x<3\}$.}
\solend
\fi

\ansspace[6]

\item 已知集合 $A=\left\{x\left|\dfrac{2x-1}{x+1}<1,x\in\R\right.\right\}$，集合 $B=\{x\mid(x-a-1)(x-a+1)\leqslant0\}$.

（1）求集合 $A$；

（2）若 $B\subseteq A$，求实数 $a$ 的取值范围.

\ifanswer
\solbegin
\stepm{(1)}{由 $\dfrac{2x-1}{x+1}<1$，得 $\dfrac{x-2}{x+1}<0$ 所以 $(x-2)(x+1)<0$，解得 $-1<x<2$，}
\step{即 $A=(-1,2)$；}
\stepm{(2)}{由 $(x-a-1)(x-a+1)\leqslant0$，得 $a-1\leqslant x\leqslant a+1$，}
\step{所以 $B=\{x\mid a-1\leqslant x\leqslant a+1\}$，}
\step{又 $B\subseteq A$，且 $a+1>a-1$ 即 $B\neq\varnothing$，所以 $\begin{cases}a-1>-1\\ a+1<2\end{cases}$，所以 $0<a<1$，}
\step{所以实数 $a$ 的取值范围是 $\{a\mid0<a<1\}$.}
\solend
\fi

\ansspace[6]

\item 已知集合 $A=\{x\mid x^2-4x+3=0\}$，$B=\{x\mid x^2-ax+3=0\}$.

（1）若 $A\cup B=B$，求实数 $a$ 的值；

（2）若 $A\cap B=B$，求实数 $a$ 的取值范围.

\ifanswer
\solbegin
\stepm{(1)}{集合 $A=\{x\mid x^2-4x+3=0\}=\{1,3\}$，$B=\{x\mid x^2-ax+3=0\}$.}
\step{$A\cup B=B$，所以 $1\in B$，$3\in B$，所以 $\begin{cases}1-a+3=0\\ 9-3a+3=0\end{cases}$，解得实数 $a=4$；}
\stepm{(2)}{因为 $A\cap B=B$，所以 $B\subseteq A$，}
\step{当 $B=\varnothing$ 时，$\Delta=(-a)^2-12<0$，解得 $-2\sqrt3<a<2\sqrt3$，}
\step{当 $B$ 是单元素集时，$\Delta=a^2-12=0$，解得 $a=2\sqrt3$，此时 $B=\{\sqrt3\}$，}
\step{不合题意；}
\step{当 $B=\{1,3\}$ 时，由 (1) 得 $a=4$，}
\step{综上，实数 $a$ 的取值范围是 $(-2\sqrt3,2\sqrt3)\cup\{4\}$.}
\solend
\fi

\ansspace[6]

\item 已知全集为 $R$，实数 $a$、$b$ 满足 $a\neq b$，集合 $M=\{x\mid x^2-3x-4\leqslant0\}$，集合 $A=\{x\mid(x+a)(x+b)>0\}$，集合 $B=\{x\mid x^2+(a-2)x-2a>0\}$.

（1）若 $\overline A=M$，求 $a$、$b$ 的值；

（2）若 $a>b>-1$，求 $A\cap B$；

（3）若 $a^2+a\in\overline B$，求 $a$ 的取值范围.

\ifanswer
\solbegin
\stepm{(1)}{$M=\{x\mid x^2-3x-4\leqslant0\}=\{x\mid-1\leqslant x\leqslant4\}$，$A=\{x\mid(x+a)(x+b)\leqslant0\}$，}
\step{因为 $\overline A=M$，所以 $a=1$，$b=-4$ 或 $a=-4$，$b=1$.}
\stepm{(2)}{因为 $a>b>-1$，所以 $-a<-b<1$，所以 $A=\{x\mid x<-a\text{ 或 }x>-b\}$，}
\step{因为 $B=\{x\mid x^2+(a-2)x-2a>0\}=\{x\mid x<-a\text{ 或 }x>2\}$，}
\step{所以 $A\cap B=\{x\mid x<-a\text{ 或 }x>2\}$，}
\stepm{(3)}{$\overline B=\{x\mid(x-2)(x+a)\leqslant0\}$，}
\step{由 $a^2+a\in\overline B$，得 $(a^2+a-2)(a^2+2a)\leqslant0\Leftrightarrow a(a-1)(a-2)^2\leqslant0$，}
\step{解得 $a=2$ 或 $0\leqslant a\leqslant1$，故 $a$ 的取值范围是 $[0,1]\cup\{2\}$.}
\solend
\fi

\ansspace[8]

\item 对平面直角坐标系的第一象限内的任意两点作如下定义：若 $\dfrac ab>\dfrac cd$，那么称点 $(a,b)$ 是点 $(c,d)$ 的“上位点”，同时点 $(c,d)$ 是点 $(a,b)$ 的“下位点”.

（1）试写出点 $(2,4)$ 的一个“上位点”坐标；

（2）已知点 $(a,b)$ 是点 $(c,d)$ 的“上位点”，判断点 $P(a+c,b+d)$ 是点 $(c,d)$ 的“上位点”还是“下位点”，并证明你的结论；

（3）设正整数 $n$ 满足以下条件：对集合 $\{t\mid0<t<2020,t\in\Z\}$ 内的任意元素 $m$，总存在正整数 $k=2m+1$，使得点 $(n,k)$ 既是点 $(2025,m)$ 的“下位点”，又是点 $(2026,m+1)$ 的“上位点”，问这样的正整数 $n$ 是否存在？若存在，求正整数 $n$ 的值；若不存在，请说明理由.

\ifanswer
\solbegin
\stepm{(1)}{由定义，点 $(a,b)$ 是点 $(2,4)$ 的“上位点”，}
\step{当且仅当 $\dfrac ab>\dfrac24=\dfrac12$ 且 $a>0,b>0$.}
\step{取 $(a,b)=(1,1)$，因为 $\dfrac11>\dfrac12$，}
\step{所以点 $(1,1)$ 是点 $(2,4)$ 的一个“上位点”.（答案不唯一）}
\stepm{(2)}{点 $P(a+c,b+d)$ 是点 $(c,d)$ 的“上位点”.}
\step{$(a,b)$ 是点 $(c,d)$ 的“上位点”，所以 $\dfrac ab>\dfrac cd$.}
\step{因为 $a,b,c,d>0$，两边同乘 $bd$ 得 $ad>bc$，即 $ad-bc>0$.}
\step{作差比较 $\dfrac{a+c}{b+d}$ 与 $\dfrac cd$：$\dfrac{a+c}{b+d}-\dfrac cd=\dfrac{d(a+c)-c(b+d)}{d(b+d)}=\dfrac{ad-bc}{d(b+d)}>0$，}
\step{所以点 $P(a+c,b+d)$ 是点 $(c,d)$ 的“上位点”.}
\stepm{(3)}{由定义，点 $(n,k)$ 是点 $(2025,m)$ 的下位点即 $\dfrac{2025}m>\dfrac nk$；}
\step{点 $(n,k)$ 是点 $(2026,m+1)$ 的上位点即 $\dfrac nk>\dfrac{2026}{m+1}$.}
\step{代入 $k=2m+1$，得对任意 $m\in\{1,2,\cdots,2019\}$，}
\step{都有 $4052-\dfrac{2026}{m+1}<n<4050+\dfrac{2025}m$.}
\step{右边 $4050+\dfrac{2025}m$ 随 $m$ 增大而减小，故最近处在 $m=2019$，}
\step{得 $n<4050+\dfrac{2025}{2019}<4051.01$，所以 $n\leqslant4051$.}
\step{左边 $4052-\dfrac{2026}{m+1}$ 随 $m$ 增大而增大，故最近处也在 $m=2019$，}
\step{得 $n>4052-\dfrac{2026}{2020}>4050.99$，所以 $n\geqslant4051$.}
\step{综上，$n=4051$.}
\step{检验：当 $n=4051$ 时，对任意 $1\leqslant m\leqslant2019,k=2m+1$ 是正整数，}
\step{且 $\dfrac{4051}{2m+1}<\dfrac{2025}m$,$4051m<4050m+2025,m<2025$，成立；}
\step{$\dfrac{4051}{2m+1}>\dfrac{2026}{m+1}$,$4051(m+1)>4052m+2026,m<2025$，成立.}
\step{又各点横纵坐标均为正，都在第一象限内，符合定义，}
\step{所以满足条件的正整数 $n$ 存在，其值为 $n=4051$.}
\solend
\fi

\ansspace*[10]

\end{enumerate}

\end{document}
"
%
% 原始材料：微信公众号“魔都数学加油站”文章，作者破晓，连载“27学年高一59”；
% 原文标题“27学年高一59——2026-2027年上海市华实高中高一上10月月考”；
% 原文链接：https://mp.weixin.qq.com/s/XzyBCmrvpRuMe-NuNk7gkA。
% 页面标注发布时间 2026-10-13 21:45（上海）。
% 正文为 9 张 PNG 页：第 01、02、03、05、06、07、08、09 页 4762×6735，第 04 页 2560×3620
% （**同一篇各页分辨率不同**），均按页序读取。9 页全为试卷正文页（卷首标题在第 1 页页顶，
% 卷末解答题第 21 题解析止于第 9 页），无封面页、目录页、二维码或宣传页。
% 粉色斜排“魔都数学加油站”为卷面水印，与公众号页脚一并未录入。
%
% 卷首标题：照卷面印作“2026--2027 年华实高中高一上 10 月月考”（公众号标题中的“上海市”
% 卷面标题行没有，本稿照卷面），年份区间按全稿惯例排 en dash。原卷只有一行标题、无副标题。
%
% 篇幅与题号：原卷自第 3 题起编号（第 1、2 题不在本卷内）。填空题 3—12（10 题）、
% 选择题 13—16（4 题）、解答题 17—21（5 题），共 19 题。本稿沿用原节与原题号，
% 只改排为全稿统一的大节样式（原卷小节标题为“一、填空题”“二、选择题”“三、解答题”）。
%
% 讲评版：原卷每题下直接印【答案】或【解析】。**只印【答案】的 1 题**（第 3 题），本稿只给 \ans{}；
% 其余 18 题都印【解析】，按全稿口径这类题不另提 \ans{}——第 13、14、15、16 题的结论写在解析末句
% （“故选 D”等），亦照录在解析内。本卷**没有**印在题干括号内的答案。
%
% 副标题：原卷未印副标题。本稿按“知识模块”口径标作“集合与常用逻辑用语、不等式”：
% 19 题中集合与常用逻辑用语 13 题（第 3、9、10、11、12、14、15、16、17、18、19、20、21 题），
% 不等式 10 题（含等式与不等式性质、一元二次不等式、恒成立、基本应用：第 4、5、6、7、8、
% 13、17、18、19、20 题），两类均远超三分之一，故并标。口径同上。
%
% 与原卷的排版差异：
%   ① 原卷每题下直接印【答案】/【解析】，本稿按全稿惯例将答案与解析置于答案版，试卷版保持干净；
%   ② 原卷未印副标题，本稿按“知识模块”口径补出；
%   ③ 原卷填空横线约两个字宽，本稿按全稿惯例统一排 \blankline[4.4em]；
%   ④ 原卷实数集排粗体 $\mathbf{R}$（第 12、13、20 题等），整数集排粗体 $\mathbf{Z}$（第 11、12 题），
%      自然数集排斜体 $N$ 或粗体 $\mathbf{Z}$ 处按本稿惯例改用 $\R$、$\Z$、$\N$；
%      空集记号统一为 \varnothing；
%   ⑤ 原卷在数学式内多处用半角逗号（如 $x_1,x_2$、$a,b\in R$），均照录未强改；
%      句末句点按全稿惯例排西文句点；“漂亮集”“上位点”“下位点”等自定义名词照录；
%   ⑥ 第 17—21 题的子问原卷各自另起一行，本稿按全稿惯例用 enumerate 的 (1)(2)(3)；
%      解析中的子问标号一律用 \stepm{(1)}{…}；第 9、11 题解析中的①、②（③④）分情形或分结论，
%      按全稿惯例作 \stepm{①}{…}、\stepm{②}{…}（第 11 题解析末“故正确的是②④”照录为行内）；
%   ⑦ 第 12 题题干里的“$A_p$”下标原卷作斜体 $p$，照录；
%   ⑧ 每道解答题原卷未留作答空白，本稿按全稿惯例加 \ansspace。
%
% 原卷笔误/存疑（均照录未改，正文对应处另有 % 注）：本卷首轮逐页转录，暂未发现影响答案的笔误。
%
% 插图：无（全卷 9 页无任何图形、无“如图”字样）。
%
\documentclass[a4paper,11pt]{article}
\usepackage{common}

\begin{document}

\examtitlest[3]{2026--2027 年华实高中高一上 10 月月考}{集合与常用逻辑用语、不等式}

\bigsec{一、填空题}

\begin{enumerate}
\setcounter{enumi}{2}

\item “$x\geqslant1$ 且 $y\geqslant1$”的否定形式为\blankline[4.4em].

\ifanswer
\ans{$x<1$ 或 $y<1$}
\fi

\item 已知方程 $2x^2+4x-3=0$ 的两个根为 $x_1,x_2$，则 $\dfrac1{x_1}+\dfrac1{x_2}=$\blankline[4.4em].

\ifanswer
\solbegin
\step{利用韦达定理得 $\begin{cases}x_1+x_2=-2\\ x_1\cdot x_2=-\dfrac32\end{cases}$，则 $\dfrac1{x_1}+\dfrac1{x_2}=\dfrac{x_1+x_2}{x_1\cdot x_2}=\dfrac43$.}
\solend
\fi

\item 设 $a>b>0$，比较大小：$\dfrac1{a-b}$\blankline[2.4em]$\dfrac1a$.

\ifanswer
\solbegin
\step{由 $a>b>0$ 得 $a>0$ 且 $a-b>0$，所以 $\dfrac1{a-b}-\dfrac1a=\dfrac b{a(a-b)}>0$，}
\step{故 $\dfrac1{a-b}>\dfrac1a$.}
\solend
\fi

\item 若不等式 $2ax^2-ax+1>0$ 的解集为 $R$，则实数 $a$ 的取值范围是\blankline[4.4em].

\ifanswer
\solbegin
\step{当 $a\neq0$ 时，由题意得 $\begin{cases}a>0\\ \Delta<0\end{cases}$，即 $\begin{cases}a>0\\ a^2-8a<0\end{cases}$，解得 $0<a<8$.}
\step{当 $a=0$ 时，恒有 $1>0$，不等式也成立.}
\step{综上所述，$a$ 的取值范围是 $[0,8)$.}
\solend
\fi

\item 集合 $A=\{x\mid m-1\leqslant x\leqslant m+7\}$，集合 $B=\{x\mid x^2-3x-10\geqslant0\}$，若 $A\cup B=R$，则实数 $m$ 的取值范围为\blankline[4.4em].

\ifanswer
\solbegin
\step{$A=\{x\mid m-1\leqslant x\leqslant m+7\}$，$B=\{x\mid x^2-3x-10\geqslant0\}=\{x\mid x\leqslant-2\text{ 或 }x\geqslant5\}$，}
\step{若 $A\cup B=R$，则 $\begin{cases}m-1\leqslant-2\\ m+7\geqslant5\end{cases}$，解得 $-2\leqslant m\leqslant-1$，}
\step{所以实数 $m$ 的取值范围为 $[-2,-1]$.}
\solend
\fi

\item 若关于 $x$ 的不等式 $ax^2+bx+c>0$ 的解集为 $(-1,3)$，则关于 $x$ 的不等式 $cx^2+bx+a>0$ 的解集为\blankline[4.4em].

\ifanswer
\solbegin
\step{因为不等式 $ax^2+bx+c>0$ 的解集为 $(-1,3)$，}
\step{所以 $-1$ 和 $3$ 是方程 $ax^2+bx+c=0$ 的实数解，}
\step{所以 $\begin{cases}-1+3=-\dfrac ba\\ -1\times3=\dfrac ca\\ a<0\end{cases}$，解得 $b=-2a$，$c=-3a$，且 $a<0$，}
\step{所以不等式 $cx^2+bx+a>0$ 可化为 $-3ax^2-2ax+a>0$，即 $3x^2+2x-1>0$，}
\step{解得 $x<-1$ 或 $x>\dfrac13$，所以该不等式的解集为 $(-\infty,-1)\cup\left(\dfrac13,+\infty\right)$.}
\solend
\fi

\item 若集合 $A=\{x\mid(a-1)x^2-2x+1=0\}$ 有且仅有 $2$ 个子集，则实数 $a$ 的值为\blankline[4.4em].

\ifanswer
\solbegin
\step{由题意得方程 $(a-1)x^2-2x+1=0$ 只有一个根，}
\stepm{①}{当 $a-1=0$ 时，$a=1$，此时方程化为 $-2x+1=0$，解得 $x=\dfrac12$，符合题意，}
\stepm{②}{当 $a-1\neq0$，即 $a\neq1$ 时，$\Delta=4-4(a-1)=0$，解得 $a=2$，}
\step{综上所述，实数 $a$ 的值为 $1$ 或 $2$.}
\solend
\fi

\item “$x+y\neq5$”是“$x\neq3$ 或 $y\neq2$”的\blankline[4.4em]条件.

\ifanswer
\solbegin
\step{分析充分性：}
\step{若 $x+y\neq5$ 成立，若 $x=3$，则 $y\neq2$，若 $y=2$，则 $x\neq3$，}
\step{所以一定能说明 $x\neq3$ 和 $y\neq2$ 至少有一个成立，即 $x\neq3$ 或 $y\neq2$ 成立，}
\step{充分性成立；}
\step{分析必要性：}
\step{若 $x\neq3$ 或 $y\neq2$，即 $x\neq3$ 和 $y\neq2$ 二者至少一个成立，}
\step{举例：$x=1,y=4$，此时 $x+y=5$，所以必要性不成立.}
\step{所以“$x+y\neq5$”是“$x\neq3$ 或 $y\neq2$”的充分不必要条件.}
\solend
\fi

\item 设 $P$ 为非空实数集满足：对任意给定的 $x,y\in P$（$x,y$ 可以相同），都有 $x+y\in P,x-y\in P,xy\in P$，则称 $P$ 为漂亮集. 关于漂亮集有下列结论：

①集合 $P=\{-2,-1,0,1,2\}$ 为漂亮集；②集合 $P=\{x\mid x=2n,n\in\Z\}$ 为漂亮集；

③若集合 $P_1,P_2$ 为漂亮集，则 $P_1\cup P_2$ 为漂亮集；④若集合 $P$ 为漂亮集，则一定有 $0\in P$.

其中所有正确结论的序号是\blankline[4.4em].

\ifanswer
\solbegin
\step{对于集合 $P=\{-2,-1,0,1,2\}$，取 $x=2,y=-2$，则 $x-y=2-(-2)=4\notin P$，}
\step{不满足漂亮集的定义，所以①错误.}
\step{对于集合 $P=\{x\mid x=2n,n\in Z\}$，设 $x=2n_1,y=2n_2,n_1,n_2\in Z$.}
\step{$x+y=2n_1+2n_2=2(n_1+n_2)$，因为 $n_1+n_2\in Z$，所以 $x+y\in P$.}
\step{$x-y=2n_1-2n_2=2(n_1-n_2)$，因为 $n_1-n_2\in Z$，所以 $x-y\in P$.}
\step{$xy=2n_1\times2n_2=4n_1n_2=2(2n_1n_2)$，因为 $2n_1n_2\in Z$，所以 $xy\in P$.}
\step{满足漂亮集的定义，所以②正确.}
\step{令 $P_1=\{x\mid x=2n,n\in Z\}$,$P_2=\{x\mid x=3n,n\in Z\}$，都是漂亮集.}
\step{$P_1\cup P_2$ 中，取 $x=2(x\in P_1)$,$y=3(y\in P_2)$,$x+y=5\notin P_1\cup P_2$，}
\step{不满足漂亮集的定义，所以③错误.}
\step{因为对任意 $x\in P$，有 $x-x=0\in P$，所以若集合 $P$ 为漂亮集，}
\step{则一定有 $0\in P$，所以④正确.}
\step{故正确的是②④.}
\solend
\fi

\item 已知：集合 $A_p=\{(x,y)\mid y=3x^2-px+5p-65,x\in\R\}$,$M=\{x\mid1\leqslant x\leqslant100,x\in\Z\}$，若 $p\in M$，则所有 $A_p$ 的交集\blankline[4.4em].

\ifanswer
\solbegin
\step{由 $M=\{1,2,\cdots,100\}$，所以所有 $A_p$ 的交集就是这些集合中共同拥有的点.}
\step{若 $(x,y)$ 是所有集合的一个公共点，}
\step{那么它特别要同时属于 $p=1$ 和 $p=2$ 时对应的两个集合.}
\step{当 $p=1$ 时，$y=3x^2-x-60$；}
\step{当 $p=2$ 时，$y=3x^2-2x-55$；联立得 $x=5,y=10$.}
\step{因此，如果公共点存在，只可能是 $(5,10)$.}
\step{接下来验证这个点是否属于所有集合，}
\step{对任意 $p\in M$，把 $x=5$ 代入集合的关系式，}
\step{得 $3x^2-px+5p-65=75-5p+5p-65=10$，与 $y=10$ 相符，}
\step{所以 $(5,10)\in A_p$ 对每一个 $p$ 都成立，因此所有 $A_p$ 的交集为 $\{(5,10)\}$.}
\solend
\fi

\end{enumerate}

% 守卫须放在 \bigsec 之前，否则标题仍会落页末（见 高一40 的教训）
\needvspace{5cm}

\bigsec{二、选择题}

\begin{enumerate}
\setcounter{enumi}{12}

\item 已知 $a,b,c\in\R$ 且 $a>b$，则下列不等式一定成立的是\mbox{（\quad）}

\keepwithprev
A. $\dfrac1a<\dfrac1b$\qquad B. $a^2>b^2$

C. $a|c|>b|c|$\qquad D. $\dfrac a{c^2+1}>\dfrac b{c^2+1}$

\ifanswer
\solbegin
\step{$A$：当 $a>0>b$ 时，$\dfrac1a>\dfrac1b$，故 $A$ 错误；}
\step{$B$：当 $a=-1,b=-2$ 时，满足 $a>b$，但 $a^2>b^2$ 不成立，故 $B$ 错误；}
\step{$C$：当 $c=0$ 时，$a|c|=b|c|$，故 $C$ 错误；}
\step{$D$：由 $a>b,c^2+1>0$，得 $\dfrac a{c^2+1}>\dfrac b{c^2+1}$，故 $D$ 正确；}
\step{故选 $D$.}
\solend
\fi

\item 已知集合 $A=\{a^2,0,-1\}$，$B=\{a,b,0\}$，若 $A=B$，则 $(ab)^{2026}$ 的值为\mbox{（\quad）}

\keepwithprev
A. $0$\qquad B. $-1$

C. $1$\qquad D. $\pm1$

\ifanswer
\solbegin
\step{因为集合 $A=\{a^2,0,-1\}$，$B=\{a,b,0\}$，$A=B$，所以 $\begin{cases}a=-1\\ a^2=b\end{cases}$，}
\step{所以 $a=-1$，$b=1$，所以 $(ab)^{2026}=1$，故选 $C$.}
\solend
\fi

\item 设全集 $U=\R$，$A=\{x\mid x<-1\text{ 或 }x>2\}$,$B=\{y\mid y=x^2+a\}$，若 $\overline A\subseteq B$，则实数 $a$ 的取值范围是\mbox{（\quad）}

\keepwithprev
A. $a\leqslant-1$\qquad B. $a>2$

C. $a<-1$\qquad D. $a\geqslant-2$

\ifanswer
\solbegin
\step{因为 $A=\{x\mid x<-1\text{ 或 }x>2\}$，则 $\overline A=\{x\mid-1\leqslant x\leqslant2\}$，}
\step{又 $B=\{y\mid y=x^2+a\}=\{y\mid y\geqslant a\}$，且 $\overline A\subseteq B$，所以 $a\leqslant-1$，故选 $A$.}
\solend
\fi

\item 设 $P$ 和 $Q$ 是两个集合，定义集合 $P-Q=\{x\mid x\in P$，且 $x\notin Q\}$，如果 $Q=\{x\mid(1-x)(x-3)>0\}$，$P=\left\{x\left|\dfrac x{2-x}>0\right.\right\}$，那么 $P-Q$ 等于\mbox{（\quad）}

\keepwithprev
A. $\{x\mid0<x<1\}$\qquad B. $\{x\mid1\leqslant x<2\}$

C. $\{x\mid0<x\leqslant1\}$\qquad D. $\{x\mid2\leqslant x<3\}$

\ifanswer
\solbegin
\step{因为 $P=\left\{x\left|\dfrac x{2-x}>0\right.\right\}=\{x\mid x(x-2)<0\}=\{x\mid0<x<2\}$，}
\step{$Q=\{x\mid(1-x)(x-3)>0\}=\{x\mid(x-3)(x-1)<0\}=\{x\mid1<x<3\}$，}
\step{所以 $P-Q=\{x\mid0<x\leqslant1\}$，故选 $C$.}
\solend
\fi

\end{enumerate}

\bigsec{三、解答题}

\begin{enumerate}
\setcounter{enumi}{16}

\item 已知集合 $A=\{x\mid x^2-9\geqslant0\}$，$B=\{x\mid|x-4|<2\}$，$C=\left\{x\left|\dfrac{x-8}{x+2}<0\right.\right\}$.

（1）求 $A\cup C$；

（2）若全集 $U=R$，求 $\overline A\cap B$.

\ifanswer
\solbegin
\stepm{(1)}{集合 $A=\{x\mid x^2-9\geqslant0\}=\{x\mid x\geqslant3\text{ 或 }x\leqslant-3\}$，}
\step{$B=\{x\mid|x-4|<2\}=\{x\mid2<x<6\}$，$C=\left\{x\left|\dfrac{x-8}{x+2}<0\right.\right\}=\{x\mid-2<x<8\}$，}
\step{所以 $A\cup C=\{x\mid x\leqslant-3\text{ 或 }x>-2\}$；}
\stepm{(2)}{全集 $U=R$，集合 $A=\{x\mid x^2-9\geqslant0\}=\{x\mid x\geqslant3\text{ 或 }x\leqslant-3\}$，}
\step{$B=\{x\mid|x-4|<2\}=\{x\mid2<x<6\}$，}
\step{所以 $\overline A=\{x\mid-3<x<3\}$，所以 $\overline A\cap B=\{x\mid2<x<3\}$.}
\solend
\fi

\ansspace[6]

\item 已知集合 $A=\left\{x\left|\dfrac{2x-1}{x+1}<1,x\in\R\right.\right\}$，集合 $B=\{x\mid(x-a-1)(x-a+1)\leqslant0\}$.

（1）求集合 $A$；

（2）若 $B\subseteq A$，求实数 $a$ 的取值范围.

\ifanswer
\solbegin
\stepm{(1)}{由 $\dfrac{2x-1}{x+1}<1$，得 $\dfrac{x-2}{x+1}<0$ 所以 $(x-2)(x+1)<0$，解得 $-1<x<2$，}
\step{即 $A=(-1,2)$；}
\stepm{(2)}{由 $(x-a-1)(x-a+1)\leqslant0$，得 $a-1\leqslant x\leqslant a+1$，}
\step{所以 $B=\{x\mid a-1\leqslant x\leqslant a+1\}$，}
\step{又 $B\subseteq A$，且 $a+1>a-1$ 即 $B\neq\varnothing$，所以 $\begin{cases}a-1>-1\\ a+1<2\end{cases}$，所以 $0<a<1$，}
\step{所以实数 $a$ 的取值范围是 $\{a\mid0<a<1\}$.}
\solend
\fi

\ansspace[6]

\item 已知集合 $A=\{x\mid x^2-4x+3=0\}$，$B=\{x\mid x^2-ax+3=0\}$.

（1）若 $A\cup B=B$，求实数 $a$ 的值；

（2）若 $A\cap B=B$，求实数 $a$ 的取值范围.

\ifanswer
\solbegin
\stepm{(1)}{集合 $A=\{x\mid x^2-4x+3=0\}=\{1,3\}$，$B=\{x\mid x^2-ax+3=0\}$.}
\step{$A\cup B=B$，所以 $1\in B$，$3\in B$，所以 $\begin{cases}1-a+3=0\\ 9-3a+3=0\end{cases}$，解得实数 $a=4$；}
\stepm{(2)}{因为 $A\cap B=B$，所以 $B\subseteq A$，}
\step{当 $B=\varnothing$ 时，$\Delta=(-a)^2-12<0$，解得 $-2\sqrt3<a<2\sqrt3$，}
\step{当 $B$ 是单元素集时，$\Delta=a^2-12=0$，解得 $a=2\sqrt3$，此时 $B=\{\sqrt3\}$，}
\step{不合题意；}
\step{当 $B=\{1,3\}$ 时，由 (1) 得 $a=4$，}
\step{综上，实数 $a$ 的取值范围是 $(-2\sqrt3,2\sqrt3)\cup\{4\}$.}
\solend
\fi

\ansspace[6]

\item 已知全集为 $R$，实数 $a$、$b$ 满足 $a\neq b$，集合 $M=\{x\mid x^2-3x-4\leqslant0\}$，集合 $A=\{x\mid(x+a)(x+b)>0\}$，集合 $B=\{x\mid x^2+(a-2)x-2a>0\}$.

（1）若 $\overline A=M$，求 $a$、$b$ 的值；

（2）若 $a>b>-1$，求 $A\cap B$；

（3）若 $a^2+a\in\overline B$，求 $a$ 的取值范围.

\ifanswer
\solbegin
\stepm{(1)}{$M=\{x\mid x^2-3x-4\leqslant0\}=\{x\mid-1\leqslant x\leqslant4\}$，$A=\{x\mid(x+a)(x+b)\leqslant0\}$，}
\step{因为 $\overline A=M$，所以 $a=1$，$b=-4$ 或 $a=-4$，$b=1$.}
\stepm{(2)}{因为 $a>b>-1$，所以 $-a<-b<1$，所以 $A=\{x\mid x<-a\text{ 或 }x>-b\}$，}
\step{因为 $B=\{x\mid x^2+(a-2)x-2a>0\}=\{x\mid x<-a\text{ 或 }x>2\}$，}
\step{所以 $A\cap B=\{x\mid x<-a\text{ 或 }x>2\}$，}
\stepm{(3)}{$\overline B=\{x\mid(x-2)(x+a)\leqslant0\}$，}
\step{由 $a^2+a\in\overline B$，得 $(a^2+a-2)(a^2+2a)\leqslant0\Leftrightarrow a(a-1)(a-2)^2\leqslant0$，}
\step{解得 $a=2$ 或 $0\leqslant a\leqslant1$，故 $a$ 的取值范围是 $[0,1]\cup\{2\}$.}
\solend
\fi

\ansspace[8]

\item 对平面直角坐标系的第一象限内的任意两点作如下定义：若 $\dfrac ab>\dfrac cd$，那么称点 $(a,b)$ 是点 $(c,d)$ 的“上位点”，同时点 $(c,d)$ 是点 $(a,b)$ 的“下位点”.

（1）试写出点 $(2,4)$ 的一个“上位点”坐标；

（2）已知点 $(a,b)$ 是点 $(c,d)$ 的“上位点”，判断点 $P(a+c,b+d)$ 是点 $(c,d)$ 的“上位点”还是“下位点”，并证明你的结论；

（3）设正整数 $n$ 满足以下条件：对集合 $\{t\mid0<t<2020,t\in\Z\}$ 内的任意元素 $m$，总存在正整数 $k=2m+1$，使得点 $(n,k)$ 既是点 $(2025,m)$ 的“下位点”，又是点 $(2026,m+1)$ 的“上位点”，问这样的正整数 $n$ 是否存在？若存在，求正整数 $n$ 的值；若不存在，请说明理由.

\ifanswer
\solbegin
\stepm{(1)}{由定义，点 $(a,b)$ 是点 $(2,4)$ 的“上位点”，}
\step{当且仅当 $\dfrac ab>\dfrac24=\dfrac12$ 且 $a>0,b>0$.}
\step{取 $(a,b)=(1,1)$，因为 $\dfrac11>\dfrac12$，}
\step{所以点 $(1,1)$ 是点 $(2,4)$ 的一个“上位点”.（答案不唯一）}
\stepm{(2)}{点 $P(a+c,b+d)$ 是点 $(c,d)$ 的“上位点”.}
\step{$(a,b)$ 是点 $(c,d)$ 的“上位点”，所以 $\dfrac ab>\dfrac cd$.}
\step{因为 $a,b,c,d>0$，两边同乘 $bd$ 得 $ad>bc$，即 $ad-bc>0$.}
\step{作差比较 $\dfrac{a+c}{b+d}$ 与 $\dfrac cd$：$\dfrac{a+c}{b+d}-\dfrac cd=\dfrac{d(a+c)-c(b+d)}{d(b+d)}=\dfrac{ad-bc}{d(b+d)}>0$，}
\step{所以点 $P(a+c,b+d)$ 是点 $(c,d)$ 的“上位点”.}
\stepm{(3)}{由定义，点 $(n,k)$ 是点 $(2025,m)$ 的下位点即 $\dfrac{2025}m>\dfrac nk$；}
\step{点 $(n,k)$ 是点 $(2026,m+1)$ 的上位点即 $\dfrac nk>\dfrac{2026}{m+1}$.}
\step{代入 $k=2m+1$，得对任意 $m\in\{1,2,\cdots,2019\}$，}
\step{都有 $4052-\dfrac{2026}{m+1}<n<4050+\dfrac{2025}m$.}
\step{右边 $4050+\dfrac{2025}m$ 随 $m$ 增大而减小，故最近处在 $m=2019$，}
\step{得 $n<4050+\dfrac{2025}{2019}<4051.01$，所以 $n\leqslant4051$.}
\step{左边 $4052-\dfrac{2026}{m+1}$ 随 $m$ 增大而增大，故最近处也在 $m=2019$，}
\step{得 $n>4052-\dfrac{2026}{2020}>4050.99$，所以 $n\geqslant4051$.}
\step{综上，$n=4051$.}
\step{检验：当 $n=4051$ 时，对任意 $1\leqslant m\leqslant2019,k=2m+1$ 是正整数，}
\step{且 $\dfrac{4051}{2m+1}<\dfrac{2025}m$,$4051m<4050m+2025,m<2025$，成立；}
\step{$\dfrac{4051}{2m+1}>\dfrac{2026}{m+1}$,$4051(m+1)>4052m+2026,m<2025$，成立.}
\step{又各点横纵坐标均为正，都在第一象限内，符合定义，}
\step{所以满足条件的正整数 $n$ 存在，其值为 $n=4051$.}
\solend
\fi

\ansspace*[10]

\end{enumerate}

\end{document}
"
%
% 原始材料：微信公众号“魔都数学加油站”文章，作者破晓，连载“27学年高一59”；
% 原文标题“27学年高一59——2026-2027年上海市华实高中高一上10月月考”；
% 原文链接：https://mp.weixin.qq.com/s/XzyBCmrvpRuMe-NuNk7gkA。
% 页面标注发布时间 2026-10-13 21:45（上海）。
% 正文为 9 张 PNG 页：第 01、02、03、05、06、07、08、09 页 4762×6735，第 04 页 2560×3620
% （**同一篇各页分辨率不同**），均按页序读取。9 页全为试卷正文页（卷首标题在第 1 页页顶，
% 卷末解答题第 21 题解析止于第 9 页），无封面页、目录页、二维码或宣传页。
% 粉色斜排“魔都数学加油站”为卷面水印，与公众号页脚一并未录入。
%
% 卷首标题：照卷面印作“2026--2027 年华实高中高一上 10 月月考”（公众号标题中的“上海市”
% 卷面标题行没有，本稿照卷面），年份区间按全稿惯例排 en dash。原卷只有一行标题、无副标题。
%
% 篇幅与题号：原卷自第 3 题起编号（第 1、2 题不在本卷内）。填空题 3—12（10 题）、
% 选择题 13—16（4 题）、解答题 17—21（5 题），共 19 题。本稿沿用原节与原题号，
% 只改排为全稿统一的大节样式（原卷小节标题为“一、填空题”“二、选择题”“三、解答题”）。
%
% 讲评版：原卷每题下直接印【答案】或【解析】。**只印【答案】的 1 题**（第 3 题），本稿只给 \ans{}；
% 其余 18 题都印【解析】，按全稿口径这类题不另提 \ans{}——第 13、14、15、16 题的结论写在解析末句
% （“故选 D”等），亦照录在解析内。本卷**没有**印在题干括号内的答案。
%
% 副标题：原卷未印副标题。本稿按“知识模块”口径标作“集合与常用逻辑用语、不等式”：
% 19 题中集合与常用逻辑用语 13 题（第 3、9、10、11、12、14、15、16、17、18、19、20、21 题），
% 不等式 10 题（含等式与不等式性质、一元二次不等式、恒成立、基本应用：第 4、5、6、7、8、
% 13、17、18、19、20 题），两类均远超三分之一，故并标。口径同上。
%
% 与原卷的排版差异：
%   ① 原卷每题下直接印【答案】/【解析】，本稿按全稿惯例将答案与解析置于答案版，试卷版保持干净；
%   ② 原卷未印副标题，本稿按“知识模块”口径补出；
%   ③ 原卷填空横线约两个字宽，本稿按全稿惯例统一排 \blankline[4.4em]；
%   ④ 原卷实数集排粗体 $\mathbf{R}$（第 12、13、20 题等），整数集排粗体 $\mathbf{Z}$（第 11、12 题），
%      自然数集排斜体 $N$ 或粗体 $\mathbf{Z}$ 处按本稿惯例改用 $\R$、$\Z$、$\N$；
%      空集记号统一为 \varnothing；
%   ⑤ 原卷在数学式内多处用半角逗号（如 $x_1,x_2$、$a,b\in R$），均照录未强改；
%      句末句点按全稿惯例排西文句点；“漂亮集”“上位点”“下位点”等自定义名词照录；
%   ⑥ 第 17—21 题的子问原卷各自另起一行，本稿按全稿惯例用 enumerate 的 (1)(2)(3)；
%      解析中的子问标号一律用 \stepm{(1)}{…}；第 9、11 题解析中的①、②（③④）分情形或分结论，
%      按全稿惯例作 \stepm{①}{…}、\stepm{②}{…}（第 11 题解析末“故正确的是②④”照录为行内）；
%   ⑦ 第 12 题题干里的“$A_p$”下标原卷作斜体 $p$，照录；
%   ⑧ 每道解答题原卷未留作答空白，本稿按全稿惯例加 \ansspace。
%
% 原卷笔误/存疑（均照录未改，正文对应处另有 % 注）：本卷首轮逐页转录，暂未发现影响答案的笔误。
%
% 插图：无（全卷 9 页无任何图形、无“如图”字样）。
%
\documentclass[a4paper,11pt]{article}
\usepackage{common}

\begin{document}

\examtitlest[3]{2026--2027 年华实高中高一上 10 月月考}{集合与常用逻辑用语、不等式}

\bigsec{一、填空题}

\begin{enumerate}
\setcounter{enumi}{2}

\item “$x\geqslant1$ 且 $y\geqslant1$”的否定形式为\blankline[4.4em].

\ifanswer
\ans{$x<1$ 或 $y<1$}
\fi

\item 已知方程 $2x^2+4x-3=0$ 的两个根为 $x_1,x_2$，则 $\dfrac1{x_1}+\dfrac1{x_2}=$\blankline[4.4em].

\ifanswer
\solbegin
\step{利用韦达定理得 $\begin{cases}x_1+x_2=-2\\ x_1\cdot x_2=-\dfrac32\end{cases}$，则 $\dfrac1{x_1}+\dfrac1{x_2}=\dfrac{x_1+x_2}{x_1\cdot x_2}=\dfrac43$.}
\solend
\fi

\item 设 $a>b>0$，比较大小：$\dfrac1{a-b}$\blankline[2.4em]$\dfrac1a$.

\ifanswer
\solbegin
\step{由 $a>b>0$ 得 $a>0$ 且 $a-b>0$，所以 $\dfrac1{a-b}-\dfrac1a=\dfrac b{a(a-b)}>0$，}
\step{故 $\dfrac1{a-b}>\dfrac1a$.}
\solend
\fi

\item 若不等式 $2ax^2-ax+1>0$ 的解集为 $R$，则实数 $a$ 的取值范围是\blankline[4.4em].

\ifanswer
\solbegin
\step{当 $a\neq0$ 时，由题意得 $\begin{cases}a>0\\ \Delta<0\end{cases}$，即 $\begin{cases}a>0\\ a^2-8a<0\end{cases}$，解得 $0<a<8$.}
\step{当 $a=0$ 时，恒有 $1>0$，不等式也成立.}
\step{综上所述，$a$ 的取值范围是 $[0,8)$.}
\solend
\fi

\item 集合 $A=\{x\mid m-1\leqslant x\leqslant m+7\}$，集合 $B=\{x\mid x^2-3x-10\geqslant0\}$，若 $A\cup B=R$，则实数 $m$ 的取值范围为\blankline[4.4em].

\ifanswer
\solbegin
\step{$A=\{x\mid m-1\leqslant x\leqslant m+7\}$，$B=\{x\mid x^2-3x-10\geqslant0\}=\{x\mid x\leqslant-2\text{ 或 }x\geqslant5\}$，}
\step{若 $A\cup B=R$，则 $\begin{cases}m-1\leqslant-2\\ m+7\geqslant5\end{cases}$，解得 $-2\leqslant m\leqslant-1$，}
\step{所以实数 $m$ 的取值范围为 $[-2,-1]$.}
\solend
\fi

\item 若关于 $x$ 的不等式 $ax^2+bx+c>0$ 的解集为 $(-1,3)$，则关于 $x$ 的不等式 $cx^2+bx+a>0$ 的解集为\blankline[4.4em].

\ifanswer
\solbegin
\step{因为不等式 $ax^2+bx+c>0$ 的解集为 $(-1,3)$，}
\step{所以 $-1$ 和 $3$ 是方程 $ax^2+bx+c=0$ 的实数解，}
\step{所以 $\begin{cases}-1+3=-\dfrac ba\\ -1\times3=\dfrac ca\\ a<0\end{cases}$，解得 $b=-2a$，$c=-3a$，且 $a<0$，}
\step{所以不等式 $cx^2+bx+a>0$ 可化为 $-3ax^2-2ax+a>0$，即 $3x^2+2x-1>0$，}
\step{解得 $x<-1$ 或 $x>\dfrac13$，所以该不等式的解集为 $(-\infty,-1)\cup\left(\dfrac13,+\infty\right)$.}
\solend
\fi

\item 若集合 $A=\{x\mid(a-1)x^2-2x+1=0\}$ 有且仅有 $2$ 个子集，则实数 $a$ 的值为\blankline[4.4em].

\ifanswer
\solbegin
\step{由题意得方程 $(a-1)x^2-2x+1=0$ 只有一个根，}
\stepm{①}{当 $a-1=0$ 时，$a=1$，此时方程化为 $-2x+1=0$，解得 $x=\dfrac12$，符合题意，}
\stepm{②}{当 $a-1\neq0$，即 $a\neq1$ 时，$\Delta=4-4(a-1)=0$，解得 $a=2$，}
\step{综上所述，实数 $a$ 的值为 $1$ 或 $2$.}
\solend
\fi

\item “$x+y\neq5$”是“$x\neq3$ 或 $y\neq2$”的\blankline[4.4em]条件.

\ifanswer
\solbegin
\step{分析充分性：}
\step{若 $x+y\neq5$ 成立，若 $x=3$，则 $y\neq2$，若 $y=2$，则 $x\neq3$，}
\step{所以一定能说明 $x\neq3$ 和 $y\neq2$ 至少有一个成立，即 $x\neq3$ 或 $y\neq2$ 成立，}
\step{充分性成立；}
\step{分析必要性：}
\step{若 $x\neq3$ 或 $y\neq2$，即 $x\neq3$ 和 $y\neq2$ 二者至少一个成立，}
\step{举例：$x=1,y=4$，此时 $x+y=5$，所以必要性不成立.}
\step{所以“$x+y\neq5$”是“$x\neq3$ 或 $y\neq2$”的充分不必要条件.}
\solend
\fi

\item 设 $P$ 为非空实数集满足：对任意给定的 $x,y\in P$（$x,y$ 可以相同），都有 $x+y\in P,x-y\in P,xy\in P$，则称 $P$ 为漂亮集. 关于漂亮集有下列结论：

①集合 $P=\{-2,-1,0,1,2\}$ 为漂亮集；②集合 $P=\{x\mid x=2n,n\in\Z\}$ 为漂亮集；

③若集合 $P_1,P_2$ 为漂亮集，则 $P_1\cup P_2$ 为漂亮集；④若集合 $P$ 为漂亮集，则一定有 $0\in P$.

其中所有正确结论的序号是\blankline[4.4em].

\ifanswer
\solbegin
\step{对于集合 $P=\{-2,-1,0,1,2\}$，取 $x=2,y=-2$，则 $x-y=2-(-2)=4\notin P$，}
\step{不满足漂亮集的定义，所以①错误.}
\step{对于集合 $P=\{x\mid x=2n,n\in Z\}$，设 $x=2n_1,y=2n_2,n_1,n_2\in Z$.}
\step{$x+y=2n_1+2n_2=2(n_1+n_2)$，因为 $n_1+n_2\in Z$，所以 $x+y\in P$.}
\step{$x-y=2n_1-2n_2=2(n_1-n_2)$，因为 $n_1-n_2\in Z$，所以 $x-y\in P$.}
\step{$xy=2n_1\times2n_2=4n_1n_2=2(2n_1n_2)$，因为 $2n_1n_2\in Z$，所以 $xy\in P$.}
\step{满足漂亮集的定义，所以②正确.}
\step{令 $P_1=\{x\mid x=2n,n\in Z\}$,$P_2=\{x\mid x=3n,n\in Z\}$，都是漂亮集.}
\step{$P_1\cup P_2$ 中，取 $x=2(x\in P_1)$,$y=3(y\in P_2)$,$x+y=5\notin P_1\cup P_2$，}
\step{不满足漂亮集的定义，所以③错误.}
\step{因为对任意 $x\in P$，有 $x-x=0\in P$，所以若集合 $P$ 为漂亮集，}
\step{则一定有 $0\in P$，所以④正确.}
\step{故正确的是②④.}
\solend
\fi

\item 已知：集合 $A_p=\{(x,y)\mid y=3x^2-px+5p-65,x\in\R\}$,$M=\{x\mid1\leqslant x\leqslant100,x\in\Z\}$，若 $p\in M$，则所有 $A_p$ 的交集\blankline[4.4em].

\ifanswer
\solbegin
\step{由 $M=\{1,2,\cdots,100\}$，所以所有 $A_p$ 的交集就是这些集合中共同拥有的点.}
\step{若 $(x,y)$ 是所有集合的一个公共点，}
\step{那么它特别要同时属于 $p=1$ 和 $p=2$ 时对应的两个集合.}
\step{当 $p=1$ 时，$y=3x^2-x-60$；}
\step{当 $p=2$ 时，$y=3x^2-2x-55$；联立得 $x=5,y=10$.}
\step{因此，如果公共点存在，只可能是 $(5,10)$.}
\step{接下来验证这个点是否属于所有集合，}
\step{对任意 $p\in M$，把 $x=5$ 代入集合的关系式，}
\step{得 $3x^2-px+5p-65=75-5p+5p-65=10$，与 $y=10$ 相符，}
\step{所以 $(5,10)\in A_p$ 对每一个 $p$ 都成立，因此所有 $A_p$ 的交集为 $\{(5,10)\}$.}
\solend
\fi

\end{enumerate}

% 守卫须放在 \bigsec 之前，否则标题仍会落页末（见 高一40 的教训）
\needvspace{5cm}

\bigsec{二、选择题}

\begin{enumerate}
\setcounter{enumi}{12}

\item 已知 $a,b,c\in\R$ 且 $a>b$，则下列不等式一定成立的是\mbox{（\quad）}

\keepwithprev
A. $\dfrac1a<\dfrac1b$\qquad B. $a^2>b^2$

C. $a|c|>b|c|$\qquad D. $\dfrac a{c^2+1}>\dfrac b{c^2+1}$

\ifanswer
\solbegin
\step{$A$：当 $a>0>b$ 时，$\dfrac1a>\dfrac1b$，故 $A$ 错误；}
\step{$B$：当 $a=-1,b=-2$ 时，满足 $a>b$，但 $a^2>b^2$ 不成立，故 $B$ 错误；}
\step{$C$：当 $c=0$ 时，$a|c|=b|c|$，故 $C$ 错误；}
\step{$D$：由 $a>b,c^2+1>0$，得 $\dfrac a{c^2+1}>\dfrac b{c^2+1}$，故 $D$ 正确；}
\step{故选 $D$.}
\solend
\fi

\item 已知集合 $A=\{a^2,0,-1\}$，$B=\{a,b,0\}$，若 $A=B$，则 $(ab)^{2026}$ 的值为\mbox{（\quad）}

\keepwithprev
A. $0$\qquad B. $-1$

C. $1$\qquad D. $\pm1$

\ifanswer
\solbegin
\step{因为集合 $A=\{a^2,0,-1\}$，$B=\{a,b,0\}$，$A=B$，所以 $\begin{cases}a=-1\\ a^2=b\end{cases}$，}
\step{所以 $a=-1$，$b=1$，所以 $(ab)^{2026}=1$，故选 $C$.}
\solend
\fi

\item 设全集 $U=\R$，$A=\{x\mid x<-1\text{ 或 }x>2\}$,$B=\{y\mid y=x^2+a\}$，若 $\overline A\subseteq B$，则实数 $a$ 的取值范围是\mbox{（\quad）}

\keepwithprev
A. $a\leqslant-1$\qquad B. $a>2$

C. $a<-1$\qquad D. $a\geqslant-2$

\ifanswer
\solbegin
\step{因为 $A=\{x\mid x<-1\text{ 或 }x>2\}$，则 $\overline A=\{x\mid-1\leqslant x\leqslant2\}$，}
\step{又 $B=\{y\mid y=x^2+a\}=\{y\mid y\geqslant a\}$，且 $\overline A\subseteq B$，所以 $a\leqslant-1$，故选 $A$.}
\solend
\fi

\item 设 $P$ 和 $Q$ 是两个集合，定义集合 $P-Q=\{x\mid x\in P$，且 $x\notin Q\}$，如果 $Q=\{x\mid(1-x)(x-3)>0\}$，$P=\left\{x\left|\dfrac x{2-x}>0\right.\right\}$，那么 $P-Q$ 等于\mbox{（\quad）}

\keepwithprev
A. $\{x\mid0<x<1\}$\qquad B. $\{x\mid1\leqslant x<2\}$

C. $\{x\mid0<x\leqslant1\}$\qquad D. $\{x\mid2\leqslant x<3\}$

\ifanswer
\solbegin
\step{因为 $P=\left\{x\left|\dfrac x{2-x}>0\right.\right\}=\{x\mid x(x-2)<0\}=\{x\mid0<x<2\}$，}
\step{$Q=\{x\mid(1-x)(x-3)>0\}=\{x\mid(x-3)(x-1)<0\}=\{x\mid1<x<3\}$，}
\step{所以 $P-Q=\{x\mid0<x\leqslant1\}$，故选 $C$.}
\solend
\fi

\end{enumerate}

\bigsec{三、解答题}

\begin{enumerate}
\setcounter{enumi}{16}

\item 已知集合 $A=\{x\mid x^2-9\geqslant0\}$，$B=\{x\mid|x-4|<2\}$，$C=\left\{x\left|\dfrac{x-8}{x+2}<0\right.\right\}$.

（1）求 $A\cup C$；

（2）若全集 $U=R$，求 $\overline A\cap B$.

\ifanswer
\solbegin
\stepm{(1)}{集合 $A=\{x\mid x^2-9\geqslant0\}=\{x\mid x\geqslant3\text{ 或 }x\leqslant-3\}$，}
\step{$B=\{x\mid|x-4|<2\}=\{x\mid2<x<6\}$，$C=\left\{x\left|\dfrac{x-8}{x+2}<0\right.\right\}=\{x\mid-2<x<8\}$，}
\step{所以 $A\cup C=\{x\mid x\leqslant-3\text{ 或 }x>-2\}$；}
\stepm{(2)}{全集 $U=R$，集合 $A=\{x\mid x^2-9\geqslant0\}=\{x\mid x\geqslant3\text{ 或 }x\leqslant-3\}$，}
\step{$B=\{x\mid|x-4|<2\}=\{x\mid2<x<6\}$，}
\step{所以 $\overline A=\{x\mid-3<x<3\}$，所以 $\overline A\cap B=\{x\mid2<x<3\}$.}
\solend
\fi

\ansspace[6]

\item 已知集合 $A=\left\{x\left|\dfrac{2x-1}{x+1}<1,x\in\R\right.\right\}$，集合 $B=\{x\mid(x-a-1)(x-a+1)\leqslant0\}$.

（1）求集合 $A$；

（2）若 $B\subseteq A$，求实数 $a$ 的取值范围.

\ifanswer
\solbegin
\stepm{(1)}{由 $\dfrac{2x-1}{x+1}<1$，得 $\dfrac{x-2}{x+1}<0$ 所以 $(x-2)(x+1)<0$，解得 $-1<x<2$，}
\step{即 $A=(-1,2)$；}
\stepm{(2)}{由 $(x-a-1)(x-a+1)\leqslant0$，得 $a-1\leqslant x\leqslant a+1$，}
\step{所以 $B=\{x\mid a-1\leqslant x\leqslant a+1\}$，}
\step{又 $B\subseteq A$，且 $a+1>a-1$ 即 $B\neq\varnothing$，所以 $\begin{cases}a-1>-1\\ a+1<2\end{cases}$，所以 $0<a<1$，}
\step{所以实数 $a$ 的取值范围是 $\{a\mid0<a<1\}$.}
\solend
\fi

\ansspace[6]

\item 已知集合 $A=\{x\mid x^2-4x+3=0\}$，$B=\{x\mid x^2-ax+3=0\}$.

（1）若 $A\cup B=B$，求实数 $a$ 的值；

（2）若 $A\cap B=B$，求实数 $a$ 的取值范围.

\ifanswer
\solbegin
\stepm{(1)}{集合 $A=\{x\mid x^2-4x+3=0\}=\{1,3\}$，$B=\{x\mid x^2-ax+3=0\}$.}
\step{$A\cup B=B$，所以 $1\in B$，$3\in B$，所以 $\begin{cases}1-a+3=0\\ 9-3a+3=0\end{cases}$，解得实数 $a=4$；}
\stepm{(2)}{因为 $A\cap B=B$，所以 $B\subseteq A$，}
\step{当 $B=\varnothing$ 时，$\Delta=(-a)^2-12<0$，解得 $-2\sqrt3<a<2\sqrt3$，}
\step{当 $B$ 是单元素集时，$\Delta=a^2-12=0$，解得 $a=2\sqrt3$，此时 $B=\{\sqrt3\}$，}
\step{不合题意；}
\step{当 $B=\{1,3\}$ 时，由 (1) 得 $a=4$，}
\step{综上，实数 $a$ 的取值范围是 $(-2\sqrt3,2\sqrt3)\cup\{4\}$.}
\solend
\fi

\ansspace[6]

\item 已知全集为 $R$，实数 $a$、$b$ 满足 $a\neq b$，集合 $M=\{x\mid x^2-3x-4\leqslant0\}$，集合 $A=\{x\mid(x+a)(x+b)>0\}$，集合 $B=\{x\mid x^2+(a-2)x-2a>0\}$.

（1）若 $\overline A=M$，求 $a$、$b$ 的值；

（2）若 $a>b>-1$，求 $A\cap B$；

（3）若 $a^2+a\in\overline B$，求 $a$ 的取值范围.

\ifanswer
\solbegin
\stepm{(1)}{$M=\{x\mid x^2-3x-4\leqslant0\}=\{x\mid-1\leqslant x\leqslant4\}$，$A=\{x\mid(x+a)(x+b)\leqslant0\}$，}
\step{因为 $\overline A=M$，所以 $a=1$，$b=-4$ 或 $a=-4$，$b=1$.}
\stepm{(2)}{因为 $a>b>-1$，所以 $-a<-b<1$，所以 $A=\{x\mid x<-a\text{ 或 }x>-b\}$，}
\step{因为 $B=\{x\mid x^2+(a-2)x-2a>0\}=\{x\mid x<-a\text{ 或 }x>2\}$，}
\step{所以 $A\cap B=\{x\mid x<-a\text{ 或 }x>2\}$，}
\stepm{(3)}{$\overline B=\{x\mid(x-2)(x+a)\leqslant0\}$，}
\step{由 $a^2+a\in\overline B$，得 $(a^2+a-2)(a^2+2a)\leqslant0\Leftrightarrow a(a-1)(a-2)^2\leqslant0$，}
\step{解得 $a=2$ 或 $0\leqslant a\leqslant1$，故 $a$ 的取值范围是 $[0,1]\cup\{2\}$.}
\solend
\fi

\ansspace[8]

\item 对平面直角坐标系的第一象限内的任意两点作如下定义：若 $\dfrac ab>\dfrac cd$，那么称点 $(a,b)$ 是点 $(c,d)$ 的“上位点”，同时点 $(c,d)$ 是点 $(a,b)$ 的“下位点”.

（1）试写出点 $(2,4)$ 的一个“上位点”坐标；

（2）已知点 $(a,b)$ 是点 $(c,d)$ 的“上位点”，判断点 $P(a+c,b+d)$ 是点 $(c,d)$ 的“上位点”还是“下位点”，并证明你的结论；

（3）设正整数 $n$ 满足以下条件：对集合 $\{t\mid0<t<2020,t\in\Z\}$ 内的任意元素 $m$，总存在正整数 $k=2m+1$，使得点 $(n,k)$ 既是点 $(2025,m)$ 的“下位点”，又是点 $(2026,m+1)$ 的“上位点”，问这样的正整数 $n$ 是否存在？若存在，求正整数 $n$ 的值；若不存在，请说明理由.

\ifanswer
\solbegin
\stepm{(1)}{由定义，点 $(a,b)$ 是点 $(2,4)$ 的“上位点”，}
\step{当且仅当 $\dfrac ab>\dfrac24=\dfrac12$ 且 $a>0,b>0$.}
\step{取 $(a,b)=(1,1)$，因为 $\dfrac11>\dfrac12$，}
\step{所以点 $(1,1)$ 是点 $(2,4)$ 的一个“上位点”.（答案不唯一）}
\stepm{(2)}{点 $P(a+c,b+d)$ 是点 $(c,d)$ 的“上位点”.}
\step{$(a,b)$ 是点 $(c,d)$ 的“上位点”，所以 $\dfrac ab>\dfrac cd$.}
\step{因为 $a,b,c,d>0$，两边同乘 $bd$ 得 $ad>bc$，即 $ad-bc>0$.}
\step{作差比较 $\dfrac{a+c}{b+d}$ 与 $\dfrac cd$：$\dfrac{a+c}{b+d}-\dfrac cd=\dfrac{d(a+c)-c(b+d)}{d(b+d)}=\dfrac{ad-bc}{d(b+d)}>0$，}
\step{所以点 $P(a+c,b+d)$ 是点 $(c,d)$ 的“上位点”.}
\stepm{(3)}{由定义，点 $(n,k)$ 是点 $(2025,m)$ 的下位点即 $\dfrac{2025}m>\dfrac nk$；}
\step{点 $(n,k)$ 是点 $(2026,m+1)$ 的上位点即 $\dfrac nk>\dfrac{2026}{m+1}$.}
\step{代入 $k=2m+1$，得对任意 $m\in\{1,2,\cdots,2019\}$，}
\step{都有 $4052-\dfrac{2026}{m+1}<n<4050+\dfrac{2025}m$.}
\step{右边 $4050+\dfrac{2025}m$ 随 $m$ 增大而减小，故最近处在 $m=2019$，}
\step{得 $n<4050+\dfrac{2025}{2019}<4051.01$，所以 $n\leqslant4051$.}
\step{左边 $4052-\dfrac{2026}{m+1}$ 随 $m$ 增大而增大，故最近处也在 $m=2019$，}
\step{得 $n>4052-\dfrac{2026}{2020}>4050.99$，所以 $n\geqslant4051$.}
\step{综上，$n=4051$.}
\step{检验：当 $n=4051$ 时，对任意 $1\leqslant m\leqslant2019,k=2m+1$ 是正整数，}
\step{且 $\dfrac{4051}{2m+1}<\dfrac{2025}m$,$4051m<4050m+2025,m<2025$，成立；}
\step{$\dfrac{4051}{2m+1}>\dfrac{2026}{m+1}$,$4051(m+1)>4052m+2026,m<2025$，成立.}
\step{又各点横纵坐标均为正，都在第一象限内，符合定义，}
\step{所以满足条件的正整数 $n$ 存在，其值为 $n=4051$.}
\solend
\fi

\ansspace*[10]

\end{enumerate}

\end{document}
